% =========================================================================
% modenanalyse_2fe2s -- Handbuch und theoretische Dokumentation
% Version 1.1.0  (Juni 2026)
%
% Lukas Knauer, AG Schuenemann, RPTU Kaiserslautern-Landau
% =========================================================================
\documentclass[11pt,a4paper,twoside]{article}

% ---------- Encoding & Sprache -------------------------------------------
\usepackage[T1]{fontenc}
\usepackage[utf8]{inputenc}
\usepackage[main=ngerman,provide=*]{babel}
\usepackage{lmodern}
\usepackage{microtype}

% ---------- Mathematik ---------------------------------------------------
\usepackage{amsmath, amssymb, amsthm}
\usepackage{mathtools}
\usepackage{bm}
\usepackage{siunitx}
\sisetup{range-phrase = --, range-units = single}

% ---------- Layout & Seite -----------------------------------------------
\usepackage[a4paper, left=2.5cm, right=2.5cm, top=2.5cm, bottom=3cm]{geometry}
\usepackage{setspace}
\onehalfspacing
\setlength{\parindent}{0pt}
\setlength{\parskip}{0.5em}

% ---------- Grafik & Farben ----------------------------------------------
\usepackage{xcolor}
\definecolor{accent}{RGB}{0,90,140}
\definecolor{codebg}{RGB}{245,245,242}
\definecolor{codeborder}{RGB}{200,200,195}
\definecolor{tomlkey}{RGB}{120,40,120}
\definecolor{tomlstr}{RGB}{0,110,0}
\definecolor{tomlcomment}{RGB}{110,110,110}
\definecolor{warnbg}{RGB}{255,248,228}
\definecolor{warnborder}{RGB}{210,160,80}
\definecolor{infobg}{RGB}{232,242,250}
\definecolor{infoborder}{RGB}{120,160,200}

% ---------- Tabellen -----------------------------------------------------
\usepackage{booktabs}
\usepackage{array}
\usepackage{tabularx}
\usepackage{longtable}

% ---------- Code-Listings ------------------------------------------------
\usepackage{listings}
\lstset{
  basicstyle=\ttfamily\small,
  backgroundcolor=\color{codebg},
  frame=single,
  rulecolor=\color{codeborder},
  breaklines=true,
  breakatwhitespace=true,
  showstringspaces=false,
  columns=flexible,
  keepspaces=true,
  xleftmargin=4pt,
  xrightmargin=4pt,
  framexleftmargin=2pt,
  framexrightmargin=2pt,
  aboveskip=0.8em,
  belowskip=0.6em,
  inputencoding=utf8,
  extendedchars=true,
  literate=%
    {ä}{{\"a}}1 {ö}{{\"o}}1 {ü}{{\"u}}1
    {Ä}{{\"A}}1 {Ö}{{\"O}}1 {Ü}{{\"U}}1
    {ß}{{\ss}}1 {µ}{{$\mu$}}1
    {→}{{$\rightarrow$}}2
    {≤}{{$\leq$}}1 {≥}{{$\geq$}}1
    {·}{{$\cdot$}}1 {–}{{--}}1 {—}{{---}}1
    {⟨}{{$\langle$}}1 {⟩}{{$\rangle$}}1
    {²}{{$^{2}$}}1 {³}{{$^{3}$}}1,
}

\lstdefinestyle{python}{
  language=Python,
  keywordstyle=\color{accent}\bfseries,
  commentstyle=\color{gray}\itshape,
  stringstyle=\color{tomlstr},
}

\lstdefinestyle{bash}{
  basicstyle=\ttfamily\small,
  keywords={},
  backgroundcolor=\color{codebg},
}

% ---------- Pseudocode --------------------------------------------------
\usepackage{algorithm}
\usepackage{algpseudocode}
\usepackage{float}
\floatname{algorithm}{Pseudocode}

% ---------- TikZ / PGFplots ---------------------------------------------
\usepackage{tikz}
\usepackage{pgfplots}
\pgfplotsset{compat=1.18}
\usetikzlibrary{arrows.meta, calc, decorations.markings, patterns}
\renewcommand{\algorithmicrequire}{\textbf{Eingabe:}}
\renewcommand{\algorithmicensure}{\textbf{Ausgabe:}}
\renewcommand{\algorithmicforall}{\textbf{f\"ur alle}}
\renewcommand{\algorithmicif}{\textbf{wenn}}
\renewcommand{\algorithmicelse}{\textbf{sonst}}
\renewcommand{\algorithmicend}{\textbf{ende}}
\renewcommand{\algorithmicfor}{\textbf{f\"ur}}
\renewcommand{\algorithmicdo}{\textbf{tue}}
\renewcommand{\algorithmicwhile}{\textbf{solange}}
\renewcommand{\algorithmicreturn}{\textbf{gib zur\"uck}}

% ---------- Custom-Macros ------------------------------------------------
\providecommand{\angstrom}{\text{\AA}}    % Angstrom-Symbol fuer math-Modus

% ---------- Ueberschriften & TOC ----------------------------------------
\usepackage{titlesec}
\titleformat{\part}[display]
  {\Huge\bfseries\color{accent}}
  {\filleft\LARGE Teil \thepart}
  {0.5em}
  {\filleft}
\titleformat{\section}
  {\Large\bfseries\color{accent}}
  {\thesection}{0.7em}{}
\titleformat{\subsection}
  {\large\bfseries}
  {\thesubsection}{0.7em}{}

% ---------- Links & Verweise --------------------------------------------
\usepackage[colorlinks=true,
            linkcolor=accent,
            urlcolor=accent,
            citecolor=accent]{hyperref}

% ---------- Boxen fuer Hinweise / Warnungen -----------------------------
\usepackage[most]{tcolorbox}
\newtcolorbox{warnbox}[1][]{
  colback=warnbg, colframe=warnborder, boxrule=0.5pt, arc=2pt,
  left=8pt, right=8pt, top=4pt, bottom=4pt, #1
}
\newtcolorbox{infobox}[1][]{
  colback=infobg, colframe=infoborder, boxrule=0.5pt, arc=2pt,
  left=8pt, right=8pt, top=4pt, bottom=4pt, #1
}

% ---------- Shortcuts fuer Mathe ----------------------------------------
\newcommand{\evec}{\mathbf{e}}
\newcommand{\rvec}{\mathbf{r}}
\newcommand{\uvec}{\mathbf{u}}
\newcommand{\nhat}{\hat{\mathbf{n}}}
\newcommand{\xhat}{\hat{\mathbf{x}}}
\newcommand{\yhat}{\hat{\mathbf{y}}}
\newcommand{\zhat}{\hat{\mathbf{z}}}
\newcommand{\fLM}{f_{\mathrm{LM}}}
\newcommand{\ER}{E_{\mathrm{R}}}
\newcommand{\kB}{k_{\mathrm{B}}}
\newcommand{\urms}{u_{\mathrm{rms}}}
\newcommand{\Fei}{\mathrm{Fe}}
\newcommand{\OOP}{\mathrm{OOP}}
\newcommand{\INP}{\mathrm{INP}}
\newcommand{\Ag}{A_{\mathrm{g}}}
\newcommand{\Bgi}[1]{B_{\mathrm{#1g}}}
\newcommand{\Bui}[1]{B_{\mathrm{#1u}}}
\newcommand{\Aui}{A_{\mathrm{u}}}
\newcommand{\PCET}{\mathrm{PCET}}
\newcommand{\ET}{\mathrm{ET}}

% --- Formel-zu-Code-Referenz-Makro ---------------------------------------
% Verweist auf die Funktion im Quellbaum, die eine Formel umsetzt.
% Format: \codref{module.function}. Der Verifikationstest
% tests/test_manual_code_refs.py prueft alle \codref-Eintraege in
% Manual.tex (das DE-Manual spiegelt das EN 1:1, daher nicht
% separat getestet). Refaktorierungen im Code, die Manual-Refs
% verwaisen lassen, schlagen dort als Testfehler auf.
\newcommand{\codref}[1]{%
  \quad\text{\textcolor{gray!70!black}{\footnotesize $\rightsquigarrow$~\texttt{#1}}}}

% --- Zusaetzliche Macros (Theorie-Sektionen) -----------------------------
\newcommand{\khat}{\hat{\mathbf{k}}}
\newcommand{\bvec}{\mathbf{B}}
\newcommand{\Eiso}{e^{2}_{\mathrm{iso}}}
\newcommand{\Miso}{M_{\mathrm{iso}}}
\newcommand{\ket}[1]{\left|#1\right\rangle}
\newcommand{\bra}[1]{\left\langle#1\right|}

% ---------- Theoremumgebungen --------------------------------------------
\newtheorem{definition}{Definition}[section]
\newtheorem{satz}{Satz}[section]
\theoremstyle{remark}
\newtheorem*{anmerkung}{Anmerkung}

% ---------- Header & Footer ----------------------------------------------
\usepackage{fancyhdr}
\setlength{\headheight}{14pt}
\pagestyle{fancy}
\fancyhf{}
\fancyhead[LE,RO]{\thepage}
\fancyhead[RE]{\leftmark}
\fancyhead[LO]{\rightmark}
\renewcommand{\headrulewidth}{0.4pt}

% =========================================================================
\begin{document}

\begin{titlepage}
  \begin{center}
    \vspace*{2cm}
    {\Huge\bfseries\color{accent} modenanalyse\_2fe2s}\\[0.4cm]
    {\large Version 1.1.0}\\[2cm]

    \rule{\textwidth}{0.4pt}\\[0.3cm]
    {\LARGE\bfseries Handbuch und theoretische Dokumentation}\\[0.3cm]
    \rule{\textwidth}{0.4pt}\\[1cm]

    {\large
      Mode-aufgel\"oste Bindungs-Reorganisations-Analyse\\
      von {[}2Fe-2S{]}-Eisen-Schwefel-Proteinen\\
      aus Gaussian-16- und ORCA-6-Frequenzrechnungen
    }\\[1.5cm]

    {\large
      Lukas Knauer\\
      AG Sch\"unemann\\
      RPTU Kaiserslautern-Landau
    }\\[1.5cm]

    \begin{tabular}{rl}
      \textbf{Cluster-Typ:}    & {[}2Fe-2S{]}, reduziert oder oxidiert,\\
                               & f\"ur Rieske-, Ferredoxin- und mitoNEET-artige Liganden\\[0.3cm]
      \textbf{Input:}          & Gaussian-16 \texttt{.log} (\texttt{freq=hpmodes} bevorzugt;\\
                               & Standard-Block als Fallback) oder ORCA-6 \texttt{.hess}\\
                               & PDB-Datei mit gleichem Cluster (optional, aber empfohlen)\\[0.3cm]
      \textbf{Output:}         & Excel-Reports (Modenanalyse, SCSD,\\
                               & Reorganisations-Energien, Modulations-Spektren,\\
                               & UMAP-Cluster), Befund (Plain-Text), Plots (PNG)\\[0.3cm]
      \textbf{Externe Pakete:} & \texttt{numpy}, \texttt{scipy}, \texttt{pandas},\\
                               & \texttt{matplotlib}, \texttt{openpyxl},\\
                               & \texttt{scikit-learn}, \texttt{umap-learn},\\
                               & \texttt{hdbscan}, \texttt{scsdpy}\\[0.3cm]
      \textbf{Referenzen:}     & Marcus, \emph{J.\ Chem.\ Phys.} \textbf{24}, 966 (1956)\\
                               & Marcus, \emph{Angew.\ Chem.\ Int.\ Ed.\ Engl.} \textbf{32}, 1111 (1993)\\
                               & Hammes-Schiffer \& Soudackov,\\
                               & \emph{J.\ Phys.\ Chem.\ B} \textbf{112}, 14108 (2008)\\
                               & Huang \& Rhys,\\
                               & \emph{Proc.\ R.\ Soc.\ Lond.\ A} \textbf{204}, 406 (1950)\\
                               & Kingsbury \& Senge, \emph{Chem.\ Sci.} \textbf{15}, 13638 (2024)\\[0.3cm]
      \textbf{Lizenz:}         & GNU General Public License v3.0 oder sp\"ater\\
                               & (\texttt{GPL-3.0-or-later}, siehe \texttt{LICENSE})\\[0.3cm]
      \textbf{Copyright:}      & \textcopyright\ 2026 Lukas Knauer,\\
                               & AG Sch\"unemann, RPTU Kaiserslautern-Landau
    \end{tabular}

    \vfill
    {\large \today}
  \end{center}
\end{titlepage}

\tableofcontents
\clearpage

% =========================================================================
% Englisches Vorwort fuer internationale Leser
% =========================================================================
\section*{Foreword (English)}
\addcontentsline{toc}{section}{Foreword (English)}

\textbf{Note for English-speaking readers:} This manual is written in
German, reflecting the package's origin in a German-speaking research
group at RPTU Kaiserslautern-Landau. A full English translation is
planned for a future release.

The English-language entry points to this software are:

\begin{itemize}
  \item \texttt{README.md} -- complete English overview, installation
        instructions, quick start, and output structure description.
  \item \texttt{CHANGELOG.md} -- English release notes for v1.0.0.
  \item \texttt{CONTRIBUTING.md} -- English contribution guidelines.
  \item Source-code public API (function signatures, class fields):
        all identifiers are in English.
  \item Command-line interface (\texttt{modenanalyse-2fe2s --help}):
        all option help texts are in English.
\end{itemize}

The internal documentation (this manual + most source-code docstrings
and inline comments) remains in German for now. Public English
documentation covers the API surface that external users need; this
manual goes much deeper into theory, validation and implementation
details for German-speaking users (typical audience: graduate students
and postdocs in a German-speaking biophysics group).

For citation, see \texttt{CITATION.cff} (English).

\clearpage

% =========================================================================
\part{Einf\"uhrung und Theorie}
% =========================================================================

\section{Einf\"uhrung}
\label{sec:einleitung}

\texttt{modenanalyse\_2fe2s} ist ein wissenschaftliches Python-Werkzeug
zur Auswertung von Normalmoden in DFT-Frequenzrechnungen
(\texttt{freq=hpmodes}) von Eisen-Schwefel-Proteinen mit
{[}2Fe-2S{]}-Cluster. Es kombiniert geometrische Klassifikation,
strenge Symmetriezerlegung und nuklear-inelastische Streuung (NIS) in
einem einheitlichen Workflow. Anwendungsfelder sind Mechanismus-Studien
zu Proton-gekoppeltem Elektronen-Transfer (PCET), Vergleich von
Rieske- und Ferredoxin-artigen Systemen sowie Mutations-Effekte
(z.\,B.\ HH$\to$CC).

\subsection{Was \texttt{modenanalyse\_2fe2s} liefert}
\label{sec:scope}

Das Werkzeug fokussiert sich auf die strukturelle Analyse von
Schwingungsmoden in DFT-Frequenzrechnungen von [2Fe-2S]-Clustern.
Die Hauptergebnisse pro Mode sind:

\begin{itemize}
  \item \textbf{Cluster-spezifische Klassifikation}: Erkennung des
        {[}2Fe-2S{]}-Kerns, Klassifikation jeder Mode nach
        Out-of-plane / In-plane-Charakter, Symmetriezerlegung in
        D$_\mathrm{2h}$-Irreps.
  \item \textbf{Liganden-Aufl\"osung}: PDB-basierte Erkennung
        koordinierender His-, Cys-, Asp/Glu-Reste; Fe-N-, Fe-S- und
        Fe-O-Streckkennzahlen pro Mode.
  \item \textbf{Marcus-Hush-Reorganisationsenergien}: pro Mode
        $\Delta r_X$ und $\lambda_X$ fuer die Bindungs-Kanaele FeFe,
        FeN, FeS, NH und HA. System-Aggregate (Total-Reorg,
        Modulations-Spektren, kumulative $\Lambda_X(\omega)$) fuer
        die Banden-Identifikation extern in Origin.
  \item \textbf{Multi-Cluster-Unterst\"utzung}: Erkennung mehrerer
        {[}2Fe-2S{]}-Cluster im selben System (z.\,B.\ Glutaredoxin-Dimer)
        und parallele Auswertung als First-Class-Feature.
  \item \textbf{Embeddings} (UMAP, HDBSCAN-Clustering):
        zweidimensionale Projektion des Mode-Featureraums zur
        explorativen Analyse.
  \item \textbf{ORCA-Support} (eigenst\"andiger \texttt{.hess}-Parser):
        Mode-Analyse direkt aus ORCA-Frequenzrechnungen.
\end{itemize}

NIS-Spektren, Fe-pDOS, Lamb-M\"o\ss bauer-Faktor und
Fe-projizierte Thermodynamik sind \emph{nicht} Teil dieses Pakets.
Diese Gr\"o\ss en lassen sich aus derselben DFT-Frequenzrechnung mit
spezialisierten Werkzeugen separat berechnen.

\begin{infobox}
\textbf{Changelog-H\"ohepunkte seit v1.0.0.}
Diese Manual-Ausgabe beschreibt das Paket im Stand v1.1.0
(vollst\"andige Liste in \texttt{CHANGELOG.md}). Inhaltliche
\"Anderungen seit v1.0.0:
\begin{itemize}\setlength{\itemsep}{0pt}
  \item \textbf{v1.0.2}: Bugfix in \texttt{\_build\_group\_map} --
        verhindert stille Nullzeilen in den
        \texttt{Groups\_OOP/INP/Tors}-Bl\"attern, wenn Residuennummern
        mit Nicht-Protein-Eintr\"agen \"uberlappen (z.\,B.\ HOH-Wasser).
        F\"ur Proteinsysteme mit L\"osungsmittelresiduen ist
        mindestens v1.0.2 erforderlich.
  \item \textbf{v1.0.3}: neue Diagnose-Ausgaben -- das
        \texttt{Coordination}-Blatt dokumentiert die Ligand- und
        Gruppen-Atomzuordnung (\S\ref{app:wb-main-coordination});
        zwei zus\"atzliche UMAP-Einbettungen (SSE-UMAP und Ca-UMAP)
        zielen auf Protein-Dom\"anenmoden bzw.\ r\"aumlich
        lokalisierte Moden (\S\ref{app:wb-embed-ssumap},
        \S\ref{app:wb-embed-caumap}).
  \item \textbf{v1.0.4}: auditgetriebene Korrekturen der
        Unsicherheits-Fortpflanzung (\S\ref{sec:migration-v104}).
        Werte unver\"andert; die \texttt{s\_*}-Spalten f\"ur
        N-H-Stretches, SSE-Element-Amplituden und SCSD-Amplituden
        \"andern ihre Gr\"o\ss enordnung. Halb-offene
        Fenstergrenzen und defensive Warnmeldungen sind ebenfalls
        Teil dieses Release.
  \item \textbf{v1.0.5}: physikalische Korrekturen der
        Sekund\"arstrukturelement-Analyse -- die SSE-UMAP-Features
        \texttt{bending\_*} (stets Null) wurden durch die tats\"achlichen
        \texttt{lateral\_*}-Schl\"ussel ersetzt; \texttt{tilting\_angle}
        ist jetzt eine echte Starrk\"orper-Kippung; \texttt{com\_amplitude}
        (je Element und f\"ur den Clusterkern) ist massengewichtet; und die
        SSE-Hauptachse stammt aus dem C$\alpha$-R\"uckgrat. Die SSE-Bl\"atter
        und ihr UMAP \"andern sich, betroffene Systeme also einmal neu
        ausf\"uhren; der SSE-UMAP-Featuresatz wurde zudem auf eine
        dekorrelierte F\"unf-Metrik-Auswahl reduziert.
  \item \textbf{v1.1.0}: die Abk\"urzung \emph{SS} (Sekund\"arstruktur)
        wurde projektweit zu \emph{SSE} umbenannt -- Sheet-Namen, der
        Ausgabedateiname (\texttt{*\_analysis\_SSE.xlsx}),
        Konfigurationsschl\"ussel und Funktionsnamen. Das behebt die
        Mehrdeutigkeit mit Disulfid (S--S) und entspricht der \"ublichen
        Bezeichnung \emph{secondary structure element}. Alte
        TOML-Schl\"ussel (\texttt{analyze\_ss}, \texttt{ss\_chain}) werden
        weiterhin akzeptiert.
\end{itemize}
\end{infobox}

\subsection{Voraussetzungen}

\begin{itemize}
  \item Python $\geq$ 3.11 mit \texttt{numpy}, \texttt{scipy},
        \texttt{openpyxl}, \texttt{matplotlib}, \texttt{pandas},
        \texttt{scikit-learn}, \texttt{umap-learn}, \texttt{hdbscan},
        \texttt{scsdpy}. Alle werden \emph{automatisch} \"uber das
        Wheel installiert.
  \item Eine \textbf{Gaussian-16-Logdatei} aus einer
        \texttt{freq=hpmodes}-Rechnung des optimierten Clusters
        (mindestens vier optimierte Schweratome: zwei Fe und zwei
        br\"uckende S) \textbf{oder} eine \textbf{ORCA-6
        \texttt{.hess}-Datei}.
  \item Optional: \textbf{PDB-Datei} mit dem gleichen Cluster, f\"ur
        Liganden-Aufl\"osung und Sekund\"arstruktur-Analyse.
\end{itemize}

\subsection{Konventionen in diesem Handbuch}

\begin{itemize}
  \item \textbf{Achsen}: Wo immer m\"oglich wird die kanonische
        Cluster-Orientierung verwendet:
        $\xhat$ entlang Fe-Fe, $\yhat$ entlang S-S,
        $\zhat = \nhat$ als Cluster-Normale. Das ist konsistent mit
        der SCSD-Referenzgeometrie (siehe \S\ref{sec:scsd}).
  \item \textbf{Einheiten}: cm$^{-1}$ f\"ur Frequenzen,
        \AA{} f\"ur L\"angen, AMU f\"ur Massen, meV f\"ur Energien,
        Kelvin f\"ur Temperaturen. F\"ur Boltzmann-Faktoren wird
        $1/\kB = \SI{11.6048}{K/meV}$ verwendet.
  \item \textbf{Heuristisch vs.\ rigoros}: Wo eine Gr\"o\ss e aus
        einer Heuristik abgeleitet wird (z.\,B.\ die
        Bewegungsmuster-Scores Breathing, Hinge, ...), ist sie
        explizit als \emph{heuristisch} ausgewiesen. Strenge
        Symmetriezerlegung kommt aus SCSD.
\end{itemize}

\begin{infobox}
\textbf{Lese-Tipp.} Wer das Programm zum ersten Mal benutzt, sollte
\S\ref{sec:einleitung} und \S\ref{sec:konfig} lesen und mit dem
Beispiel in \S\ref{sec:workflow} starten. Die theoretischen Kapitel
\S\ref{sec:theorie} k\"onnen bei Bedarf vertieft werden, sind aber f\"ur
den Routine-Einsatz nicht zwingend.
\end{infobox}

% =========================================================================
\section{Theoretische Grundlagen}
\label{sec:theorie}

% -------------------------------------------------------------------------
\subsection{Normalmoden aus Gaussian: \texttt{hpmodes}-Format}
\label{sec:normalmoden}

Eine DFT-Frequenzrechnung in Gaussian-16 liefert nach erfolgreicher
Geometrieoptimierung die zweite Ableitung der Energie nach den
kartesischen Atomkoordinaten -- die mass-gewichtete Hesse-Matrix
$\mathbf{F}$. Diagonalisierung von $\mathbf{F}$ liefert die
Eigenvektoren $\mathbf{L}$ und Eigenwerte $\omega_l^2$ der
Normalschwingungen \cite{WilsonDeciusCross1955}:

\begin{equation}
  \mathbf{F}\,\mathbf{L}_l = \omega_l^2 \, \mathbf{L}_l ,
  \qquad l = 1, \ldots, 3N-6
  \label{eq:fg-eigenproblem}
\end{equation}

mit $N$ als Zahl der Atome im Cluster (genauer: dem optimierten
Modell-Bereich). Die kartesischen Verschiebungen pro Atom $i$ in
Mode $l$ erh\"alt man durch Demassengewichtung
\cite{WilsonDeciusCross1955}:

\begin{equation}
  \evec_{l,i} \;=\; \frac{1}{\sqrt{m_i}}\,\mathbf{L}_{l,i} ,
  \label{eq:demass}
\end{equation}

wobei $m_i$ die atomare Masse von Atom $i$ ist.

\subsubsection*{HP-Block versus Standard-Block}

Mit dem Schl\"usselwort \texttt{freq=hpmodes} schreibt Gaussian die
Eigenvektoren in einem zweiten Block (\emph{High-Precision Modes}) mit
f\"unf Nachkommastellen aus. Der \emph{Standard-Block} dagegen
verwendet zwei Nachkommastellen und ist f\"ur sehr kleine
Verschiebungen (z.\,B.\ Modes unter 100~cm$^{-1}$ in gro\ss en
Systemen) unzureichend.

\texttt{modenanalyse\_2fe2s} liest beide Bl\"ocke. Die HP-Eigenvektoren
werden bevorzugt; nur wenn ein einzelner HP-Block fehlt (selten, z.\,B.\
bei Versionsmischung im Logfile), greift der Reader auf den
Standard-Block zur\"uck. Diese Fallback-Logik ist in
\texttt{logio.read\_blocks} implementiert.

\subsubsection*{Konsistenzcheck zwischen HP und Standard}
\label{sec:hp-std-check}

Beide Bl\"ocke geben die Frequenzen $\omega_l$ unabh\"angig aus.
\texttt{modenanalyse\_2fe2s} pr\"uft beim Lauf, ob die HP- und
Standard-Frequenzen f\"ur dieselben Modennummern \"ubereinstimmen
(\texttt{logio.check\_hp\_std\_frequency\_consistency}, Toleranz
$\Delta\omega \leq 0{,}01$~cm$^{-1}$). Eine Abweichung deutet auf:

\begin{itemize}
  \item einen unvollst\"andig geschriebenen Logfile (Job abgebrochen),
  \item einen Mismatch zwischen Lese-Offset und tats\"achlichem
        Block-Anfang (z.\,B.\ wenn ein Job mit \texttt{Restart} fortgesetzt
        wurde),
  \item oder eine Verwechslung zweier Logfiles.
\end{itemize}

Konsistenzfehler werden als Warnung in den Befund geschrieben; das
Programm l\"auft weiter, kennzeichnet die betroffenen Modes aber im
Befund.

\subsubsection*{Skalierung der Eigenvektoren}

F\"ur alle nachfolgenden Auswertungen werden die rohen
Gaussian-Eigenvektoren mit der RMS-Amplitude
$\urms(\omega_l, m_{\mathrm{red}}, T)$ skaliert
(\texttt{core.compute\_thermal\_amplitude})
\cite{Achterhold2002,WilsonDeciusCross1955}:

\begin{equation}
  \urms(\omega, m_{\mathrm{red}}, T) \;=\;
  \sqrt{\frac{\hbar}{2\,m_{\mathrm{red}}\,\omega} \,
        \coth\!\left(\frac{\hbar\omega}{2\kB T}\right)} .
  \codref{core.compute\_thermal\_amplitude}
  \label{eq:urms}
\end{equation}

$m_{\mathrm{red}}$ ist die reduzierte Masse, die Gaussian f\"ur jede
Mode mitliefert. Bei $T=0$ reduziert sich Gleichung~(\ref{eq:urms}) zur
Nullpunkts-Amplitude
$\urms = \sqrt{\hbar / (2 m_{\mathrm{red}} \omega)}$.

Die mit $\urms$ multiplizierten Eigenvektoren haben Einheit \AA{} und
beschreiben die thermisch (oder bei $T \to 0$ rein quantenmechanisch)
gemittelte Auslenkung pro Mode. Sie sind die Eingabe f\"ur s\"amtliche
geometrischen Auswertungen ab \S\ref{sec:oop-inp}.

% -------------------------------------------------------------------------
\subsection{Unsicherheits-Propagation}
\label{sec:uncertainty}

Jede numerische Observable, die das Werkzeug ausgibt
(\texttt{stretch}, \texttt{bend}, \texttt{kern\_d},
\texttt{amplitude\_mean}, \texttt{hn\_stretch}, die SCSD-Amplituden,
\ldots), wird mit einer Begleitunsicherheit versehen
(\texttt{s\_stretch}, \texttt{s\_bend}, \texttt{s\_kern\_d},
\ldots). Diese werden in erster Ordnung Gauss-propagiert aus zwei
zugrundeliegenden Rauschquellen:
\begin{itemize}\setlength{\itemsep}{0pt}
  \item $\sigma_e := $ \texttt{cfg.sigma\_eigvec} (Standard $5\cdot 10^{-4}$):
        numerisches Rauschen pro Eigenvektor-Komponente, in
        kartesischen Einheiten des unskalierten Gaussian-Eigenvektors.
        Steuert die Schwelle ``unbedeutender Beitrag'' f\"ur
        verschiebungsbasierte Observablen.
  \item $\sigma_c := $ \texttt{cfg.sigma\_coord} (Standard $10^{-3}$~\AA):
        Koordinatenrauschen in \AA{}, f\"ur die SCSD-Referenz.
\end{itemize}

\subsubsection*{Konvention: thermisch skaliertes Sigma}
\label{sec:uncertainty-convention}

Alle geometrischen Observablen verwenden den \emph{thermisch
skalierten} Eigenvektor $\evec_{l,i}\,\urms(\omega_l)$ mit
Einheit \AA{}. Das begleitende Sigma wird mit derselben Skalierung
propagiert:
\begin{equation}
  \sigma_e^\text{therm}(l) \;=\; \sigma_e \cdot \urms(\omega_l),
  \codref{core.analyze\_his\_hn, core.analyze\_fe\_ligand}
  \label{eq:sigma-therm}
\end{equation}
so dass das dimensionslose \emph{Signifikanzverh\"altnis}
$|\text{Observable}|\,/\,\sigma$ unter der Wahl von $\urms$ invariant
ist. Diese Konvention ist essenziell f\"ur die Wert/Sigma-Verh\"altnisse,
die in der Farbkodierung der Excel-Bl\"atter verwendet werden.

\paragraph{Konkrete Beispiele.} F\"ur eine typische Niederfrequenz-Mode
mit $\urms \approx 0{,}05$~\AA{}:
\begin{itemize}\setlength{\itemsep}{0pt}
  \item Fe-X-Stretch ($i = $ Ligand, $j = $ Fe):
        $|\evec_i - \evec_j|$ ist eine Differenz zweier unabh\"angiger
        Gauss-Rauschkomponenten, also
        $\sigma(\text{stretch}) = \sigma_e^\text{therm}\sqrt{2}
        \approx 3{,}5\cdot 10^{-5}$~\AA.
  \item Bend INP/OOP: gleicher Faktor $\sqrt{2}$ (Pythagor\"aische
        Zerlegung \"andert die Rauschausbreitung entlang einer
        einzelnen Achse nicht), also $\sigma_\text{bend\_inp}
        = \sigma_\text{bend\_oop} = \sigma_e^\text{therm}\sqrt{2}$.
  \item N-H-Stretch (\texttt{analyze\_his\_hn},
        \texttt{include\_hn\_vibration}) folgt derselben Konvention:
        $\sigma_\text{hn} = \sigma_e^\text{therm}\sqrt{2}$.
  \item SSE-Element-Amplitude:
        $\sigma_\text{amp} = \sigma_e^\text{therm}\sqrt{N_\text{atoms}}$,
        wobei das $\sqrt{N}$-Wachstum aus der Mittelung \"uber
        $N_\text{atoms}$ im Sekund\"arstruktur-Element kommt.
  \item SCSD-Amplituden: getrennte Sigmas f\"ur ``Referenz'' und
        ``Verzerrung''. Die Referenz ist die statische
        Geometrie-Unsicherheit $\sigma_c\sqrt{2}$, die Verzerrung
        kombiniert Referenz- und Thermisches Rauschen:
        $\sigma_\text{SCSD,dist}
         = \sqrt{2}\sqrt{\sigma_c^2 + (\sigma_e \cdot \urms)^2}$.
\end{itemize}

\paragraph{\"Anderung in v1.0.4.} Drei dieser Sigma-Propagationen
waren zuvor inkonsistent implementiert
(\S\ref{sec:migration-v104}): dem N-H-Stretch fehlte die thermische
Skalierung, das SSE-Element verwendete eine fest verdrahtete Konstante
statt \texttt{cfg.sigma\_eigvec}, und SCSD verwendete fest verdrahtete
Referenzkonstanten statt \texttt{cfg.sigma\_coord} und
\texttt{cfg.sigma\_eigvec}. Mit v1.0.4 respektieren alle Sigmas
Gleichung~\ref{eq:sigma-therm} und die TOML-Einstellungen
\texttt{cfg.sigma\_eigvec}/\texttt{sigma\_coord}.

% -------------------------------------------------------------------------
\subsection{Migrationshinweise zu v1.0.4}
\label{sec:migration-v104}

Version 1.0.4 ist ein Bugfix-Release, das die \emph{ausgegebenen
Unsicherheiten} einiger Observablen ver\"andert, w\"ahrend die
zugrundeliegenden \emph{Werte} unver\"andert bleiben. Falls Sie
v1.0.2- oder v1.0.3-Ergebnisse besitzen (z.\,B.\ f\"ur ein
Paper oder eine Dissertation), bitte vor dem Ziehen von
Signifikanz-Schlussfolgerungen diesen Abschnitt lesen.

\paragraph{Was sich nicht ge\"andert hat.} Alle geometrischen und
energetischen Observablen -- \texttt{stretch}, \texttt{bend},
\texttt{kern\_d}, \texttt{lambda\_X}, SCSD-Amplituden,
\texttt{hn\_stretch}, SSE-Element \texttt{amplitude\_mean},
Einbettungs-Koordinaten -- liefern \emph{numerisch identische} Werte
\"uber v1.0.2, v1.0.3 und v1.0.4 f\"ur dasselbe Eingabe-Log.
Keine physikalische Interpretation bestehender Werte wird
invalidiert.

\paragraph{Was sich ge\"andert hat.}

\begin{itemize}\setlength{\itemsep}{2pt}
  \item \texttt{s\_hn\_stretch}: skaliert jetzt mit $\urms$
        (Gl.~\ref{eq:sigma-therm}). F\"ur typische Niederfrequenz-Moden
        ($\urms \sim 0{,}05$~\AA) macht das \texttt{s\_hn\_stretch}
        etwa \textbf{20\,$\times$ kleiner} als in v1.0.2/v1.0.3.
        Konsequenz: viele N-H-Stretches, die unter der alten
        Konvention ``unter dem Rauschen'' erschienen, sind nun
        korrekt als signifikante PCET-Signale identifizierbar.
  \item \texttt{s\_amplitude\_mean} / \texttt{s\_axial\_amplitude} /
        \texttt{s\_tilting\_angle} \ldots\ in
        \texttt{analyze\_sse\_element}: vor v1.0.4 wurde eine fest
        verdrahtete Konstante $1\cdot 10^{-4}$ statt
        \texttt{cfg.sigma\_eigvec} verwendet. Mit dem Standard
        $\sigma_e = 5\cdot 10^{-4}$ sind die neuen Sigmas
        etwa \textbf{5\,$\times$ gr\"o\ss er}. Reklassifikation:
        SSE-Element-Eintr\"age, die unter dem alten Code als
        ``signifikant'' markiert waren, k\"onnen nun als
        ``marginal'' erscheinen.
  \item \texttt{s\_SCSD\_*}: vor v1.0.4 wurden fest verdrahtete
        Konstanten $5\cdot 10^{-7}$ / $5\cdot 10^{-6}$ statt
        \texttt{cfg.sigma\_coord} und \texttt{cfg.sigma\_eigvec}
        verwendet. Neue Sigmas sind etwa \textbf{1000\,$\times$
        gr\"o\ss er}; die vorigen Werte waren physikalisch
        unplausibel.
  \item Fenstergrenzen-Behandlung: der Multi-Window-Modus
        (\texttt{freq\_windows}) verwendete \emph{abgeschlossen-abgeschlossene}
        Intervalle $[lo, hi]$, sodass eine Mode mit Frequenz exakt
        auf einer Grenze in \emph{zwei} benachbarten Fenstern
        erschien. Ab v1.0.4 sind die Grenzen \emph{halb-offen}
        $[lo, hi)$, wobei das letzte Fenster $[lo, hi]$ bleibt.
        Praktische Auswirkung: keine bei typischen Float-Frequenz-
        Rechnungen, aber $B$-Faktor und
        $\Lambda(\omega_\text{cut})$ sind nun garantiert ohne
        Doppelz\"ahlung.
  \item Defensive Warnmeldungen: einige interne Lookup-Fehler
        (Fe-Ligand $c2l$-Miss, His-H $c2l$-Miss, fehlender
        PCET-Acceptor-Index, Null-Zeilen in den Excel-Bl\"attern,
        mehrdeutige PDB-Gaussian-Atomzuordnungen) waren zuvor
        stille Durchl\"aufe. Mit v1.0.4 wird f\"ur jeden dieser
        F\"alle bei seinem ersten Auftreten in einem Lauf ein
        \texttt{UserWarning} ausgegeben (pro Session dedupliziert).
        Falls Sie in einem zuvor stillen Log neue Warnungen sehen,
        pr\"ufen Sie die betroffenen Liganden- oder
        Atomzuordnungen, bevor Sie den entsprechenden Excel-Zeilen
        vertrauen.
  \item Interpolations-Randanker (FUND~12, vom Anwender gemeldet):
        wenn \texttt{freq\_max} gesetzt ist und die n\"achste Mode
        oberhalb mehr als \texttt{interp\_context\_cm1} (Standard
        $5$~cm$^{-1}$) entfernt liegt, wurde vor v1.0.4 keine
        Kontextmode geladen und die interpolierte pDOS fiel knapp
        unterhalb \texttt{freq\_max} steil ab. Mit v1.0.4 wird
        genau \emph{eine} zus\"atzliche Mode -- die n\"achste
        oberhalb \texttt{freq\_max} -- als Fallback-Anker geladen,
        sodass $\texttt{np.interp}$ unabh\"angig von der Modendichte
        immer einen Randanker hat. Derselbe Fallback wirkt
        symmetrisch bei \texttt{freq\_min}. Das Run-Log dokumentiert,
        wann der Fallback ausgel\"ost wurde.
  \item Synthetischer Null-Anker am oberen Rand (FUND~13): wenn
        keine reale Mode oberhalb \texttt{freq\_max} existiert
        (das DFT-Spektrum endet innerhalb des Analysebereichs),
        f\"ugt v1.0.4 eine \emph{synthetische Null-Mode} bei
        $\texttt{freq\_max} + \texttt{interp\_context\_cm1}$ mit
        allen Observablen $= 0$ ein. Das ergibt einen glatten Abfall
        der interpolierten pDOS \"uber die Pufferzone statt einer
        Stufe am oberen Rand. Die synthetische Mode ist mit
        \texttt{mode\_type = "synthetic\_zero"} und
        \texttt{precision = "synthetic"} markiert und tr\"agt per
        Konstruktion null zu $B$-Faktoren,
        $\Lambda_X^\text{total}$ und den Modulationsspektren bei.
\end{itemize}

\paragraph{Post-Release-Audit-Patches (Mai 2026).} Nach dem
v1.0.4-Release identifizierte ein Audit eines realen
[2Fe-2S]-Cys$_2$His$_2$-Batch-Runs (Rieske-Typ in beiden
Protonierungszust\"anden) vier zus\"atzliche Warn- und Reporting-Probleme. Keines beeinflusst
numerische Ausgaben. Vollst\"andige Beschreibungen sind im Supplement
(FUND 14--17); in K\"urze:

\begin{itemize}\setlength{\itemsep}{2pt}
  \item FUND 14: Die Warnung
        ``\texttt{His\_HN NICHT erkannt}'' wurde pro deprotoniertem
        His \emph{zweimal} statt einmal ausgegeben. Behoben durch
        Verschiebung der Warnung aus der inneren
        PDB-vs-Gaussian-Fallback-Schleife heraus.
  \item FUND 15: F\"ur Systeme mit deprotonierten His-Liganden
        (z.\,B.\ Rieske-Typ Cys$_2$His$_2$ im deprotonierten
        Zustand oder mitoNEET-Typ Cys$_3$His nach Deprotonierung) sagte die PCET-Disabled-
        Meldung ``\texttt{no His ligands at cluster}'', obwohl
        \texttt{Fe\_N\_His}-Bl\"atter und \texttt{Lambda\_FeN}
        klar belegten, dass His-Liganden vorhanden waren. Der
        Entscheidungstext unterscheidet jetzt zwischen ``keine His
        \"uberhaupt'' und ``His vorhanden, aber alle deprotoniert''.
  \item FUND 16: Die Warnung
        ``\texttt{interp\_boundary\_mode='context' but no context
        modes available}'' wurde auch bei Full-Spectrum-Runs
        ausgegeben (\texttt{freq\_min}/\texttt{freq\_max} beide
        \texttt{None}), wo gar kein Context-Loading versucht wird.
        Die Warnung wird jetzt nur emittiert, wenn der Benutzer
        explizit ein Fenster gesetzt hat.
  \item FUND 17: REPORT.txt zeigte \texttt{freq\_filter: - - -
        cm-1}, wenn nur \texttt{freq\_windows} gesetzt war, und
        versteckte damit den aktiven Filter vom Audit-Trail. Der
        Report emittiert jetzt eine zus\"atzliche
        \texttt{freq\_windows: [...]}-Zeile.
\end{itemize}

\paragraph{Re-Berechnung vorhandener Daten.} Empfehlung: f\"uhren Sie
vorhandene TOML-Konfigurationen unter v1.0.4 erneut aus und
vergleichen Sie die \texttt{s\_*}-Spalten von
\texttt{Fe\_S\_Cys.xlsx}, \texttt{His\_HN.xlsx}, \texttt{SCSD.xlsx}
und die SSE-Element-Bl\"atter. Die Werte sind identisch; nur Sigmas
und die Signifikanz-F\"arbung \"andern sich.

% -------------------------------------------------------------------------
\subsection{OOP/INP-Klassifikation}
\label{sec:oop-inp}

Eine zentrale Frage f\"ur jeden Cluster-Mode lautet: \emph{Wie viel
Anteil der Schwingung steht senkrecht zur Cluster-Ebene
(out-of-plane, OOP), und wie viel liegt in der Ebene (in-plane, INP)?}
Reine Wippbewegungen der vier Cluster-Atome (Umbrella, Hinge) sind
$\OOP$-dominiert, w\"ahrend Streckmoden (Fe-Fe, S-S, Atmen) in der
Ebene leben.

\subsubsection*{Definition der Cluster-Normalen}

Sei $\rvec_i$ ($i \in \{\Fei_1, \Fei_2, \mathrm{S}_1, \mathrm{S}_2\}$)
die Gleichgewichtsposition der vier Cluster-Atome. Die Normalen-
Richtung der Ebene ist die Richtung, die die Summe der quadrierten
Senkrechtkomponenten minimiert. Konkret aus der Singul\"arwertzerlegung
der zentrierten Cluster-Koordinaten
$\tilde{\rvec}_i = \rvec_i - \bar{\rvec}$ \cite{Kabsch1976}:

\begin{equation}
  \tilde{\mathbf{R}} \;=\; \mathbf{U}\,\boldsymbol{\Sigma}\,\mathbf{V}^\top ,
  \qquad
  \nhat \;=\; \mathbf{V}_{:,3}\,.
  \codref{geometry.cluster\_normal}
  \label{eq:cluster-normal-svd}
\end{equation}

Im Code ist das \texttt{geometry.cluster\_normal}.
Per Konvention zeigt $\nhat$ in $+\zhat$-Richtung des Gaussian-Frames,
sodass \"uber alle Modes hinweg ein einheitliches Vorzeichen entsteht.

\subsubsection*{Anteil der Mode senkrecht zur Ebene}

F\"ur Mode $l$ und ein Atom $i$ ist die OOP-Komponente die Projektion
der Verschiebung auf $\nhat$ (Standard-Zerlegung, vgl.\
\cite{Gee2021}):

\begin{equation}
  \evec_{l,i}^{\OOP} \;=\; (\evec_{l,i} \cdot \nhat)\, \nhat ,
  \qquad
  \evec_{l,i}^{\INP} \;=\; \evec_{l,i} - \evec_{l,i}^{\OOP} .
  \label{eq:oop-projektion}
\end{equation}

Der \emph{OOP-Anteil} der Mode auf einer Atomgruppe $G$ ist das
Verh\"altnis der Quadratsummen \cite{Gee2021}:

\begin{equation}
  \alpha_{\OOP}(G; l) \;=\;
  \frac{\sum_{i \in G} \lVert\evec_{l,i}^{\OOP}\rVert^2}
       {\sum_{i \in G} \lVert\evec_{l,i}\rVert^2} \,\cdot\, 100\,\%
  \,.
  \codref{core.classify\_oop\_inp}
  \label{eq:oop-anteil}
\end{equation}

\begin{infobox}
\textbf{Hinweis zur physikalischen Interpretation.}
Gaussians \texttt{hpmodes}-Block schreibt
\emph{kartesische Verschiebungs-Eigenvektoren}~$\evec_{l,i}$ (bereits
massen-entzogen via Gl.~\ref{eq:demass}), normiert so, dass
$\sum_i \lVert\evec_{l,i}\rVert^2 \approx 1$ statt der mass-gewichteten
Norm $\sum_i m_i \lVert\evec_{l,i}\rVert^2$.
Gleichung~\ref{eq:oop-anteil} misst daher den
\emph{geometrischen kartesischen Out-of-Plane-Anteil}: sie sagt, welcher
Anteil der quadrierten kartesischen Verschiebungsamplituden aus der
Cluster-Ebene zeigt. Sie entspricht \emph{nicht} dem
kinetischen-Energie-Anteil der senkrecht zur Ebene flie\ss enden Mode,
welcher die mass-gewichtete Form
$\sum_i m_i \lVert\evec_{l,i}^{\OOP}\rVert^2 / \sum_i m_i \lVert\evec_{l,i}\rVert^2$
erfordern w\"urde. F\"ur einen [2Fe-2S]-Cluster sind beide Gr\"o\ss en
numerisch \"ahnlich (nur Fe und S tragen bei, mit vergleichbaren
Massen), aber die Unterscheidung ist wichtig beim Vergleich von
$\alpha_{\OOP}$-Werten zwischen chemisch verschiedenen Systemen.
\end{infobox}

\subsubsection*{Bin\"are und feine Klassifikation}

Klassisch wird $\alpha_{\OOP}(G_{\text{Cluster}}; l)$ einer
Schwellenpr\"ufung unterzogen:

\begin{equation}
  \text{mode\_type}(l) \;=\;
  \begin{cases}
    \OOP, & \alpha_{\OOP} \geq \theta_{\text{OOP}} \\
    \INP, & \alpha_{\OOP} \leq 100\,\% - \theta_{\text{OOP}} \\
    \text{Gemischt}, & \text{sonst}
  \end{cases}
\end{equation}

mit Default $\theta_{\text{OOP}} = 60\,\%$ (per Config justierbar via
\texttt{cfg.mode\_type\_threshold}).

Zus\"atzlich existiert eine \emph{feine Klassifikation} mit drei
Schwellen $(\theta_1, \theta_2, \theta_3)$
(\texttt{core.classify\_oop\_inp}). Default $(60, 75, 90)$:

\begin{center}
\begin{tabularx}{\textwidth}{l X}
\toprule
\textbf{Stufe} & \textbf{Definition} \\
\midrule
\textbf{Pur OOP}        & $\alpha_{\OOP} \geq \theta_3 = 90\,\%$ \\
\textbf{Stark OOP}      & $\theta_2 \leq \alpha_{\OOP} < \theta_3$, also $75-90\,\%$ \\
\textbf{Mehrheitlich OOP} & $\theta_1 \leq \alpha_{\OOP} < \theta_2$, also $60-75\,\%$ \\
\textbf{Gemischt}       & $|{\alpha_{\OOP} - 50\,\%}| < \theta_1 - 50\,\%$ \\
\textbf{Mehrheitlich INP} & $100\,\% - \theta_2 < \alpha_{\OOP} \leq 100\,\% - \theta_1$, also $25-40\,\%$ \\
\textbf{Stark INP}      & $100\,\% - \theta_3 < \alpha_{\OOP} \leq 100\,\% - \theta_2$, also $10-25\,\%$ \\
\textbf{Pur INP}        & $\alpha_{\OOP} \leq 100\,\% - \theta_3 = 10\,\%$ \\
\bottomrule
\end{tabularx}
\end{center}

Die feine Klassifikation gibt im Excel-Sheet eine pr\"azisere Aussage
\"uber den Charakter einer Mode -- insbesondere im OOP-Hinge-Bereich,
wo eine bin\"are 60\,\%-Schwelle entweder zu permissiv ist (Modes
zwischen 60 und 75\,\% sehen sehr unterschiedlich aus) oder eine
unscharfe Antwort liefert.

\begin{infobox}
\textbf{Welche Schwelle f\"ur welche Frage?} F\"ur den Vergleich mit
NRVS-Spektren in Polarisationsgeometrie liefert die bin\"are 60\,\%-Schwelle
gen\"ugend Trennsch\"arfe. F\"ur Manuskript-Tabellen mit
Mode-Eigenschaften (z.\,B.\ Symmetrie-Kandidaten f\"ur D$_\mathrm{2h}$)
ist die feine 7-Stufen-Klassifikation aussagekr\"aftiger; mit ihr
lassen sich die Modes "`pur"' OOP / INP zuverl\"assig identifizieren.
\end{infobox}

\subsubsection*{Worked Example: OOP\%-Berechnung an einer Cluster-Atmungsmode}

Zur Veranschaulichung berechnen wir $\alpha_{\OOP}$ explizit f\"ur eine
hypothetische Cluster-Atmungsmode in einem Cys$_4$-Modell-Cluster.

\paragraph{Schritt 1 -- Geometrie und Cluster-Normale.}
Die vier Cluster-Atome sitzen idealisiert bei (Einheit $\angstrom$):
\begin{align*}
  \mathbf{r}_{\text{Fe}_1} &= (+1.50,\, 0.00,\, 0.00) \\
  \mathbf{r}_{\text{Fe}_2} &= (-1.50,\, 0.00,\, 0.00) \\
  \mathbf{r}_{\text{S}_1}  &= ( 0.00,\, +1.10,\, 0.00) \\
  \mathbf{r}_{\text{S}_2}  &= ( 0.00,\, -1.10,\, 0.00)
\end{align*}
Schwerpunkt $\bar{\mathbf{r}} = \mathbf{0}$. Die zentrierten Koordinaten
liegen alle in der $xy$-Ebene; SVD liefert
$\hat{\mathbf{n}} = (0, 0, 1)$ (Cluster-Normale parallel zur $z$-Achse).

\paragraph{Schritt 2 -- Mode-Eigenvektoren (thermisch skaliert, $T=5$\,K).}
Eine reine \textbf{Atmungsmode} streckt/staucht beide Fe-S-Diagonalen
synchron in der Ebene. Ein typischer thermisch skalierter Eigenvektor
sieht so aus (nach $\mathbf{e} \to \mathbf{e}\cdot u_\text{rms}$):
\begin{align*}
  \mathbf{e}_{\text{Fe}_1} &= (+0.012,\, +0.009,\, 0.000) \\
  \mathbf{e}_{\text{Fe}_2} &= (-0.012,\, -0.009,\, 0.000) \\
  \mathbf{e}_{\text{S}_1}  &= ( 0.000,\, -0.018,\, 0.000) \\
  \mathbf{e}_{\text{S}_2}  &= ( 0.000,\, +0.018,\, 0.000)
\end{align*}

\paragraph{Schritt 3 -- OOP-Komponenten extrahieren.}
F\"ur jedes Atom $i$:
$\mathbf{e}_i^\text{OOP} = (\mathbf{e}_i \cdot \hat{\mathbf{n}})\,
\hat{\mathbf{n}}$. Da alle Eigenvektoren $z$-Komponente $0$ haben,
ist $\mathbf{e}_i^\text{OOP} = \mathbf{0}$ f\"ur alle $i$.

\paragraph{Schritt 4 -- $\alpha_{\OOP}$ ausrechnen.}
\begin{align*}
  \sum_{i \in G} \|\mathbf{e}_i^\text{OOP}\|^2 &= 0 \\
  \sum_{i \in G} \|\mathbf{e}_i\|^2 &= (0.012^2 + 0.009^2)\cdot 2 + (0.018^2)\cdot 2 \\
   &= 2\cdot 0.000225 + 2\cdot 0.000324 = 0.001098
\end{align*}
\begin{equation*}
  \alpha_{\OOP} = \frac{0}{0.001098}\cdot 100\,\% = 0\,\%
\end{equation*}
Klassifikation: \textbf{Pur INP} ($\alpha_{\OOP} \leq 10\,\%$). Die
Cluster-Atmung ist eine in-plane-Mode -- wie erwartet.

\paragraph{Schritt 5 -- Vergleich mit Hinge-Mode.}
Eine Cluster-Hinge-Mode (das "`Buchaufklappen"' der zwei Fe-S-Bonds
um eine S-S-Achse) hat dagegen die Form:
\begin{align*}
  \mathbf{e}_{\text{Fe}_1} &= (0.000,\, 0.000,\, +0.020) \\
  \mathbf{e}_{\text{Fe}_2} &= (0.000,\, 0.000,\, +0.020) \\
  \mathbf{e}_{\text{S}_1}  &= (0.000,\, 0.000,\, -0.020) \\
  \mathbf{e}_{\text{S}_2}  &= (0.000,\, 0.000,\, -0.020)
\end{align*}
Hier gilt $\mathbf{e}_i^\text{OOP} = \mathbf{e}_i$ f\"ur alle $i$,
also $\alpha_{\OOP} = 100\,\%$ -- klassifiziert als \textbf{Pur OOP}.

\paragraph{Schritt 6 -- Was im realen DFT-Lauf passiert.}
In einem echten Cys$_4$- oder Cys$_2$His$_2$-Cluster sind reine
Atmungs- und Hinge-Modes selten -- meist mischen sich Komponenten,
und $\alpha_{\OOP}$ liegt zwischen 20 und 80\,\%. Die Klassifikation
nach den Default-Schwellen $(60, 75, 90)$ liefert eine 7-stufige
Aufteilung, die etwa f\"ur einen typischen [2Fe-2S]-Lauf
folgende Verteilung produziert:
\begin{itemize}\setlength{\itemsep}{0pt}
  \item Pur INP: $\sim 30$\,\% der Cluster-Modes
  \item Mehrheitlich/Stark INP: $\sim 35$\,\%
  \item Gemischt: $\sim 20$\,\%
  \item Mehrheitlich/Stark OOP: $\sim 12$\,\%
  \item Pur OOP: $\sim 3$\,\%
\end{itemize}
(Diese Zahlen variieren stark zwischen Systemen -- in starren Cluster-
Geometrien dominieren INP-Modes, in flexiblen Liganden-Sph\"aren
finden sich mehr Gemischte.)

% -------------------------------------------------------------------------
\subsection{Fe-Ligand-Bindungs-Zerlegung: Stretch, Bend INP, Bend OOP}
\label{sec:fe-lig-decomp}

Die OOP/INP-Zerlegung aus \S\ref{sec:oop-inp} gilt f\"ur den
\emph{Cluster-Kern} (Fe$_1$, Fe$_2$, S$_1$, S$_2$). Eine zweite
Zerlegung wird f\"ur jede Fe-Ligand-Bindung
(Fe-S$_\text{Cys}$, Fe-N$_\text{His}$, Fe-O$_\text{Asp/Glu}$)
durchgef\"uhrt: die Relativ\-verschiebung der beiden Bindungspartner
wird in eine \emph{Stretch}-Komponente (entlang der Bindungsachse)
und eine \emph{Bend}-Komponente (senkrecht dazu) zerlegt, wobei der
Bend weiter in INP und OOP relativ zur Cluster-Ebene aufgeteilt wird.

Seien $\evec_\text{Fe}$ und $\evec_\text{Ligand}$ die thermisch
skalierten Verschiebungen des Fe-Atoms und des Ligand-Donor-Atoms,
und sei
\begin{equation}
  \hat{\mathbf{b}} = \frac{\rvec_\text{Ligand} - \rvec_\text{Fe}}
                            {\lVert\rvec_\text{Ligand} - \rvec_\text{Fe}\rVert}
\end{equation}
der Einheitsvektor von Fe zum Ligand-Donor in der Gleichgewichts-
Geometrie. Definiere die Relativ\-verschiebung
\begin{equation}
  \Delta\evec = \evec_\text{Ligand} - \evec_\text{Fe}.
  \label{eq:fe-lig-delta}
\end{equation}

\paragraph{Zerlegung entlang und senkrecht zur Bindung.}
Die vorzeichenbehaftete Stretch-Komponente ist die Projektion auf
$\hat{\mathbf{b}}$, der Bend-Vektor ist der Rest:
\begin{equation}
  \Delta r_\text{stretch} = \Delta\evec \cdot \hat{\mathbf{b}},
  \qquad
  \Delta\evec_\perp = \Delta\evec
                       - \Delta r_\text{stretch}\hat{\mathbf{b}}.
  \label{eq:fe-lig-perp}
\end{equation}

\paragraph{Zerlegung des Bends in INP und OOP.}
Die Cluster-Normale $\hat{\mathbf{n}}$ (\S\ref{sec:oop-inp}) zerlegt
$\Delta\evec_\perp$ weiter in eine Komponente entlang $\hat{\mathbf{n}}$
(aus der Cluster-Ebene) und einen Rest (in der Cluster-Ebene):
\begin{align}
  \Delta r_\text{bend\_oop} &= \Delta\evec_\perp \cdot \hat{\mathbf{n}},
  \qquad\text{(vorzeichenbehaftet)} \\
  \Delta r_\text{bend\_inp} &= \lVert\Delta\evec_\perp
       - \Delta r_\text{bend\_oop}\hat{\mathbf{n}}\rVert.
  \codref{core.analyze\_fe\_ligand}
  \label{eq:fe-lig-bend-split}
\end{align}
Die ausgegebenen \texttt{bend}-, \texttt{bend\_inp}- und
\texttt{bend\_oop}-Spalten in \texttt{Fe\_S\_Cys.xlsx} /
\texttt{Fe\_N\_His.xlsx} sind die \emph{Absolutwerte} (L\"angen)
dieser drei Gr\"o\ss en, alle in \AA{}. Nach Pythagoras gilt
\begin{equation}
  \lVert\Delta\evec_\perp\rVert^2
    = (\Delta r_\text{bend\_inp})^2 + (\Delta r_\text{bend\_oop})^2,
\end{equation}
was der nat\"urlichen Energie-Zerlegung der Bend-Mode entspricht.

\paragraph{Bend-INP/OOP-Prozente.}
F\"ur Vergleiche von Moden mit sehr unterschiedlichen
Gesamt-Amplituden werden die absoluten L\"angen in
\emph{Quadratamplituden-Anteile} relativ zum Gesamt-Bend umgewandelt:
\begin{equation}
  \text{bend\_inp\_pct} = 100 \cdot
     \frac{(\Delta r_\text{bend\_inp})^2}
          {(\Delta r_\text{bend\_inp})^2 + (\Delta r_\text{bend\_oop})^2},
\end{equation}
und analog \texttt{bend\_oop\_pct}. Die beiden summieren sich
exakt zu 100\,\%. F\"ur sehr kleine Gesamt-Bends (unter dem
Rauschniveau $2\sigma_e^\text{therm}\sqrt{2}$,
\S\ref{sec:uncertainty}) werden die Prozente auf NaN gesetzt,
um numerisches Rauschen nicht zu verst\"arken.

\paragraph{Sigma-Propagation.}
Alle drei Gr\"o\ss en (stretch, bend\_inp, bend\_oop) sind
Differenzen zweier unabh\"angiger Gauss-Rauschkomponenten und tragen
daher das Standard-$\sqrt{2}\,\sigma_e^\text{therm}$-Sigma:
\begin{equation}
  \sigma_\text{stretch} = \sigma_\text{bend\_inp}
  = \sigma_\text{bend\_oop}
  = \sqrt{2}\,\sigma_e \cdot \urms(\omega_l).
\end{equation}
Die Sigmas werden in die
\texttt{Fe\_S\_Cys.xlsx}-/\texttt{Fe\_N\_His.xlsx}-Dateien unter
\texttt{s\_stretch}, \texttt{s\_bend\_inp}, \texttt{s\_bend\_oop}
geschrieben und f\"ur die Signifikanz-F\"arbung der Bend-Zellen
verwendet.

\paragraph{Physikalische Interpretation.}
F\"ur PCET-Anwendungen ist die INP/OOP-Aufteilung des
Fe-N(His)-Bends besonders informativ:
\begin{itemize}\setlength{\itemsep}{0pt}
  \item \emph{bend\_inp dominant} ($>$70\,\%): der His-Ring wippt
        in der Cluster-Ebene -- eine Bewegung, die den
        Fe-Fe-Abstand und den Eisen-Schwefel-Kern stark moduliert,
        klassisch ET-aktiv.
  \item \emph{bend\_oop dominant} ($>$70\,\%): der His-Ring kippt
        aus der Ebene heraus, was die N-H-Bindungsrichtung
        moduliert -- Kopplungs-Kanal f\"ur PT zu einem vergrabenen
        Akzeptor.
  \item \emph{gemischt} (30--70\,\%): die Mode tr\"agt gleichzeitig
        zu ET- und PT-Pfaden bei -- typische PCET-Signatur.
\end{itemize}

% -------------------------------------------------------------------------
\subsection{SCSD nach Kingsbury \& Senge}
\label{sec:scsd}

Die OOP/INP-Zerlegung der vorigen Sektion ist eine
\emph{geometrische Heuristik} -- sie kennt nur die Cluster-Ebene, nicht
die Symmetrie. Eine rigorose Symmetriezerlegung leistet die
\emph{Symmetry-Coordinate Structural Decomposition} (SCSD) nach
Kingsbury \& Senge \cite{Kingsbury2024}.

\subsubsection*{Idee}

Ein {[}2Fe-2S{]}-Cluster mit ideal-rhombischer Geometrie geh\"ort zur
Punktgruppe D$_\mathrm{2h}$. Die acht m\"oglichen Verschiebungsmuster
der vier Atome (1 Translation $\times$ 3 Achsen + 1 Rotation $\times$
3 Achsen + 2 Streck) verteilen sich auf die acht D$_\mathrm{2h}$-Irreps:
$\Ag, \Bgi{1}, \Bgi{2}, \Bgi{3}, \Aui, \Bui{1}, \Bui{2}, \Bui{3}$.

Eine reale Cluster-Geometrie hat aber endliche Auslenkungen -- der
Cluster ist nie ideal rhombisch. SCSD projiziert die tats\"achlichen
Atomkoordinaten auf die acht Irreps der Punktgruppe und liefert pro
Irrep eine Beitragsamplitude. Auf eine Normalmode angewandt, sagt
SCSD: \emph{Welche Symmetrie hat diese Auslenkung tats\"achlich?}

\subsubsection*{D$_\mathrm{2h}$-Charaktertafel und Symmetrieoperationen}

Die Punktgruppe D$_\mathrm{2h}$ hat acht Symmetrieoperationen
$\{E, C_2(z), C_2(y), C_2(x), i, \sigma_h, \sigma_v(xz), \sigma_v(yz)\}$
und entsprechend acht eindimensionale Irreps. Die Charaktertafel:

\begin{center}
\begin{tabular}{lccccccccl}
  \toprule
  \textbf{D$_{2h}$} & $E$ & $C_2(z)$ & $C_2(y)$ & $C_2(x)$ &
                     $i$ & $\sigma_h$ & $\sigma_v(xz)$ & $\sigma_v(yz)$ &
                     \textbf{Lineare Basis} \\
  \midrule
  $A_g$    & 1 & 1 & 1 & 1 & 1 & 1 & 1 & 1 & $x^2, y^2, z^2$ \\
  $B_{1g}$ & 1 & 1 & $-$1 & $-$1 & 1 & 1 & $-$1 & $-$1 & $xy$ \\
  $B_{2g}$ & 1 & $-$1 & 1 & $-$1 & 1 & $-$1 & 1 & $-$1 & $xz$ \\
  $B_{3g}$ & 1 & $-$1 & $-$1 & 1 & 1 & $-$1 & $-$1 & 1 & $yz$ \\
  $A_u$    & 1 & 1 & 1 & 1 & $-$1 & $-$1 & $-$1 & $-$1 & --- \\
  $B_{1u}$ & 1 & 1 & $-$1 & $-$1 & $-$1 & $-$1 & 1 & 1 & $z$ \\
  $B_{2u}$ & 1 & $-$1 & 1 & $-$1 & $-$1 & 1 & $-$1 & 1 & $y$ \\
  $B_{3u}$ & 1 & $-$1 & $-$1 & 1 & $-$1 & 1 & 1 & $-$1 & $x$ \\
  \bottomrule
\end{tabular}
\end{center}

Im [2Fe-2S]-Kontext ist die Achsenkonvention (Kingsbury 2024):
$\hat{\mathbf{x}}$ entlang Fe-Fe, $\hat{\mathbf{y}}$ entlang S-S
(Br\"ucken-Schwefel-zu-Br\"ucken-Schwefel),
$\hat{\mathbf{z}} = \hat{\mathbf{n}}$ (Cluster-Normale).

\subsubsection*{Bedeutung der einzelnen Irreps fuer die Cluster-Bewegung}

Direkt aus der Charaktertafel kann man ablesen, welche Bewegungen
welcher Irrep zugeordnet sind:

\begin{center}
\begin{tabular}{lp{8.5cm}}
  \toprule
  \textbf{Irrep} & \textbf{Cluster-Bewegung in [2Fe-2S]-Konvention} \\
  \midrule
  $A_g$    & Symmetrische Streckungen entlang Hauptachsen.
             Beispiele: Fe-Fe-Atmung, S-S-Streckung,
             Cluster-Atmung-Mischung. Klassischer
             "`Reorganisations-Mode"' f\"ur ET. \\
  $B_{1g}$ & In-Plane-Knickbewegung mit gleichzeitiger Verzerrung
             beider Diagonalen. \\
  $B_{2g}$ & In-Plane-Bewegung mit ungleichen Fe-S-Abst\"anden auf
             beiden Cluster-Seiten. \\
  $B_{3g}$ & In-Plane-Bewegung mit Drehung des Rhombus-Mittelpunkts
             relativ zu den Atomen. \\
  $A_u$    & "`Zentral-Inversions-Mode"' -- selten in starren
             Clustern, kann bei stark verzerrten Geometrien
             auftauchen. \\
  $B_{1u}$ & Reine OOP-Verschiebung aller Cluster-Atome in
             $\hat{\mathbf{z}}$-Richtung. Hinge-Mode. \\
  $B_{2u}$ & OOP-Verschiebung mit Twist um die $\hat{\mathbf{x}}$-Achse. \\
  $B_{3u}$ & OOP-Verschiebung mit Twist um die $\hat{\mathbf{y}}$-Achse. \\
  \bottomrule
\end{tabular}
\end{center}

Die $g$-Irreps sind \emph{gerade} (in-plane, da paritäts-erhaltend),
die $u$-Irreps \emph{ungerade} (typischerweise OOP-aktiv).

\subsubsection*{Kanonische Referenz und Vorzeichenwahl}

Die Projektion erfordert eine eindeutige Achsenwahl. \texttt{scsdpy}
verwendet die Kingsbury-kanonische Konvention:

\begin{enumerate}
  \item $\zhat$ ist die Cluster-Normale.
  \item Die Cluster-Hauptachse (l\"angste interatomare Distanz) wird
        auf $\xhat$ gelegt.
  \item Vorzeichen so gew\"ahlt, dass die Cluster-Atome in einer
        bestimmten Reihenfolge auf positive Quadranten kommen.
\end{enumerate}

In \texttt{modenanalyse\_2fe2s} wird die Referenzgeometrie einmal pro
Lauf gebaut und gespeichert (\texttt{core.\_get\_scsd\_model}). Sie ist
unabh\"angig davon, in welcher Orientierung Gaussian den Cluster
optimiert hat.

\subsubsection*{Anwendung auf Normalmodes}

F\"ur Mode $l$ wird das Cluster (vier Atome) um $\evec_{l,i}\,\urms$
ausgelenkt; aus dieser ausgelenkten Geometrie minus der
Referenzgeometrie ergibt sich ein 12-dimensionaler Verschiebungsvektor.
Die SCSD-Projektion liefert acht Amplituden \cite{Kingsbury2024}:

\begin{equation}
  \mathbf{a}_l \;=\;
  \mathbf{P}_{\mathrm{D}_{2\mathrm{h}}}\bigl(\evec_l^{\text{Cluster}} \,\urms\bigr) ,
  \qquad
  \mathbf{a}_l \in \mathbb{R}^8 .
  \codref{core.compute\_scsd\_for\_mode\_full}
  \label{eq:scsd-projection}
\end{equation}

Die Quadratsumme $\sum_\Gamma a_{l,\Gamma}^2$ ist gleich der
ungewichteten Summe \"uber die quadrierten Atom-Verschiebungen --
die SCSD-Zerlegung ist eine \emph{orthogonale} Basistransformation.

\subsubsection*{Dominante Irrep -- rigoros bestimmt}
\label{sec:scsd-rigoros}

Die rigorose Funktion \texttt{core.extract\_dominant\_scsd\_irreps}
bestimmt die dominante Irrep wie folgt. Sie:

\begin{enumerate}
  \item summiert die quadrierten Amplituden pro Irrep,
  \item normiert auf die Gesamtsumme,
  \item liefert \emph{prim\"are} und \emph{sekund\"are} Irrep zur\"uck,
        gemeinsam mit ihren prozentualen Anteilen.
\end{enumerate}

Diese Klassifikation kann widersprechen zur OOP/INP-Heuristik: ein
Mode kann z.\,B.\ "`mehrheitlich OOP"' (Heuristik) und gleichzeitig
"`prim\"ar $\Bgi{1}$"' (SCSD) sein -- $\Bgi{1}$ ist eine spezifische
OOP-Symmetrie, die Heuristik unterscheidet das nicht.

\begin{infobox}
\textbf{Das Sheet \texttt{SCSD} im Excel-Output} listet pro Mode die
acht Amplituden plus die zwei dominanten Irreps mit Prozentangabe. Die
Modes mit konsistenter ($>$80\,\%) prim\"arer Irrep sind diejenigen,
die in der Symmetrie-Tabelle des Manuskripts gef\"uhrt werden sollten.
\end{infobox}

\subsubsection*{Wann SCSD ausgeschaltet bleibt}

Die SCSD-Analyse ben\"otigt das externe Paket \texttt{scsdpy}. Wenn
es nicht installiert ist, wird das SCSD-Sheet \"ubersprungen und im
Befund wird ein Hinweis ausgegeben. Alle anderen Auswertungen
(OOP/INP-Heuristik, Kern-Klassifikation, NIS) sind davon nicht
betroffen.

% -------------------------------------------------------------------------
\subsection{Heuristische Kern-Klassifikation}
\label{sec:kern-klassifikation}

Neben der rigorosen SCSD-Zerlegung gibt es eine \emph{heuristische
Kern-Klassifikation}, die jede Mode anhand ihres geometrischen
Bewegungsmusters einem Bewegungstyp zuordnet
(\texttt{core.classify\_kernel\_mode\_from\_evg}). Sie ist robust auch
ohne \texttt{scsdpy} und bietet eine intuitive Kategorisierung. Die
Heuristik ist explizit als solche ausgewiesen und nicht zu verwechseln
mit der rigorosen SCSD-Antwort.

\subsubsection*{Bewegungs-Scores}

Der heuristische Klassifikator berechnet f\"ur jede Mode acht
Bewegungs-Scores:

\begin{center}
\begin{tabularx}{\textwidth}{l X}
\toprule
\textbf{Score} & \textbf{Definition} \\
\midrule
\textbf{Translation} & Schwerpunktverschiebung des Clusters
                       (sollte bei intramolekularer Mode klein sein) \\
\textbf{Umbrella}    & Mittlere OOP-Auslenkung aller vier Cluster-Atome
                       (alle gemeinsam aus der Ebene heraus) \\
\textbf{Hinge-Fe}    & OOP-Differenz Fe$_1$ minus Fe$_2$ (Fe-Atome
                       schaukeln gegeneinander) \\
\textbf{Hinge-S}     & OOP-Differenz S$_1$ minus S$_2$ (S-Atome
                       schaukeln gegeneinander) \\
\textbf{OOP-Twist}   & Kombination beider Hinge-Modi (Verdrillung) \\
\textbf{Fe-Streckung} & In-plane-Anteil entlang Fe-Fe-Achse \\
\textbf{S-Streckung} & In-plane-Anteil entlang S-S-Achse \\
\textbf{Breathing}   & Atemschwingung: alle vier Atome simultan vom
                       Schwerpunkt weg / hin zu ihm \\
\textbf{Rhombus-Scherung} & Anti-symmetrische Verzerrung des Rhombus
                       (Fe-S-Bindungswinkel) \\
\textbf{Rotation-ip} & Drehbewegung des Clusters in seiner Ebene \\
\bottomrule
\end{tabularx}
\end{center}

Jeder Score ist auf die Gesamtverschiebung der vier Cluster-Atome
normiert und liegt in $[0, 1]$.

\subsubsection*{Klassifikation}

Die zwei h\"ochsten Scores werden als \emph{prim\"ares} und
\emph{sekund\"ares} Bewegungsmuster zur\"uckgegeben. Im Excel-Sheet
\texttt{Kern\_Scores} sind alle acht Scores plus die beiden Top-Labels
gelistet. Diese Heuristik gibt eine Vor-Sortierung, die
\textit{ohne} D$_\mathrm{2h}$-Symmetrieannahme funktioniert (auch in
verzerrten Geometrien).

\subsubsection*{Heuristik vs.\ SCSD: Wann widersprechen sie sich?}

In ideal-rhombischen Clustern stimmen beide \"uberein:

\begin{center}
\begin{tabularx}{\textwidth}{l l X}
\toprule
\textbf{Heuristik} & \textbf{SCSD-Irrep} & \textbf{Bemerkung} \\
\midrule
Umbrella     & $\Bgi{2}$ oder $\Bgi{3}$ & abh\"angig von Definition der Achsen \\
Hinge-Fe     & $\Aui$           & antisymmetrisch zur Cluster-Mitte \\
Breathing    & $\Ag$            & vollsymmetrisch \\
Fe-Streckung & $\Bgi{1}$        & antisymmetrisch entlang Fe-Fe \\
Rotation-ip  & $\Bgi{1}$ (in-plane) & abh\"angig von Achsenwahl \\
\bottomrule
\end{tabularx}
\end{center}

In real verzerrten Clustern, wie sie in Proteinen vorkommen, teilen
sich die Bewegungen oft auf zwei Heuristik-Scores auf
(z.\,B.\ ein Mode mit Hinge-Fe \emph{und} Rhombus-Scherung etwa
gleicher Gr\"o\ss e). Hier ist die SCSD-Antwort schl\"ussiger: sie
ordnet eine eindeutige Symmetrie zu, die nicht von der Achsenwahl
abh\"angt.

\begin{warnbox}
\textbf{Achtung -- nicht verwechseln.} Die Heuristik-Labels (Umbrella,
Hinge, Breathing, ...) sind \emph{Bewegungsbezeichner}, keine
Symmetrie-Labels. Ein "`Hinge-Fe"'-Mode kann je nach Achsenwahl
Symmetrie $\Aui$, $\Bui{1}$ oder $\Bui{2}$ tragen. F\"ur
Manuskript-Aussagen \"uber die Symmetrie immer die SCSD-Antwort
verwenden, nicht die Heuristik.
\end{warnbox}
% =========================================================================
% Anhang an: modenanalyse_handbuch.tex
% Sections 2.5 - 2.8 + Bibliographie
% =========================================================================

% -------------------------------------------------------------------------
\subsection{Embeddings: UMAP und HDBSCAN}
\label{sec:embeddings}

Bei einem [2Fe--2S]-System mit 5\,000--19\,000 Modes ist eine
gleichzeitige Darstellung aller Mode-Eigenschaften auf
zweidimensionalen Plots nicht m\"oglich. \texttt{modenanalyse\_2fe2s}
nutzt UMAP (Uniform Manifold Approximation and Projection) als
Dimensionsreduktion, optional gefolgt von HDBSCAN-Clustering.

\subsubsection*{Wahl der Methode: warum UMAP?}

Im Vergleich zu Alternativen (PCA, t-SNE) hat sich UMAP nach
Test-L\"aufen auf protonierten und deprotonierten
Cys$_2$His$_2$-Modellsystemen sowie einer HHCC-Mutante als die
einzige Methode erwiesen, die f\"ur DFT-Mode-Daten konsistent
interpretierbare Ergebnisse liefert.

\begin{itemize}
  \item \textbf{PCA} \cite{Pearson1901,Hotelling1933,Jolliffe2002} ist
        linear und schnell, kann aber nichtlineare Mannigfaltigkeiten
        nicht aufl\"osen. Bei DFT-Modes liegen die Datenpunkte typisch
        auf gekr\"ummten Untermannigfaltigkeiten im Featureraum
        (Frequenz-OOP-INP-Score-Beziehungen sind nichtlinear); PCA
        zeigt nur eine Schatten-Projektion. PC1 liefert in unseren
        Daten typischerweise nur 25--35\,\% Varianzerkl\"arung -- zu
        wenig f\"ur eine eigenst\"andige Visualisierung.

  \item \textbf{t-SNE} \cite{vanDerMaaten2008,vanDerMaaten2014} bewahrt
        lokale Nachbarschaften, aber globale Distanzen sind
        bedeutungslos: ein scheinbar gro\ss er Abstand zwischen zwei
        Cluster-Regionen in der t-SNE-Darstellung entspricht nicht
        notwendigerweise einem gro\ss en Abstand im Featureraum
        \cite{vanDerMaaten2008}. Au\ss erdem ist t-SNE stochastisch,
        die Ergebnisse h\"angen vom Zufalls-Seed ab. F\"ur Cluster-
        Identifikation l\"asst t-SNE Cluster oft auch dann erscheinen,
        wenn die zugrundeliegenden Daten kontinuierlich verteilt sind.

  \item \textbf{UMAP} \cite{McInnes2018,Becht2019} bewahrt sowohl
        lokale Nachbarschaft als auch globale Topologie auf
        kontrollierbare Weise und ist deutlich schneller als t-SNE
        bei vergleichbarer Qualit\"at \cite{Becht2019}. Der
        \texttt{n\_neighbors}-Parameter erlaubt eine explizite
        Steuerung der Lokalit\"at-Globalit\"ats-Balance, was f\"ur die
        Mode-Analyse essenziell ist (kleine n\_neighbors bevorzugen
        feinere Symmetrie-Klassen, gro\ss e n\_neighbors die globale
        Frequenzstruktur).
\end{itemize}

Eine PCA-Anwender-Sicht (Varianz-Erkl\"arung pro Achse) ist hier nicht
mehr erforderlich, weil die Feature-Bedeutung schon vor dem Embedding
durch die kuratierte Auswahl (siehe Abschnitt \emph{Feature-Matrix})
explizit ist. Die anderen Sheets im Excel-Output
(\texttt{Modenanalyse}, \texttt{Reorganisationsenergie},
\texttt{Reorg\_Total}) liefern die quantitative Information, die PCA
zus\"atzlich beigetragen h\"atte.

\subsubsection*{UMAP: mathematische Grundlagen}

UMAP konstruiert eine niederdimensionale Einbettung in zwei Schritten,
basierend auf der Theorie topologischer Mannigfaltigkeiten und
unscharfer simplizialer Mengen \cite{McInnes2018}.

\paragraph{Schritt 1: Hochdimensionaler Nachbarschaftsgraph.}

Im urspr\"unglichen Featureraum $\mathbb{R}^p$ (mit $p = 31$ Features
) wird f\"ur jeden Datenpunkt $x_i$ ein lokaler
Nachbarschaftsgraph konstruiert. Pro Punkt werden die $k$ n\"achsten
Nachbarn bestimmt (Parameter \texttt{n\_neighbors}). Die Distanz zu
diesen Nachbarn wird mittels einer adaptiven Metrik
$d_i(x_i, x_j)$ in einen Wahrscheinlichkeitswert $\mu_{ij} \in [0,1]$
umgerechnet:
\begin{equation}
  \mu_{ij} = \exp\!\left(\frac{-\max(0,\,
              d(x_i, x_j) - \rho_i)}{\sigma_i}\right),
  \label{eq:umap-mu}
\end{equation}
wobei $\rho_i$ der Abstand zum n\"achsten Nachbarn von $x_i$ ist
(garantiert, dass jeder Punkt mit Wahrscheinlichkeit 1 zu seinem
unmittelbaren Nachbarn verbunden ist) und $\sigma_i$ so gew\"ahlt wird,
dass $\sum_j \mu_{ij} = \log_2(k)$ erf\"ullt ist (lokale
Skalen-Adaption).

Die Asymmetrie in $\mu_{ij}$ wird durch eine \emph{unscharfe Vereinigung}
zu einer symmetrischen Adjazenzmatrix $w_{ij}$ kombiniert
\cite{McInnes2018}:
\begin{equation}
  w_{ij} = \mu_{ij} + \mu_{ji} - \mu_{ij}\,\mu_{ji}.
  \label{eq:umap-w}
\end{equation}

\paragraph{Schritt 2: Niedrigdimensionale Einbettung.}

Im 2D-Zielraum (\texttt{n\_components=2}) werden Positionen $y_i$
gesucht, die einen \"ahnlichen Nachbarschaftsgraph $v_{ij}$ erzeugen.
Die Wahrscheinlichkeit im Zielraum ist
\begin{equation}
  v_{ij} = \bigl(1 + a\,\|y_i - y_j\|^{2b}\bigr)^{-1},
\end{equation}
mit Konstanten $a \approx 1.577$ und $b \approx 0.895$ aus der
Default-Konfiguration (\texttt{min\_dist=0.1}). Die Verlustfunktion
ist die fuzzy Cross-Entropy:
\begin{equation}
  \mathcal{L} = \sum_{ij} \Bigl[w_{ij}\log\frac{w_{ij}}{v_{ij}}
              + (1 - w_{ij})\log\frac{1 - w_{ij}}{1 - v_{ij}}\Bigr],
  \label{eq:umap-loss}
\end{equation}
minimiert mittels stochastischem Gradientenabstieg.

\paragraph{Bewahrte Eigenschaften.}

Im Vergleich zu t-SNE bewahrt UMAP nicht nur lokale, sondern auch
globale Strukturen \cite{Becht2019}. Distanzen zwischen
Cluster-Regionen in der UMAP-Projektion korrelieren mit Distanzen im
Featureraum -- mit der Einschr\"ankung, dass gro\ss e Distanzen leicht
\emph{komprimiert} werden (UMAP ist nicht abstandserhaltend, aber
topologie-erhaltend).

\subsubsection*{Parameter}

\begin{itemize}
  \item \texttt{n\_neighbors} (Auto-Heuristik:
        $\max(5, \min(15, N/100))$): Steuert die Balance lokal vs.\
        global. Kleine Werte (5--10) bewahren feinste lokale Struktur;
        gro\ss e Werte (20--50) betonen globale Topologie. Bei einem
        typischen [2Fe-2S]-Protein mit $N \approx 5\,000$--$10\,000$ Modes ergibt die
        Auto-Heuristik $\texttt{n\_neighbors} = 15$, was dem
        Paper-Standard \cite{McInnes2018} entspricht.

  \item \texttt{min\_dist} = 0.1 (Default): minimaler Abstand
        zwischen Punkten in der Einbettung. Kleinere Werte erzeugen
        kompaktere Cluster.

  \item \texttt{n\_components} = 2: 2D-Visualisierung.

  \item \texttt{random\_state} = 42 (deterministisch reproduzierbar).
\end{itemize}

Konfigurierbar \"uber \texttt{cfg.umap\_n\_neighbors} (manueller
Override).

\subsubsection*{Feature-Matrix}
\label{sec:embedding-features}

Pro Mode werden 31 numerische Features extrahiert. Zwei urspr\"unglich
geplante \texttt{Cys *\_OOP}-Features sind per geometrischer Konstruktion
immer null (Cys-S$_\gamma$ liegt in der Cluster-Ebene, hat keine
Out-of-plane-Komponente) und wurden ausgeschlossen. Die verbleibenden 31
Features sind:

\begin{enumerate}\setlength{\itemsep}{0pt}\setlength{\parsep}{0pt}
  \item[\textbf{Cluster-Ring-Anteile}] (3 Features, siehe
        Abschnitt~\ref{sec:scsd}):
        \begin{itemize}\setlength{\itemsep}{0pt}
          \item \texttt{lig\_OOP\%} -- Anteil der Mode-Bewegung in der
                Liganden-Sph\"are, der OOP zur Cluster-Ebene erfolgt
          \item \texttt{kern\_OOP\%} -- dasselbe f\"ur den
                Cluster-Kern (Fe + Cluster-S)
          \item \texttt{kern\_|d|} -- absolute Cluster-Auslenkung in
                \AA{} (Norm)
        \end{itemize}

  \item[\textbf{Kern-Scores / D$_{2h}$-Irreps}] (10 Features, siehe
        Abschnitt~\ref{sec:scsd}):
        \begin{itemize}\setlength{\itemsep}{0pt}
          \item \texttt{Translation}, \texttt{Umbrella},
                \texttt{Hinge-Fe}, \texttt{Hinge-S},
                \texttt{OOP-Twist}, \texttt{Fe-Streckung},
                \texttt{S-Streckung}, \texttt{Breathing},
                \texttt{Rhombus-Scherung}, \texttt{Rotation-ip}
        \end{itemize}
        Diese sind Projektionen der Mode auf die irreduziblen
        Darstellungen der idealen $D_{2h}$-Cluster-Symmetrie.

  \item[\textbf{Liganden OOP-Anteile}] (2 Features):
        \begin{itemize}\setlength{\itemsep}{0pt}
          \item \texttt{His 255\_OOP}, \texttt{His 259\_OOP}
        \end{itemize}
        Verteilung der Liganden-OOP-Bewegung auf die zwei His. Die
        Werte erf\"ullen \texttt{His 255\_OOP + His 259\_OOP = 1}; sie
        sind perfekt antikorreliert. UMAP gewichtet jedoch
        antikorrelierte Features nicht st\"arker, und die symmetrische
        Behandlung beider Liganden ist konzeptionell sauberer.
        \emph{Nicht} aufgenommen: \texttt{Cys 207\_OOP},
        \texttt{Cys 216\_OOP} -- siehe oben.

  \item[\textbf{Liganden INP-Anteile}] (4 Features):
        \begin{itemize}\setlength{\itemsep}{0pt}
          \item \texttt{Cys 207\_INP}, \texttt{Cys 216\_INP},
                \texttt{His 255\_INP}, \texttt{His 259\_INP}
        \end{itemize}
        Dasselbe f\"ur In-plane-Anteile. Auch hier perfekte
        Antikorrelation paarweise, alle behalten.

  \item[\textbf{Pro-Liganden-Atom-Bewegungen}] (12 Features, siehe
        Sheets \texttt{Fe\_S\_Cys}, \texttt{Fe\_N\_His}):
        \begin{itemize}\setlength{\itemsep}{0pt}
          \item Pro Cys (207, 216): \texttt{stretch},
                \texttt{bend\_inp}, \texttt{bend\_oop}
          \item Pro His (255, 259): \texttt{stretch},
                \texttt{bend\_inp}, \texttt{bend\_oop}
        \end{itemize}
        Stretch und Bend des Fe-S-/Fe-N-Bindungsvektors mit
        Aufspaltung der Bend-Komponente in INP- und OOP-Anteile
        bez\"uglich der Cluster-Ebene.
\end{enumerate}

Vor UMAP wird die Feature-Matrix standardisiert
(\texttt{sklearn.preprocessing.StandardScaler}; pro Feature Mittelwert
0, Varianz 1). Damit hat keine einzelne Variable wegen ihres
absoluten Wertebereichs Vorrang.

\paragraph{Was bewusst nicht in der Matrix steht.}

\begin{itemize}
  \item \emph{Frequenz} ($\omega_l$): wird nicht als Feature
        verwendet, weil sonst die UMAP-Karte trivial nach Frequenz
        sortiert w\"are. Frequenz steht als Metadaten-Spalte daneben
        und kann zum Einf\"arben des UMAP-Plots verwendet werden.

  \item \emph{Reorganisations-Lambda} ($\lambda_X$): nicht als Feature
        eingef\"uhrt, weil die zugrundeliegenden geometrischen
        Modulationen (\texttt{dr\_FeFe}, \texttt{dr\_NH},
        \texttt{dr\_FeN}) in den Per-Liganden-stretch-/bend-Features
        bereits indirekt enthalten sind. Doppelerfassung w\"urde die
        Reorg-relevanten Modes \"uberbetonen.

  \item \emph{Sekund\"arstruktur} (HELIX/SHEET-Anteile): wird im
        \emph{separaten} Sekund\"arstruktur-UMAP behandelt
        (Sheet \texttt{SSE\_UMAP\_Cluster}). Mode-UMAP und
        SSE-UMAP sind konzeptionell unterschiedliche Auswertungen.
\end{itemize}

\subsubsection*{HDBSCAN: dichtebasierte Cluster-Findung}

HDBSCAN \cite{Campello2013,McInnes2017} ist ein hierarchischer
dichtebasierter Cluster-Algorithmus. Im Unterschied zu k-Means braucht
er keine vorgegebene Cluster-Anzahl; er identifiziert Cluster als
zusammenh\"angende Regionen hoher Dichte und markiert isolierte
Punkte als Rauschen (Label $-1$).

\paragraph{min\_cluster\_size.} Der einzige wichtige Parameter:
minimale Cluster-Gr\"o\ss e. \texttt{modenanalyse\_2fe2s} verwendet eine
Auto-Heuristik
\[
  \texttt{mcs} = \max\bigl(5, \min(30, \sqrt{N})\bigr),
\]
F\"ur typische [2Fe-2S]-Protein-Systeme ($N \approx 5\,000$--$10\,000$) ergibt das
$\texttt{mcs} = 30$.

\begin{infobox}
\textbf{Wenn HDBSCAN viele Cluster oder viel Rauschen findet --
das hei\ss t was?}

UMAP auf einem typischen Cys$_2$His$_2$-System mit 31 Features und
\texttt{n\_neighbors=15} liefert typisch 5--8 Cluster ohne Rauschen.
Wenn HDBSCAN signifikant mehr Cluster findet (z.\,B.\ 80--100), liegt
das meist an einer zu kleinen \texttt{mcs}-Wahl oder an einer
fragmentierten UMAP-Topologie. Wenn HDBSCAN 0 Cluster findet
(alles Rauschen), ist die Mode-Verteilung im Featureraum zu glatt --
das ist erwartetes Verhalten f\"ur DFT-Mode-Daten ohne harte
\"Uberg\"ange zwischen Mode-Klassen \cite{Campello2013}.

\textbf{Konsequenz.} Die UMAP-Visualisierung selbst ist auch ohne
HDBSCAN-Cluster aussagekr\"aftig -- man sieht Mode-Dichten und
\"Uberg\"ange durch Einf\"arbung nach Frequenz, Mode-Typ oder
Kern-Score. HDBSCAN ist zus\"atzliche Information, kein notwendiger
Bestandteil der Auswertung.
\end{infobox}

\paragraph{Erforderlich.} \texttt{hdbscan} und \texttt{umap-learn}
m\"ussen installiert sein. Ohne diese Pakete werden die UMAP-Sheets
\"ubersprungen.

\subsubsection*{Praktische Empfehlung}

\begin{enumerate}
  \item \textbf{UMAP-Plot anschauen}: das Sheet
        \texttt{Projektionen} enth\"alt die UMAP-Koordinaten f\"ur
        alle Modes. Mit Origin/matplotlib einf\"arben nach Frequenz,
        Mode-Typ, oder Kern-Primary-Score.

  \item \textbf{Cluster-Sheets pr\"ufen}: das Sheet
        \texttt{UMAP\_Cluster} zeigt HDBSCAN-Labels und
        Feature-Werte pro Cluster. \texttt{UMAP\_Cluster\_Profil}
        zeigt Z-Score-Profile und die Top-Modes pro Cluster.

  \item \textbf{Vergleich zwischen Strukturen} (z.\,B.\ prot vs deprot,
        vs Mut-Variante) erfolgt extern: zwei UMAP-Plots f\"ur die
        zwei Strukturen nebeneinander, mit gleicher Frequenz-Skala
        einf\"arben. Innerhalb einer Auswertung wird kein direkter
        Vergleich durchgef\"uhrt -- ein System pro Lauf.
\end{enumerate}


% =========================================================================
% Theorie-Kapitel zur mode-aufgel\"osten Reorg-Pipeline
% =========================================================================


\subsection{Marcus-Hush-Theorie der Reorganisation}
\label{sec:marcus-hush}

Das Reorg-Kapitel \S\ref{sec:reorg} berechnet pro Mode Beitr\"age zur
Reorganisations-Energie. Bevor wir dort hineingehen, brauchen wir die
klassische Marcus-Hush-Theorie als \emph{Vergleichs-Anker} -- damit klar
ist, was unsere mode-aufgel\"oste Variante mit der historischen
Marcus-Theorie gemeinsam hat und wo sie sich unterscheidet.

\subsubsection*{Klassisches Bild: zwei diabatische Energiefl\"achen}

Marcus' Theorie des Elektronentransfers \cite{Marcus1956,Marcus1993}
beschreibt eine Reaktion zwischen einem Reaktant- und einem
Produkt-Zustand durch zwei \emph{diabatische} Energiefl\"achen
$U_R(\mathbf{Q})$ (Reaktant) und $U_P(\mathbf{Q})$ (Produkt). Die
Reaktionskoordinaten $\mathbf{Q} = (Q_1, Q_2, \ldots, Q_N)$ sind die
Schwingungsmoden des reaktiven Komplexes -- \emph{alle}, nicht
einzelne.

In der harmonischen Approximation sind beide Fl\"achen Paraboloide:
\begin{equation}
  U_R(\mathbf{Q}) = U_R^0 + \frac{1}{2}\sum_i \mu_i\,\omega_i^2\,
                       (Q_i - Q_i^R)^2
\end{equation}
\begin{equation}
  U_P(\mathbf{Q}) = U_P^0 + \frac{1}{2}\sum_i \mu_i\,\omega_i^2\,
                       (Q_i - Q_i^P)^2
\end{equation}
wobei $\mu_i$ die reduzierte Masse, $\omega_i$ die Eigenfrequenz und
$Q_i^R, Q_i^P$ die Gleichgewichtsverschiebungen der Mode $i$ in den
beiden Zust\"anden sind.

Die zentrale Annahme von Marcus: \textbf{die Frequenzen $\omega_i$
sind in beiden Zust\"anden gleich}, nur die Verschiebungen
$Q_i^R \neq Q_i^P$. Das ist die "`harmonisch-gleich"'-Approximation;
sie gilt, wenn Reaktant und Produkt sich elektronisch \"ahneln (z.\,B.\
zwei Oxidationsstufen desselben Komplexes).

\subsubsection*{Reorganisations-Energie als zentrale Gr\"o{\ss}e}

Die \textbf{Reorganisations-Energie} $\lambda$ ist definiert als die
Energie, die das System zahlt, wenn es im Reaktant-Zustand bleibt,
aber zur Produkt-Geometrie verschoben wird:
\begin{equation}
  \lambda = U_R(\mathbf{Q}^P) - U_R(\mathbf{Q}^R)
          = \frac{1}{2}\sum_i \mu_i\,\omega_i^2\,(\Delta Q_i)^2
  \quad\text{mit}\quad \Delta Q_i = Q_i^P - Q_i^R.
  \label{eq:lambda-marcus-classical}
\end{equation}
Dies ist die ber\"uhmte Marcus-Hush-Standardform. Sie ist eine
\emph{Gleichgewichts-Eigenschaft} des Systems -- keine
Mode-Eigenschaft, keine Aktivierungs-Energie, keine kinetische
Gr\"o{\ss}e. Die thermische Population einzelner Modes spielt in
\eqref{eq:lambda-marcus-classical} keine Rolle.

\subsubsection*{Aktivierungs-Energie und Rate}

Aus $\lambda$ und der freien Reaktionsenergie $\Delta G^0 = U_P^0 -
U_R^0$ folgt nach kurzer Geometrie die Marcus-Aktivierungs-Energie:
\begin{equation}
  \Delta G^\ddagger = \frac{(\lambda + \Delta G^0)^2}{4\lambda}
\end{equation}
und damit die Marcus-Rate:
\begin{equation}
  k_\text{ET} = \frac{2\pi}{\hbar}|H_{RP}|^2 \cdot
                \frac{1}{\sqrt{4\pi\lambda k_BT}} \cdot
                \exp\!\left(-\frac{\Delta G^\ddagger}{k_BT}\right)
  \label{eq:marcus-rate}
\end{equation}
mit $H_{RP}$ dem elektronischen Kopplungselement zwischen Reaktant-
und Produkt-Diabaten. Drei Charakteristika:

\begin{itemize}\setlength{\itemsep}{0pt}
  \item Bei kleinem $|\Delta G^0|$: $\Delta G^\ddagger \approx
        \lambda/4$. \emph{Reorganisations-Energie ist die Ma{\ss}zahl
        f\"ur die Reaktionsbarriere.}
  \item Bei $\Delta G^0 = -\lambda$: barrierfrei, $k$ maximal
        ($\Delta G^\ddagger = 0$).
  \item Bei $|\Delta G^0| > \lambda$: \emph{Marcus-inverter Bereich},
        Rate sinkt wieder, weil die Reaktion zu exotherm wird.
\end{itemize}

In Eisen-Schwefel-Proteinen sind alle drei Bereiche relevant:
ferredoxinartige ET-Schritte sind oft im Marcus-Normalbereich;
PCET-Schritte in mitoNEET und Rieske k\"onnen die invertierte Region
ber\"uhren.

\subsubsection*{Was diese Theorie liefert -- und was sie nicht liefert}

\paragraph{Was Marcus-Hush leistet:}
\begin{itemize}\setlength{\itemsep}{0pt}
  \item Eine vollst\"andige analytische Beschreibung der
        ET/PCET-Aktivierungsenergie, falls die harmonische
        Approximation g\"ultig ist.
  \item Eine direkt messbare Vorhersage der Rate $k_\text{ET}$
        (Glg.~\ref{eq:marcus-rate}), wenn $\lambda$, $\Delta G^0$
        und $|H_{RP}|^2$ bekannt sind.
  \item Eine Begr\"undung, warum $\lambda$ system-spezifisch ist
        (kein Universalwert, h\"angt von der Solvatationsumgebung,
        von der Cluster-Geometrie und von den Liganden ab).
\end{itemize}

\paragraph{Was Marcus-Hush nicht direkt liefert:}
\begin{itemize}\setlength{\itemsep}{0pt}
  \item \emph{Welche} Modes wesentlich zu $\lambda$ beitragen.
        Die Summe in \eqref{eq:lambda-marcus-classical} ist global;
        Promoter-Modes (s.\,u.) mussten in den 1990ern erst durch
        Hammes-Schiffer und Andere herausdestilliert werden.
  \item \emph{Direkten Zugang aus DFT-Frequenzrechnungen}. Klassische
        Berechnung von $\lambda$ erfordert zwei separate
        Frequenzrechnungen (Reaktant + Produkt) plus ein robustes
        Mode-Mapping zwischen beiden -- ein nicht-trivialer
        Aufwand.
  \item \emph{Reorganisations-Energie als spektroskopische Observable}.
        $\lambda$ aus \eqref{eq:lambda-marcus-classical} ist eine
        Aktivierungs-Energie f\"ur eine konkrete Reaktion, keine
        Bandenfrequenz, kein Linienprofil.
\end{itemize}

Diese drei Lücken motivieren die Erweiterung in zwei Richtungen:

\begin{enumerate}
  \item \textbf{Mode-aufgel\"oste Sicht} (Hammes-Schiffer, Cukier,
        Stuchebrukhov): zerlegt $\lambda$ in Beitr\"age einzelner
        Modes und identifiziert Promoter-Modes mit besonderer
        Bedeutung f\"ur die Rate. Die Multi-Mode-Sicht macht
        \eqref{eq:lambda-marcus-classical} zu einer
        Frequenz-aufgel\"osten Information.
  \item \textbf{Spektroskopisch zug\"angliche Variante}: ersetzt die
        statischen Verschiebungen $\Delta Q_i$ durch thermische
        Bindungs-Modulationen $\Delta r_X(i)$. Damit wird aus
        $\lambda$ eine Gr\"o{\ss}e, die direkt aus einer einzigen
        Frequenzrechnung folgt -- und die Banden in NRVS-Spektren
        zuordnet.
\end{enumerate}

Die zweite Erweiterung ist genau das, was \texttt{modenanalyse\_2fe2s}
in \S\ref{sec:reorg} berechnet. Was wir dort $\lambda_X$ nennen, ist
\emph{nicht} das klassische Marcus-Hush-$\lambda$ aus
\eqref{eq:lambda-marcus-classical}, sondern eine verwandte aber
unterscheidbare Gr\"o{\ss}e -- die thermische Bindungs-Reorg-Energie pro
Mode und Bindung. Der Zusammenhang zur Huang-Rhys-Theorie macht klar,
warum diese Gr\"o{\ss}e eine wohldefinierte spektroskopische Observable
ist (n\"achster Abschnitt).

\subsubsection*{Hammes-Schiffer-Erweiterung: Multi-Mode-PCET}

Hammes-Schiffer und Soudackov \cite{HammesSchifferSoudackov2008}
zeigen, dass PCET in komplexen Systemen wie Proteinen \emph{nicht}
durch eine einzelne Promoter-Mode beschrieben werden kann. Stattdessen
ist die Reaktion an viele Modes gleichzeitig gekoppelt; das Spektrum
der Beitr\"age $\lambda_i = \tfrac{1}{2}\mu_i\omega_i^2(\Delta Q_i)^2$
ist breit verteilt.

Die naive Standard-Marcus-Theorie kann diese Multi-Mode-Realit\"at
durch eine \emph{effektive} Reorganisations-Energie
$\lambda_\text{eff} = \sum_i \lambda_i$ darstellen, aber die
mode-aufgel\"oste Verteilung geht dabei verloren. F\"ur die
NRVS-Bandenidentifikation ist genau diese Verteilung wichtig:

\begin{itemize}\setlength{\itemsep}{0pt}
  \item Welche Frequenzbereiche tragen am st\"arksten zur Reorg bei?
  \item Welche Bindungen sind in diesen Bereichen aktiv?
  \item Welche Mutationen sollten welche Bandenstruktur ver\"andern?
\end{itemize}

Die Multi-Mode-Sicht ist auch der Grund, warum globale
"`Top-Mode-Klassifikations"'-Ans\"atze (PCET-Score, $L_\text{PCET}$,
filter\_variant) konzeptionell unbefriedigend waren: sie suchten nach
\emph{einer} dominierenden Mode, w\"ahrend die Realit\"at eine breite
Verteilung ist. Die mode-aufgel\"oste Reorg-Pipeline macht
diese Verteilung explizit.

\subsubsection*{Was bedeutet $\Delta Q$ vs.\ $\Delta r_X$?}

Eine zentrale Subtilit\"at zur Vorbereitung des n\"achsten Kapitels: die
klassische Marcus-Verschiebung $\Delta Q_i$ (massengewichtete
Mode-Koordinate) ist \emph{nicht} dasselbe wie eine
Bindungs-Modulation $\Delta r_X$ (kartesische Distanz-\"Anderung in
$\angstrom$).

\begin{itemize}\setlength{\itemsep}{0pt}
  \item $\Delta Q$: Verschiebung der \emph{Eigenmode-Koordinate} zwischen
        zwei Gleichgewichtsgeometrien. Hat Einheit
        $\sqrt{\text{amu}}\cdot\angstrom$. Tritt in
        $\eqref{eq:lambda-marcus-classical}$ auf.
  \item $\Delta r_X(i)$: \emph{Bindungs-Modulation} entlang der
        $X$-Bindung, induziert durch Mode $i$ mit thermisch
        skalierten Auslenkungen. Hat Einheit $\angstrom$. Wird in
        \texttt{reorganization.compute\_mode\_modulations} berechnet.
\end{itemize}

Beide haben dieselbe \emph{Form} in der Standard-Glg.\
$\lambda = \tfrac{1}{2}\mu\omega^2(\cdot)^2$, aber unterschiedliche
\emph{physikalische Bedeutung}. Die Br\"ucke zwischen beiden ist die
Huang-Rhys-Theorie -- siehe n\"achster Abschnitt.

\subsection{Huang-Rhys-Theorie und unser $\lambda_X$}
\label{sec:huang-rhys}

In \S\ref{sec:reorg} berechnen wir
$\lambda_X(i) = \tfrac{1}{2}\mu_i\omega_i^2(\Delta r_X(i))^2$ mit
\emph{thermisch skalierten} Bindungs-Modulationen. Dieser Abschnitt
zeigt, dass das exakt der \textbf{Huang-Rhys-Faktor} ist (in einer
bestimmten Konvention), und dass diese Gr\"o{\ss}e eine wohldefinierte
spektroskopische Bedeutung hat.

\subsubsection*{Was Huang-Rhys urspr\"unglich beschrieben hat}

Huang und Rhys \cite{HuangRhys1950} berechneten optische Linienprofile
von Farbzentren in Festk\"orpern. Ihre Frage: wenn ein elektronischer
\"Ubergang die Gleichgewichtsposition einer harmonischen Schwingung
verschiebt, wie sieht dann das resultierende Vibronen-Spektrum aus?

Antwort: das Spektrum ist eine Poisson-Verteilung
\begin{equation}
  P(n) = \frac{S^n e^{-S}}{n!}
\end{equation}
\"uber Phononen-Anzahl $n$, parametrisiert durch den
\textbf{Huang-Rhys-Faktor}
\begin{equation}
  S = \frac{\mu\,\omega\,(\Delta q)^2}{2\hbar},
  \label{eq:huang-rhys}
\end{equation}
wobei $\Delta q$ die \"ubergangs-induzierte Gleichgewichts-Verschiebung
einer Mode mit reduzierter Masse $\mu$ und Frequenz $\omega$ ist.

\paragraph{Bedeutung von $S$:}
\begin{itemize}\setlength{\itemsep}{0pt}
  \item $S \ll 1$: \"Ubergang ist ann\"ahernd vertikal,
        Linienspektrum dominiert durch 0-0-Bande.
  \item $S \approx 1$: starke Mode-Kopplung, Vibronen-Bande mit
        mehreren Phononen-Anregungen.
  \item $S \gg 1$: Gauss-f\"ormige Bandenstruktur (Stokes-Verschiebung
        gleich $S\hbar\omega$).
\end{itemize}

In moderner Sprache: $S$ misst die Mode-Kopplung \emph{quantitativ},
in Einheiten der Nullpunkts-Auslenkung
$u_\text{ZP} = \sqrt{\hbar/(2m\omega)}$:
\begin{equation}
  S = \frac{(\Delta q)^2}{2 u_\text{ZP}^2}.
  \label{eq:huang-rhys-units}
\end{equation}

\subsubsection*{Verbindung zu Marcus-Hush}

Vergleiche \eqref{eq:huang-rhys} mit der Marcus-Reorganisations-Energie
einer einzelnen Mode (aus \eqref{eq:lambda-marcus-classical} mit
$N=1$):
\begin{equation}
  \lambda^\text{Marcus}_\text{single} = \frac{1}{2}\mu\,\omega^2\,(\Delta q)^2
  = S \cdot \hbar\omega.
  \label{eq:lambda-vs-S}
\end{equation}

Damit gilt: \textbf{$\lambda^\text{Marcus} = S \cdot \hbar\omega$}.
Die beiden Gr\"o{\ss}en sind ineinander \"uberf\"uhrbar.

Bei einer Multi-Mode-Verteilung:
$\lambda^\text{Marcus}_\text{total} = \sum_i S_i \cdot \hbar\omega_i$.

\subsubsection*{Was ist $\Delta q$ in unserem Fall?}

Der Knackpunkt: in der klassischen Marcus-Theorie ist $\Delta q$ eine
\emph{statische Verschiebung} der Mode-Koordinate zwischen zwei
elektronischen Zust\"anden. \texttt{modenanalyse\_2fe2s} hat aber
\emph{nur einen} elektronischen Zustand pro Lauf -- die einzelne
DFT-Frequenzrechnung gibt nur eine Gleichgewichtsgeometrie und ihre
Eigenmodes.

L\"osung: wir ersetzen die statische Verschiebung durch eine
\emph{thermische Auslenkung}:
\begin{equation}
  (\Delta q)^2 \to \langle q^2 \rangle = \frac{\hbar}{2\mu\omega}
  \coth\!\left(\frac{\hbar\omega}{2k_BT}\right)
  = u_\text{rms}^2.
  \label{eq:thermal-replacement}
\end{equation}

Mit der Konvention $\Delta r_X(i) = u_\text{rms}(i) \cdot \alpha_X(i)$,
wobei $\alpha_X(i)$ der Mode-Anteil auf der Bindung $X$ ist, wird
unser $\lambda_X(i)$:
\begin{equation}
  \lambda_X(i) = \frac{1}{2}\mu_i\,\omega_i^2\,
                  u_\text{rms}^2(i)\,\alpha_X^2(i).
\end{equation}

Setzt man \eqref{eq:thermal-replacement} ein und nutzt
$\hbar\omega/(k_BT) \to \infty$ bei $T \to 0$ (also $\coth \to 1$),
erh\"alt man:
\begin{equation}
  \lambda_X(i) \xrightarrow{T \to 0}
  \frac{1}{2}\mu_i\,\omega_i^2\,\frac{\hbar}{2\mu_i\omega_i}\,\alpha_X^2(i)
  = \frac{\hbar\omega_i}{4}\,\alpha_X^2(i).
  \label{eq:lambda-zero-T}
\end{equation}

\paragraph{Verbindung zum Huang-Rhys-Faktor.}
Aus \eqref{eq:lambda-vs-S} folgt mit $(\Delta q)^2 = u_\text{rms}^2 \cdot \alpha_X^2$:
\begin{equation}
  S_X(i) = \frac{\mu_i\omega_i u_\text{rms}^2(i)\,\alpha_X^2(i)}{2\hbar}
  \xrightarrow{T \to 0} \frac{\alpha_X^2(i)}{4}.
\end{equation}
Im T → 0-Limit ist also der Huang-Rhys-Faktor pro Bindung
einfach \emph{ein Viertel des quadrierten Mode-Anteils}. F\"ur eine
reine Streckmode der Bindung ($\alpha_X = 1$) gilt
$S_X = 1/4$ -- also \emph{schwache Kopplung} im Huang-Rhys-Sinn,
typisches Linienspektrum.

\subsubsection*{Bedeutung der Gr\"o{\ss}enordnung}

Setzen wir konkrete Zahlen ein. F\"ur eine typische
Fe-S-Streckschwingung bei $\omega = 350$~cm$^{-1}$ und $T = 5$\,K:
\begin{itemize}\setlength{\itemsep}{0pt}
  \item Nullpunkts-Energie: $\hbar\omega/4 = 87.5$~cm$^{-1}$.
  \item F\"ur eine reine Streckmode ($\alpha_X = 1$):
        $\lambda_X(T \to 0) = 87.5$~cm$^{-1}$.
  \item Bei thermischer Population (5\,K $\ll \hbar\omega/k_B = 503$\,K):
        $\coth \approx 1$, also nahezu Nullpunkts-Wert.
  \item Bei Raumtemperatur (300\,K, $\coth \approx 1.04$):
        $\lambda_X \approx 91$~cm$^{-1}$, kaum Erh\"ohung.
\end{itemize}

Dieselbe Mode bei $\omega = 100$~cm$^{-1}$:
\begin{itemize}\setlength{\itemsep}{0pt}
  \item $\hbar\omega/4 = 25$~cm$^{-1}$.
  \item Bei 5\,K: $\coth \approx 1.0$, $\lambda_X \approx 25$~cm$^{-1}$.
  \item Bei 300\,K: $\coth \approx 4.2$, $\lambda_X \approx 105$~cm$^{-1}$.
\end{itemize}
Niederfrequente Modes sind also bei Raumtemperatur stark thermisch
populiert -- ihre $\lambda_X$-Beitr\"age werden gegen\"uber dem
Nullpunkts-Wert um Faktor $\coth(\hbar\omega/2k_BT)$ verst\"arkt.

\paragraph{Konsequenz f\"ur Tieftemperatur-NRVS.}
Bei Tieftemperatur (5--40\,K, $\hbar\omega \gg k_BT$ f\"ur typische
NRVS-Frequenzen) ist $u_\text{rms}^2 \approx \hbar/(2\mu\omega)$, also
$T$-unabh\"angig. Damit gilt
\begin{equation}
  \lambda_X(i) \approx \frac{\hbar\omega_i}{4}\alpha_X^2(i)
  \quad\text{f\"ur Tieftemperatur-NRVS}.
\end{equation}
Diese Form ist \emph{rein quantenmechanisch} (Nullpunkts-Energie),
unabh\"angig von der genauen Probentemperatur. Das ist genau das, was
NRVS misst: die quantenmechanische Mode-Aktivit\"at, nicht eine
thermische Boltzmann-Mittelung.

\subsubsection*{Was unser $\lambda_X$ also ist -- abschlie{\ss}ende Definition}

Mit allen Vorarbeiten k\"onnen wir jetzt eine pr\"azise Definition
formulieren:

\begin{tcolorbox}[colback=accent!5,colframe=accent,
  title={Was $\lambda_X$ in \texttt{modenanalyse\_2fe2s} ist}]
$\lambda_X(i)$ ist die \textbf{thermische Bindungs-Reorg-Energie}
der Mode $i$ entlang der Bindungs-Klasse $X$:
\begin{itemize}\setlength{\itemsep}{0pt}
  \item Form: identisch zu Marcus-Hush ($\tfrac{1}{2}\mu\omega^2 r^2$),
        aber mit thermisch skaliertem $r$ statt statischem
        Reaktant-Produkt-Versatz.
  \item Verwandt mit dem Huang-Rhys-Faktor: $S_X(i) =
        \lambda_X(i)/(\hbar\omega_i)$.
  \item Im Tieftemperatur-Limit:
        $\lambda_X(i) \to (\hbar\omega_i/4)\,\alpha_X^2(i)$
        -- rein quantenmechanisch, $T$-unabh\"angig.
  \item Direkte spektroskopische Bedeutung: $\lambda_X(i)$ ist die
        Energie der Bindung-$X$-Modulation in Mode $i$, gemessen am
        QHO-Erwartungswert.
\end{itemize}

\textbf{Diese Gr\"o{\ss}e ist verwandt mit, aber nicht identisch zu, der
klassischen Marcus-Hush-Aktivierungs-Energie f\"ur eine Reaktion
zwischen zwei elektronischen Zust\"anden.} F\"ur letztere br\"auchte man
zwei separate DFT-Rechnungen.
\end{tcolorbox}

\subsection{NRVS-Theorie soweit relevant f\"ur Bandenidentifikation}
\label{sec:nrvs}

Das Tool berechnet \emph{nicht} NRVS-Spektren -- diese werden mit
spezialisierten externen Werkzeugen aus derselben DFT-Frequenzrechnung
synthetisiert. Aber die in
\S\ref{sec:reorg} berechneten Modulations-Spektren $M_X(\omega)$ und
kumulativen Reorg-Kurven $\Lambda_X(\omega)$ sind direkt mit
NRVS-Banden vergleichbar -- darum brauchen wir die NRVS-Theorie als
Interpretations-Rahmen.

\subsubsection*{Was misst NRVS?}

NRVS \cite{Sturhahn2004} misst die
\emph{$^{57}$Fe-projizierte Vibrations-Zustandsdichte}:
\begin{equation}
  D_\text{Fe}(\omega) = \sum_i |\mathbf{e}_\text{Fe}^{(i)}|^2 \cdot
                          \delta(\omega - \omega_i)
  \label{eq:fe-pdos}
\end{equation}
gewichtet mit dem Lamb-M\"o{\ss}bauer-Faktor $f_\text{LM}$ und einer
multi-Phonon-Konvolution. F\"ur die Bandenidentifikation reicht meist
die Fe-pDOS ohne Multi-Phonon -- das nennen wir $g_\text{Fe}(\omega)$.

Wichtige Eigenschaften:
\begin{itemize}\setlength{\itemsep}{0pt}
  \item NRVS ist \textbf{element-selektiv}: nur Modes mit
        Fe-Beteiligung erscheinen. Modes mit reinen Liganden-Bewegungen
        (z.\,B.\ C-H-Stretching) sind nicht sichtbar.
  \item NRVS ist \textbf{nicht symmetrieselektiv}: anders als
        Raman/IR sind alle Modes sichtbar, unabh\"angig von ihrer
        D$_{2h}$-Symmetrie. Das macht NRVS besonders n\"utzlich f\"ur
        Eisen-Schwefel-Cluster, wo viele Modes per
        Auswahlregeln in IR/Raman dunkel sind.
  \item NRVS ist \textbf{quantitativ}: das Integral der pDOS ist auf
        $1$ normiert (oder $3 N_\text{Fe}$ in einigen
        Konventionen). Bandenintensit\"aten sind direkt
        Mode-Kopplungen.
\end{itemize}

\subsubsection*{Bandenstruktur in [2Fe-2S]-Proteinen}

Die typische [2Fe-2S]-NRVS-Bandenstruktur (Gee 2021 \cite{Gee2021}; 
eine umfassende \"Ubersicht der NRVS-Modi geben Wang \emph{et al.}\ 
2025 \cite{Wang2025}) hat drei Hauptbereiche:

\paragraph{Niederfrequenz-Bereich (50--200\,cm$^{-1}$).}
Skelett-Schwingungen des gesamten Proteins. F\"ur das Cluster relevant:
Cluster-Ebenen-Bewegungen (translatorisch + rotatorisch in der
Protein-Matrix), Cluster-OOP-Buckel (akustische Mode-Mischung mit
Liganden). $\lambda_X$ ist hier in allen Kan\"alen klein
($\lambda_X \approx \hbar\omega/4 \approx 12$--50\,cm$^{-1}$ pro
Mode), aber die \emph{Anzahl} der Modes ist gro{\ss}.

\paragraph{Mittlerer Bereich (200--400\,cm$^{-1}$).}
Hier liegen die meisten Cluster-Bindungs-Streckschwingungen:
\begin{itemize}\setlength{\itemsep}{0pt}
  \item Fe-N(His)-Stretching:
        260--290\,cm$^{-1}$ in mitoNEET-typischer Cys$_3$His$_1$-
        Koordination (Gee 2021 \cite{Gee2021}: berechnet 223\,cm$^{-1}$
        f\"ur Fe(II)-N$_\delta$(His) im reduzierten Zustand, ca.\
        270--290\,cm$^{-1}$ im oxidierten Zustand). F\"ur klassische
        Rieske-Cys$_2$His$_2$-Cluster: Doppelbande bei 266 + 274\,cm$^{-1}$
        (Kuila \emph{et al.}\ 1992 \cite{Kuila1992}, Fe(III)-N), zugeordnet
        zu protoniertem (266) bzw.\ deprotoniertem (274) Imidazol.
  \item Fe-Fe-Cluster-Atmung: $\sim 200$\,cm$^{-1}$, je nach Cluster-
        Geometrie und Liganden.
  \item Fe-S$_\text{bridge}$-symmetrische Streckung: 300--340\,cm$^{-1}$.
  \item Fe-S$_\text{Cys}$-Stretching: 280--320\,cm$^{-1}$.
\end{itemize}
F\"ur eine reine Fe-X-Streckmode mit $\alpha_X \approx 1$ ist
$\lambda_X \approx \hbar\omega/4 \approx 60$--$100$\,cm$^{-1}$ -- gro{\ss}er
Beitrag pro Mode.

\paragraph{Hochfrequenz-Bereich (400--800\,cm$^{-1}$).}
\begin{itemize}\setlength{\itemsep}{0pt}
  \item Fe-S$_\text{bridge}$-asymmetrische Streckung: 380--420\,cm$^{-1}$.
  \item Imidazol-Ring-Modes mit Fe-Beteiligung: 600--700\,cm$^{-1}$.
  \item Cluster-Liganden-Biegungen mit Fe-Anteil: 500--700\,cm$^{-1}$.
\end{itemize}

\subsubsection*{Wie man $\Lambda_X(\omega)$ mit NRVS vergleicht}

Die kumulative Reorg-Kurve $\Lambda_X(\omega) = \sum_{\omega_i \leq
\omega} \lambda_X(i)$ ist eine \textbf{Stufenfunktion}: jede Mode mit
nicht-trivialem $\lambda_X$-Beitrag erzeugt einen Sprung der H\"ohe
$\lambda_X(i)$ bei $\omega = \omega_i$.

Im Vergleich mit dem NRVS-Spektrum:
\begin{itemize}\setlength{\itemsep}{0pt}
  \item Wo eine Bande in der Fe-pDOS liegt, sollte ein Sprung in
        $\Lambda_X(\omega)$ sein -- f\"ur den Kanal $X$, der diese
        Bande dominiert.
  \item Die H\"ohe des Sprungs ist proportional zu $\hbar\omega \cdot
        \alpha_X^2$. Ein Sprung von 50\,cm$^{-1}$ bei
        $\omega = 350$\,cm$^{-1}$ entspricht $\alpha_X^2 = 4 \cdot 50/350
        = 0.57$ -- also einer Mode mit etwa 57\% Bindungs-Anteil
        auf $X$.
  \item Wenn das NRVS-Spektrum eine Bande zeigt, aber alle
        $\Lambda_X(\omega)$-Kurven flach bleiben, dann ist die Mode
        in keinem unserer 5 Kan\"ale stark aktiv -- sie ist
        wahrscheinlich eine Liganden-interne Mode mit nur leichter
        Fe-Beteiligung.
\end{itemize}

Diese Korrespondenz macht $\Lambda_X(\omega)$ zu einem direkten
\textbf{Bandenzuordnungs-Werkzeug}. Wenn ein Reviewer fragt:
"`Welche Bindung ist in Bande X aktiv?"', plottet man
$\Lambda_\text{FeFe}, \Lambda_\text{FeN}, \Lambda_\text{FeS},
\Lambda_\text{NH}, \Lambda_\text{HA}$ ueber dem NRVS-Spektrum und
liest die Antwort direkt ab.

\subsubsection*{Tieftemperatur-Konsistenz}

NRVS-Messungen werden typischerweise bei 5--40\,K durchgef\"uhrt
(fluess.\ He), um die Lamb-M\"o{\ss}bauer-Auflöesung zu maximieren und
die thermische Linienverbreiterung zu minimieren. In dieser
Tieftemperatur-Region ist unser $\lambda_X$ identisch mit dem
Nullpunkts-Wert, $\hbar\omega/4 \cdot \alpha_X^2$ -- d.\,h.\ unser
berechnetes $\lambda_X$ \emph{ist} das, was NRVS misst, modulo
Spektral-Verbreiterung und multi-Phonon-Effekten.

Konsequenz: f\"ur die Bandenidentifikation muss man NICHT bei der
exakten Mess-Temperatur rechnen. Ein Lauf bei
\texttt{temp\_k = 5} liefert die Tieftemperatur-Werte; ein Lauf bei
\texttt{temp\_k = 300} liefert die Raumtemperatur-Werte; in der
NRVS-relevanten Frequenzregion ($\omega > 200$\,cm$^{-1}$) sind
beide praktisch identisch, weil die Bose-Einstein-Population dort
verschwindet.

\subsubsection*{Was NRVS \emph{nicht} sieht}

\begin{itemize}\setlength{\itemsep}{0pt}
  \item \textbf{Reine N-H-Streckschwingungen} (Imidazol-NH bei 3000+\,cm$^{-1}$):
        au{\ss}erhalb des typischen NRVS-Messbereichs (Apparate enden
        bei 800--1000\,cm$^{-1}$), und die Mode hat ohnehin nur
        winzigen Fe-Anteil. Konsequenz: $\Lambda_\text{NH}$ ist in
        unseren typischen Filtern (\texttt{freq\_max = 800}) immer
        sehr klein -- siehe Befund-Tabelle.
  \item \textbf{Reine OH-Streckschwingungen}: dito.
  \item \textbf{Backbone-Modes} ohne Fe-Beteiligung: bleiben in der
        pDOS dunkel.
\end{itemize}

Diese Punkte sind wichtig f\"ur die Bandenzuordnungs-Begr\"undung im
Befund: wenn eine Mode in der pDOS dunkel ist, aber unser
$\Lambda_X(\omega)$ dort einen Sprung zeigt, ist das ein \emph{Hinweis}
darauf, dass die Mode existiert und Fe-Beteiligung hat, aber das
Lamb-M\"o{\ss}bauer-Faktor sehr klein ist (Mode ist nicht koh\"arent
mit Fe). Das ist gerade in PCET-Studien interessant.

% =========================================================================
% ENDE NEUE THEORIE-KAPITEL
% =========================================================================

\subsection{Mode-aufgel\"oste Bindungs-Reorganisations-Energien}
\label{sec:reorg}

Das Paket nutzt einen \emph{konzeptionellen
Pivot} vollzogen: weg von einer Klassifikation einzelner Modes als
"`PCET-Kandidat"' (mit Schwellen, Filtern, $\tanh$-S\"attigung), hin
zur direkten Berechnung mode-aufgel\"oster Bindungs-Reorganisations-
Beitr\"age, mit Marcus-Hush-Form aber thermisch skalierten
Mode-Amplituden.

Bevor wir die Pipeline beschreiben: eine Klarstellung der
physikalischen Bedeutung der berechneten Gr\"o{\ss}en.

\subsubsection*{Was $\lambda_X$ ist und was es nicht ist}

\textbf{Klassische Marcus-Hush-Reorganisationsenergie}
\cite{Marcus1956,Marcus1993} ist
\begin{equation}
  \lambda^{\text{Marcus}}_\text{total}
  = \frac{1}{2}\sum_i \mu_i\,\omega_i^2\,(\Delta Q_i)^2,
  \label{eq:marcus-hush}
\end{equation}
wobei $\Delta Q_i$ die \emph{statische} Verschiebung der harmonischen
Mode-Koordinate $i$ zwischen Reaktant- und Produktzustand ist
(Aktivierungs-Energie f\"ur die Reaktion). Die Berechnung dieser Gr\"o{\ss}e
erfordert zwei separate Frequenzrechnungen (Reaktant und Produkt) plus
ein Mode-Mapping zwischen beiden -- Aufwand, den wir hier
\emph{nicht} betreiben.

\textbf{Was wir stattdessen berechnen}, ist eine verwandte aber
verschiedene Gr\"o{\ss}e: die mode-aufgel\"oste Beitragsenergie der
\emph{thermischen Bindungs-L\"angen-Fluktuationen} pro Bindungs-Kanal.
Aus den Eigenvektoren der Frequenzrechnung wird f\"ur jede Mode $i$
und jede Bindung $X$ die thermisch skalierte Bindungs-Modulation
$\Delta r_X(i)$ berechnet (Auslenkungs-Vektoren multipliziert mit
$u_\text{rms}$ aus dem QHO-Erwartungswert), und daraus
\begin{equation}
  \lambda_X(i) = \tfrac{1}{2}\,\mu_i\,\omega_i^2\,(\Delta r_X(i))^2.
  \label{eq:lambda-mode}
\end{equation}
Die Summe \"uber alle Modes,
$\Lambda_X^\text{total} = \sum_i \lambda_X(i)$, ist die thermische
Bindungs-Reorganisations-Energie der Bindung $X$. Bei $T \to 0$
reduziert sich das (mit $\mu = m_\text{mode}$) auf
\begin{equation}
  \lambda_X(i) \to \tfrac{\hbar\omega_i}{4} \cdot \alpha_X^2(i),
\end{equation}
wobei $\alpha_X(i)$ der dimensionslose ``Mode-Anteil'' der Bindung $X$
ist, implizit definiert durch
$\Delta r_X(i) = u_\text{rms}(i) \cdot \alpha_X(i)$. F\"ur eine reine
Streckmode der Bindung $X$ gilt $\alpha_X(i) = 1$, und der
Reorganisations-Beitrag bei $T \to 0$ ist $\hbar\omega/4$
(Nullpunkts-Energie geteilt durch 2). In der Praxis kann $\alpha_X(i)$
auch gr\"o\ss er als $1$ sein f\"ur Moden, in denen die zwei
Bindungspartner stark gegeneinander schwingen
(Relativ-Verschiebungs-Verst\"arkung); das ist geometrisch legitim
und verletzt kein Erhaltungsgesetz -- die Eigenvektor-Normierung
$\sum_i \lVert\evec_i\rVert^2 \approx 1$ ist eine Summe \"uber
\emph{Atome}, nicht \"uber Bindungen.

\paragraph{Bedeutung f\"ur PCET-Studien.}
\begin{itemize}\setlength{\itemsep}{0pt}
  \item $\Lambda_X^\text{total}$ ist eine wohldefinierte
        spektroskopische Observable. Sie misst, welche Bindung in den
        thermischen Schwingungen "`am st\"arksten zappelt"', mode-
        aufgel\"ost.
  \item Geeignet f\"ur \textbf{NRVS-Bandenidentifikation}:
        Frequenzbereiche mit gro{\ss}en
        $\Lambda_X(\omega)$-Spr\"ungen sind die Banden, in denen
        Bindung $X$ stark moduliert wird. Direkter Vergleich mit
        NRVS-Spektren m\"oglich.
  \item Geeignet f\"ur \textbf{System-Vergleiche bei gleicher
        Temperatur}: WT vs. Mut, prot vs. deprot. Skaliert konsistent.
        Insbesondere: bei deprot-Variante verschwinden NH und HA
        kanalbedingt, die anderen bleiben vergleichbar.
  \item Verwandt mit dem \textbf{Huang-Rhys-Faktor}
        \cite{HuangRhys1950} und der Franck-Condon-Theorie
        vibratorischer Spektren.
\end{itemize}

\paragraph{Was die Gr\"o{\ss}e nicht liefert.}
\begin{itemize}\setlength{\itemsep}{0pt}
  \item \emph{Keine} klassischen Marcus-Theorie-Aktivierungs-Energien
        f\"ur eine spezifische ET- oder PT-Reaktion -- daf\"ur br\"auchte
        man Reaktant- und Produkt-Geometrien getrennt.
  \item \emph{Keine} Reaktions-Raten-Vorhersage.
  \item \emph{Kein} direkter Vergleich mit Literatur-Werten der
        Marcus-$\lambda$ einer Reaktion (sofern dort statische
        Reaktant-Produkt-Differenzen verwendet werden).
\end{itemize}

\paragraph{Hammes-Schiffer-Sicht.}
Hammes-Schiffer und Soudackov \cite{HammesSchifferSoudackov2008}
zeigen, dass PCET in Proteinen typischerweise \emph{viele Moden
gleichzeitig} aktiviert. Die Vorstellung einer einzelnen
"`Promoter-Mode"' ist Spezialfall (kleine Modelle), nicht Regelfall.
Unsere mode-aufgel\"oste Sicht macht genau diese Multi-Mode-Realit\"at
explizit: statt einer Top-Mode-Klassifikation liefern wir die
Verteilung der Bindungs-Modulationen \"uber alle Modes.

\subsubsection*{F\"unf Bindungs-Kan\"ale}

Pro Mode werden f\"unf Bindungs-Kan\"ale ausgewertet:

\begin{itemize}\setlength{\itemsep}{0pt}
  \item \texttt{FeFe} -- Fe-Fe-Cluster-Atmung (ET-Reorganisation).
  \item \texttt{FeN} -- Fe-N(His)-Bindung (PT-relevant in
        Rieske-artigen Systemen).
  \item \texttt{FeS} -- Fe-S(Cys)-Bindung sowie Fe-S$_\text{bridge}$
        innerhalb des Clusters.
  \item \texttt{NH} -- N-H-Stretchung der protonierten His.
  \item \texttt{HA} -- H-Akzeptor-Distanz (Wasserstoffbr\"ucke), berechnet
        im \emph{Reaktionskoordinaten-Modus} (siehe unten).
\end{itemize}

NH und HA werden bei deprotonierten His-Liganden \"ubersprungen
(kein H verf\"ugbar). FeS umfasst sowohl Fe-Cys-Bindungen als auch
Cluster-interne Fe-$\mu$S-Bindungen; die Aufschl\"usselung pro einzelne
Bindung erfolgt im Sheet \texttt{Reorg\_pro\_Bindung}.

\subsubsection*{Geometrische Modulation $\Delta r_X$}

F\"ur eine Mode mit thermisch skaliertem Auslenkungsvektor
$\mathbf{e}_a, \mathbf{e}_b$ der beiden Bindungspartner (bei
Position $\mathbf{r}_a, \mathbf{r}_b$) ist die vorzeichenbehaftete
Bindungsl\"angen-Modulation
\begin{equation}
  \Delta r_X^\text{(klass.)}
  = (\mathbf{e}_a - \mathbf{e}_b) \cdot \hat{\mathbf{n}}_{ab},
  \qquad
  \hat{\mathbf{n}}_{ab} = \frac{\mathbf{r}_a - \mathbf{r}_b}{\|\mathbf{r}_a - \mathbf{r}_b\|}.
  \codref{reorganization.signed\_dr\_along\_axis}
  \label{eq:dr-classical}
\end{equation}
Konvention: $\Delta r > 0$ heisst "`Bindung dehnt sich"'.
Diese klassische Definition wird f\"ur FeFe, FeN, FeS und NH verwendet.

\paragraph{HA als Reaktionskoordinate.}
F\"ur den HA-Kanal misst Gleichung~\ref{eq:dr-classical} \emph{nicht}
das, was f\"ur PT physikalisch relevant ist: bei einer hochfrequenten
C=O-Schwingung im Protein-R\"uckgrat bewegt sich der Akzeptor stark,
das H bleibt fest, aber die Distanz $|\mathbf{r}_H - \mathbf{r}_A|$
\"andert sich -- und w\"urde als HA-Modulation gez\"ahlt, obwohl es
keinen PT-Charakter hat.

Die HA-Modulation misst die
\emph{Bewegung des H relativ zum Donor-N}, projiziert auf die
N$\to$Akzeptor-Achse:
\begin{equation}
  \Delta r_\text{HA}^\text{(react.\ coord.)}
  = (\mathbf{e}_H - \mathbf{e}_N) \cdot \hat{\mathbf{n}}_{NA},
  \qquad
  \hat{\mathbf{n}}_{NA} = \frac{\mathbf{r}_A - \mathbf{r}_N}{\|\mathbf{r}_A - \mathbf{r}_N\|}.
  \codref{reorganization.compute\_mode\_modulations}
  \label{eq:dr-reaction-coord}
\end{equation}
Wenn nur der Akzeptor sich bewegt (Skelett-Mode), ist
$\mathbf{e}_H - \mathbf{e}_N \approx 0$ und damit
$\Delta r_\text{HA} = 0$ -- kein PT-Beitrag. Eine echte PT-Bewegung
(H bewegt sich vom N zum A) liefert dagegen $\Delta r_\text{HA} > 0$.

\paragraph{Hauptakzeptor pro H.}
Pro His-H-Atom werden zwar alle Akzeptoren in
\texttt{pcet\_hbond\_cutoff\_a} (Default 4\,\AA) gefunden, aber nur
der \emph{Hauptakzeptor} (h\"ochstes Gauss-Gewicht
$\exp\bigl(-(r_{\text{eq}} - r_0)^2/(2\sigma^2)\bigr)$ mit
$r_0 = 2.8$\,\AA, $\sigma = 0.4$\,\AA) wird als HA-Kanal verwendet.
Damit wird die Doppelz\"ahlung von Akzeptoren vermieden -- ein
H bewegt sich physikalisch nur in eine Richtung. Konkret f\"ur PCET ist
die Logik: pro His-Imidazol gibt es ein NE2-H (oder ND1-H), und an
diesem H haengt genau ein dominanter H-Bridge-Akzeptor (typischerweise
ein H$_2$O, Carbonyl-O eines Backbones, oder ein Asp/Glu-O).

\paragraph{Wirkung des Gauss-Gewichts.}
Das Gauss-Gewicht wird
$w = \exp\bigl(-(r_{\text{eq}} - r_0)^2/(2\sigma^2)\bigr)$
\emph{ausschlie{\ss}lich} f\"ur die Hauptakzeptor-Auswahl verwendet --
nicht mehr als D\"ampfungs-Faktor in der Lambda-Aggregation. Die binaere
Logik ist nun: \emph{ist ein H-Bridge-Akzeptor in
\texttt{pcet\_hbond\_cutoff\_a} Reichweite oder nicht?} Bei Akzeptoren
in Reichweite z\"ahlt die HA-Mode voll (weight $= 1.0$), bei keinem
entf\"allt der Kanal. Die geometrische D\"ampfung weiter Akzeptoren
ergibt sich automatisch aus den kleineren $\Delta r_{HA}$-Werten
(weite H-Bridges modulieren ohnehin schwach), zus\"atzliche
Multiplikation mit $w < 1$ w\"are doppelte Strafe.

\subsubsection*{Reduzierte Masse: zwei Konventionen}

F\"ur den $\lambda$-Beitrag pro Mode haben wir zwei Wahlm\"oglichkeiten
f\"ur die reduzierte Masse:

\begin{itemize}
  \item $\mu_\text{pair}$ -- bindungsspezifische reduzierte Masse,
        berechnet aus den beiden gebundenen Atomen
        ($\mu_\text{NH} \approx 0.94$\,amu,
        $\mu_\text{FeFe} \approx 27.92$\,amu usw.). Naehere
        \emph{isolierter Oszillator}: gut f\"ur Mode-zu-Mode-Vergleich
        \emph{derselben} Bindung.

  \item $\mu_\text{mode}$ -- Mode-eigene reduzierte Masse aus dem
        Gaussian-Output. Diese erf\"ullt per Konstruktion die
        Orthogonalit\"at der Eigenmodes, wodurch sich die
        $\lambda$-Beitr\"age \emph{ohne Doppelz\"ahlung}
        addieren -- das ist die \textbf{Marcus-Hush-korrekte} Wahl
        f\"ur die Summe in Gl.~\ref{eq:marcus-hush}.
\end{itemize}

Beide werden ausgegeben (Spalten \texttt{lambda\_X\_pair\_cm1}
und \texttt{lambda\_X\_mode\_cm1} im Sheet
\texttt{Reorganisationsenergie}). F\"ur die System-Aggregation in
\texttt{Reorg\_Total} wird $\mu_\text{mode}$ verwendet.

\subsubsection*{Reorg-Beitrag $\lambda_X$ pro Mode}

Die schon in \eqref{eq:lambda-mode} oben gegebene Formel
$\lambda_X(i) = \tfrac{1}{2}\,\mu_i\,\omega_i^2\,(\Delta r_X(i))^2$
wird mit den oben definierten geometrischen Modulationen
$\Delta r_X(i)$ und der gew\"ahlten reduzierten Masse-Konvention
($\mu_\text{pair}$ oder $\mu_\text{mode}$) angewandt.

Einheiten-Konversion: $\omega_\text{SI} = 2\pi c\,\omega_{\text{cm}^{-1}}$,
$\mu_\text{SI} = \mu_{\text{amu}} \cdot \text{u}$,
$\Delta r_\text{SI} = \Delta r_{\text{\AA}} \cdot 10^{-10}$\,m;
$\lambda_{\text{cm}^{-1}} = \lambda_\text{J} / (h c)$.

Wichtig: Wegen der thermischen Skalierung
$\mathbf{e}_a = \mathbf{e}_a^\text{normalized} \cdot u_\text{rms}$
liegt $\Delta r_X(i)$ skalenm\"a{\ss}ig im Bereich der QHO-Wahrscheinlichkeits-
Verteilung, nicht der Reaktant-Produkt-Differenz. F\"ur reine
Streckmoden der Bindung $X$ konvergiert $\lambda_X(i)$ bei $T \to 0$
gegen $\hbar\omega_i/4$ ($= \omega_i/4$ in cm$^{-1}$); bei gemischten
Modes wird das mit dem quadrierten Mode-Anteil $\alpha_X^2$ multipliziert.

\subsubsection*{System-Aggregation}

\begin{equation}
  \Lambda_X^\text{total} = \sum_{i=1}^{N_\text{modes}} \lambda_X(i)
  \codref{reorganization.compute\_total\_reorganization}
  \label{eq:lambda-total}
\end{equation}

pro Kanal $X \in \{\text{FeFe, FeN, FeS, NH, HA}\}$.

Die Summen werden zus\"atzlich \emph{frequenzaufgel\"ost} ausgegeben:

\begin{itemize}
  \item \textbf{Modulations-Spektren} $M_X(\omega)$ -- Gauss-
        verbreiterte Akkumulation der Bindungs-Modulationen je Mode
        entlang der Frequenzachse:
        \begin{equation}
          M_X(\omega)
          = \sum_{i=1}^{N_\text{modes}}
            \Delta r_X^\text{RSS}(i)\,
            \mathcal{G}(\omega - \omega_i, \sigma_M),
          \quad
          \mathcal{G}(x,\sigma) = \frac{e^{-x^2/(2\sigma^2)}}{\sigma\sqrt{2\pi}}.
          \codref{reorganization.compute\_modulation\_spectra}
          \label{eq:M-spectrum}
        \end{equation}
        Default-Verbreiterung: $\sigma_M = 5$\,cm$^{-1}$, einstellbar
        via \texttt{reorg\_spectrum\_sigma\_cm1}.

        \medskip
        \noindent\textbf{Was ist $\Delta r_X^\text{RSS}(i)$ genau?}
        Jeder Eltern-Kanal
        $X \in \{\text{FeFe}, \text{FeN}, \text{FeS}, \text{NH}, \text{HA}\}$
        kann mehrere \emph{Sub-Kan\"ale} haben: z.\,B.\ enth\"alt FeN
        die beiden Fe-N(His$_a$)- und Fe-N(His$_b$)-Sub-Bindungen, FeS
        enth\"alt alle vier Fe-S(Cys$_k$)-Sub-Bindungen (bzw.\ zwei
        f\"ur Cys$_2$His$_2$-Cluster). F\"ur Mode $i$ aggregieren wir
        daher die Sub-Kanal-Modulationen als gewichtete
        Wurzel-Quadratsumme
        \begin{equation}
          \Delta r_X^\text{RSS}(i)
            \;=\; \sqrt{\sum_{s \in \text{Sub-Kan\"alen}(X)}
                          w_s\,(\Delta r_X^{(s)}(i))^2},
          \codref{reorganization.aggregate\_by\_parent}
          \label{eq:dr-rss}
        \end{equation}
        wobei die Gewichte $w_s$ aus
        \texttt{build\_channels\_from\_coord\_info} stammen
        (Gauss-Gewichtung des H$\cdot\cdot\cdot$Akzeptor-Abstands
        f\"ur HA, eins f\"ur FeFe/FeN/FeS/NH). F\"ur Eltern-Kan\"ale
        mit nur einem Sub-Kanal (FeFe, NH, HA-pro-Akzeptor) reduziert
        sich das zu $\sqrt{w}\,|\Delta r_X(i)|$, also wieder die
        einfache absolute Modulation. Die RSS-Aggregation ist die
        \emph{energetisch konsistente} Wahl: sie erf\"ullt
        $\lambda_X^\text{Eltern}(i) = \tfrac{1}{2}\mu_i\omega_i^2
        \,(\Delta r_X^\text{RSS}(i))^2$, wobei
        $\lambda_X^\text{Eltern}$ die gewichtete Summe der
        Sub-Kanal-$\lambda$'s ist.

  \item \textbf{Co-Modulations-Spektren} -- geometrisches Mittel
        zweier $M_X$-Spektren, hoch nur dort, wo \emph{gleichzeitig}
        zwei Bindungen moduliert werden:
        \begin{align}
          C_\text{PCET}(\omega) &= \sqrt{M_\text{HA}(\omega)\, M_\text{FeFe}(\omega)},\\
          C_{\text{PT-FeN}}(\omega) &= \sqrt{M_\text{HA}(\omega)\, M_\text{FeN}(\omega)},\\
          C_{\text{ET-FeS}}(\omega) &= \sqrt{M_\text{FeFe}(\omega)\, M_\text{FeS}(\omega)}.
        \end{align}

  \item \textbf{Kumulative Reorg-Kurven} $\Lambda_X(\omega)$ -- monoton
        steigend, konvergieren am Spektrum-Ende gegen $\Lambda_X^\text{total}$:
        \begin{equation}
          \Lambda_X(\omega_\text{cut})
          = \sum_{i\,:\,\omega_i \le \omega_\text{cut}} \lambda_X(i).
        \end{equation}
        Spr\"unge in $\Lambda_X(\omega)$ markieren reorg-aktive
        Frequenzfenster, die direkt mit NRVS-Spektren verglichen
        werden k\"onnen.
\end{itemize}

\subsubsection*{Excel-Output}
\label{sec:reorg-excel}

Vier Sheets pro Lauf bzw.\ Frequenzfenster:

\begin{itemize}
  \item \texttt{Reorganisationsenergie} -- pro Mode 21 Spalten:
        \texttt{freq\_cm1} und f\"ur jeden Kanal vier Werte
        ($\Delta r_\text{rms}$, $\Delta r_\text{sum-signed}$,
        $\lambda_\text{pair}$, $\lambda_\text{mode}$).

  \item \texttt{Reorg\_Total} -- 5-Zeilen-Zusammenfassung: pro Kanal
        $\Lambda_X^\text{total}$ in beiden $\mu$-Konventionen plus
        Anzahl der beitragenden Modes.

  \item \texttt{Reorg\_pro\_Bindung} -- pro einzelner Bindung
        (FeS\_Cys207, FeS\_Cys216, FeS\_Cluster\_*, FeN\_His255 usw.)
        deren Beitrag, mit Akzeptor-Gewicht und Anzahl Modes.

  \item \texttt{Modulations\_Spektren} -- frequenzaufgel\"ost: f\"unf
        $M_X(\omega)$-Kurven plus drei Co-Indikator-Kurven
        $C_\text{PCET}, C_\text{PT-FeN}, C_\text{ET-FeS}$.

  \item \texttt{Lambda\_kumulativ} -- f\"unf $\Lambda_X(\omega)$-Kurven
        f\"ur direkten Vergleich mit NRVS-Spektren.
\end{itemize}


\subsection{HP/Standard-Frequenz-Konsistenzcheck}
\label{sec:hp-std}

Wie in \S\ref{sec:hp-std-check} angedeutet, schreibt Gaussian
bei \texttt{freq=hpmodes} die Frequenzen \emph{zweimal} in das
Logfile -- einmal im HP-Block (5 Nachkommastellen) und einmal im
Standard-Block (2 Nachkommastellen). Beide Bl\"ocke beziehen sich auf
dieselbe Eigenwertberechnung; ihre Frequenzen \emph{m\"ussen} also
\"ubereinstimmen (innerhalb der Schreibgenauigkeit von 0.01\,cm$^{-1}$).

\subsubsection*{Implementierung}

\texttt{logio.check\_hp\_std\_frequency\_consistency} vergleicht beim
Lauf:

\begin{enumerate}
  \item Mode-f\"ur-Mode: f\"ur jede Modennummer wird die HP-Frequenz
        mit der Standard-Frequenz verglichen.
  \item Toleranz: 0.01\,cm$^{-1}$ (default), verstellbar via
        \texttt{cfg.hp\_std\_freq\_tolerance\_cm1}.
  \item Bei Abweichungen wird die maximale Differenz, die betroffene
        Modennummer und beide Frequenzen geloggt.
\end{enumerate}

\subsubsection*{Was Abweichungen bedeuten}

\begin{center}
\begin{tabularx}{\textwidth}{l X}
\toprule
\textbf{Symptom} & \textbf{M\"ogliche Ursache} \\
\midrule
1--2 Modes mit $\Delta\omega > 0.01$ &
  Numerische Reduzierung in Gaussian-Output (selten, harmlos) \\
Viele Modes systematisch versetzt &
  Falsches Block-Pairing (HP-Block aus anderem Job, Restart-Mismatch) \\
Komplette Diskrepanz &
  Abgebrochener Job, partielles Logfile, oder zwei verschiedene
  Logfiles aneinandergeh\"angt \\
\bottomrule
\end{tabularx}
\end{center}

In allen F\"allen wird die Diskrepanz \emph{nicht} fatal behandelt --
das Programm l\"auft weiter und nutzt die HP-Werte. Die Diagnose
landet aber im Befund, sodass eine Nachuntersuchung m\"oglich ist.

\begin{warnbox}
\textbf{Praxis-Hinweis.} In sehr alten Logfiles (Gaussian 09 oder
\"alter) gibt es f\"ur die niedrigsten Frequenzen ($\omega \to 0$)
gelegentlich Vorzeichen-Diskrepanzen zwischen HP- und Standard-Block
(eine Mode, die im HP-Block als $-2.3$\,cm$^{-1}$ und im Standard-Block
als $+2.3$\,cm$^{-1}$ erscheint). Das ist ein Gaussian-Artefakt der
Translationen/Rotationen-Behandlung. Die Konsistenzpr\"ufung warnt
darauf, das Verhalten ist aber harmlos -- diese Modes werden in der
Regel ohnehin wegen $\omega < $ Threshold ausgeschlossen.
\end{warnbox}

% =========================================================================
% Bibliographie
% =========================================================================

\clearpage
\section*{Referenzen}
\addcontentsline{toc}{section}{Referenzen}
% =========================================================================

\part{Anwendung}
% =========================================================================

% =========================================================================
\section{Konfiguration}
\label{sec:konfig}

Alle Lauf-Parameter werden in einem ``Config''-Objekt zusammengefasst.
``Config'' ist eine Python-``dataclass'' mit Default-Werten f\"ur alle
optionalen Felder (\texttt{src/modenanalyse\_2fe2s/config.py}). Drei
Felder sind \emph{Pflicht}: \texttt{log\_file}, \texttt{output\_dir},
sowie -- falls Liganden-Aufl\"osung gew\"unscht ist --
\texttt{pdb\_file}.

\subsection{Beispiel-Konfiguration}

Die typische, in Spyder lauff\"ahige Form (siehe ``example\_run.py''):

\begin{lstlisting}[style=python]
from modenanalyse_2fe2s import Config, run_analysis

cfg = Config(
    log_file   = r"D:\Daten\cluster_freq.log",
    pdb_file   = r"D:\Daten\cluster.pdb",
    output_dir = r"D:\Daten\results",
    temp_k     = 40.0,
    freq_max   = 800.0,
    freq_windows = [(0, 100), (100, 300), (300, 500), (500, 700)],
)

run_analysis(cfg)
\end{lstlisting}

Alle weiteren Felder behalten ihren Default. Die folgenden Tabellen
listen die Felder gruppiert nach Funktion auf.

\subsection{Vollst\"andige Parameter-Referenz}
\label{sec:param-referenz}

Die folgende \"Ubersicht zeigt alle 48 Konfigurationsfelder mit ihren
Default-Werten. Ausf\"uhrliche Beschreibungen und Empfehlungen siehe
unten in den Themen-Abschnitten.

\begin{longtable}{p{0.32\textwidth} p{0.20\textwidth} p{0.40\textwidth}}
\caption{Vollst\"andige \"Ubersicht aller Konfigurationsfelder mit
Default-Werten. Eine kommentierte Beispiel-TOML mit allen Feldern liegt
unter \texttt{run\_example\_full.toml} im Distributions-Verzeichnis.}
\label{tab:param-overview} \\
\toprule
\textbf{Feld} & \textbf{Default} & \textbf{Sektion} \\
\midrule
\endfirsthead
\toprule
\textbf{Feld} & \textbf{Default} & \textbf{Sektion} \\
\midrule
\endhead
% ── Eingabe ─────────────────────────────────────────────
\texttt{log\_file}                  & --- (Pflicht)        & Eingabe \\
\texttt{output\_dir}                & --- (Pflicht)        & Eingabe \\
\texttt{pdb\_file}                  & \texttt{""}          & Eingabe \\
\texttt{pdb\_chain}                 & \texttt{"A"}         & Eingabe \\
\texttt{logname\_suffix}            & \texttt{""}          & Eingabe \\
\midrule
% ── Frequenz ────────────────────────────────────────────
\texttt{freq\_min}                  & \texttt{None}        & Frequenz \\
\texttt{freq\_max}                  & \texttt{None}        & Frequenz \\
\texttt{freq\_windows}              & \texttt{None}        & Frequenz \\
\midrule
% ── Thermo ──────────────────────────────────────────────
\texttt{temp\_k}                    & \texttt{None}        & Thermo \\
\texttt{amplitude}                  & \texttt{1.0}         & Thermo \\
\midrule
% ── Cluster ─────────────────────────────────────────────
\texttt{cluster\_index}             & \texttt{0}           & Cluster \\
\texttt{fe\_s\_cutoff}              & \texttt{3.0} \AA     & Cluster \\
\texttt{fe\_fe\_cutoff}             & \texttt{3.5} \AA     & Cluster \\
\texttt{strict\_cluster}            & \texttt{False}       & Cluster \\
\midrule
% ── Ligand ──────────────────────────────────────────────
\texttt{fe\_coord\_cutoff\_n}       & \texttt{2.50} \AA    & Ligand \\
\texttt{fe\_coord\_cutoff\_s}       & \texttt{2.70} \AA    & Ligand \\
\texttt{fe\_coord\_cutoff\_o}       & \texttt{2.50} \AA    & Ligand \\
\texttt{include\_hn\_vibration}     & \texttt{True}        & Ligand \\
\midrule
% ── OOP/INP ─────────────────────────────────────────────
\texttt{mode\_type\_threshold}      & \texttt{60.0\,\%}    & OOP \\
\texttt{mode\_type\_detail\_thresholds} & \texttt{(60, 75, 90)} & OOP \\
\midrule
% ── SCSD/SSE ─────────────────────────────────────────────
\texttt{analyze\_scsd}              & \texttt{True}        & SCSD \\
\texttt{analyze\_sse}                & \texttt{True}        & SCSD \\
\texttt{sse\_chain}                  & \texttt{""}          & SCSD \\
\midrule
% ── PCET / Marcus-Hush-Reorg ────────────────────────────
\texttt{pcet\_enabled}              & \texttt{True}        & PCET \\
\texttt{pcet\_hbond\_cutoff\_a}     & \texttt{4.0} \AA     & PCET \\
\texttt{pcet\_acceptor\_r0\_a}      & \texttt{2.8} \AA     & PCET \\
\texttt{pcet\_acceptor\_sigma\_a}   & \texttt{0.4} \AA     & PCET \\
\texttt{reorg\_spectrum\_sigma\_cm1} & \texttt{5.0} cm$^{-1}$  & PCET \\
\texttt{reorg\_spectrum\_step\_cm1} & \texttt{0.5} cm$^{-1}$  & PCET \\
\midrule
% ── Embedding ───────────────────────────────────────────
\texttt{export\_embedding\_plots}   & \texttt{True}        & Embedding \\
\texttt{umap\_n\_neighbors}         & \texttt{None} (auto) & Embedding \\
\midrule
% ── Numerik ────────────────────────────────────────────
\texttt{include\_hydrogen}          & \texttt{True}        & Numerik \\
\texttt{sigma\_eigvec}              & \texttt{5e-4}        & Numerik \\
\texttt{sigma\_coord}               & \texttt{1e-3}        & Numerik \\
\texttt{amplitude\_threshold}       & \texttt{0.001}       & Numerik \\
\texttt{coord\_match\_tol}          & \texttt{1.0} \AA     & Numerik \\
\texttt{scan\_chunk\_mb}            & \texttt{16}          & Numerik \\
\texttt{show\_errors\_in\_excel}    & \texttt{True}        & Numerik \\
\texttt{interp\_step}               & \texttt{0.05} cm$^{-1}$ & Numerik \\
\texttt{interp\_boundary\_mode}     & \texttt{"context"}   & Numerik \\
\texttt{interp\_edge\_extend}       & \texttt{0.5} cm$^{-1}$ & Numerik \\
\texttt{interp\_context\_cm1}       & \texttt{5.0} cm$^{-1}$ & Numerik \\
\texttt{use\_cache}                 & \texttt{True}        & Numerik \\
\bottomrule
\end{longtable}



\subsection{Eingabe und Ausgabe}

\begin{longtable}{p{0.27\textwidth} p{0.16\textwidth} p{0.50\textwidth}}
\toprule
\textbf{Feld} & \textbf{Default} & \textbf{Beschreibung} \\
\midrule
\endhead
\texttt{log\_file}      & \emph{(Pflicht)} & Pfad zur DFT-Frequenz-
  Datei. Unterst\"utzt: Gaussian-16 \texttt{.log} (mit/ohne
  \texttt{freq=hpmodes}) und ORCA \texttt{.hess} (automatisch
  erkannt an der Endung). \\
\texttt{output\_dir}    & \emph{(Pflicht)} & Zielverzeichnis. Wird
  angelegt falls nicht vorhanden. \\
\texttt{pdb\_file}      & \texttt{""}     & PDB-Datei mit demselben
  Cluster (Liganden-Aufl\"osung; optional, aber empfohlen). \\
\texttt{pdb\_chain}     & \texttt{"A"}     & Chain-Identifier in der
  PDB-Datei. \\
\texttt{logname\_suffix} & \texttt{""}     & Optionales Suffix f\"ur die
  Output-Dateien (n\"utzlich bei mehreren L\"aufen aus demselben Logfile,
  z.\,B.\ verschiedene Cluster). \\
\bottomrule
\end{longtable}

\subsection{Frequenzbereich und Fenster}

\begin{longtable}{p{0.27\textwidth} p{0.16\textwidth} p{0.50\textwidth}}
\toprule
\textbf{Feld} & \textbf{Default} & \textbf{Beschreibung} \\
\midrule
\endhead
\texttt{freq\_min}      & \texttt{None} & Untere Frequenzgrenze in
  cm$^{-1}$. \texttt{None} bedeutet keine untere Grenze. \\
\texttt{freq\_max}      & \texttt{None} & Obere Frequenzgrenze.
  F\"ur NRVS-relevante Auswertung typischerweise 800. \\
\texttt{freq\_windows}  & \texttt{None} & Liste von Fenstern wie
  \texttt{[(0, 100), (100, 300), \ldots]}. Wenn gesetzt, generiert das
  Programm \emph{ein Excel pro Fenster} (Multi-Window-Modus). \\
\bottomrule
\end{longtable}

\begin{infobox}
\textbf{Konvention.} \texttt{freq\_max} (und \texttt{freq\_min}) gelten
auch f\"ur s\"amtliche Reorg- und Aggregat-Auswertungen. Da
$^{57}$Fe-projizierte Beitr\"age oberhalb $\sim$800 cm$^{-1}$
verschwindend klein sind und NRVS-Experimente in dem Bereich messen,
ist eine 800-cm$^{-1}$-Grenze sowohl experimentell als auch numerisch
sinnvoll.
\end{infobox}

\begin{infobox}
\textbf{Fenstergrenzen-Konvention (seit v1.0.4).}
Der Multi-Window-Modus partitioniert Modes in Bins
\texttt{[(lo$_1$,hi$_1$), (lo$_2$,hi$_2$), \ldots, (lo$_N$,hi$_N$)]}.
Die Grenzen sind \emph{halb-offen von links}:
\begin{itemize}\setlength{\itemsep}{0pt}
  \item Jedes Fenster au\ss er dem letzten verwendet das Intervall
        $[\,\text{lo},\, \text{hi}\,)$. Eine Mode mit Frequenz exakt
        auf einer Fenstergrenze geh\"ort zum \emph{oberen} Fenster
        (dem, das bei diesem Wert beginnt).
  \item Das letzte Fenster der Liste verwendet
        $[\,\text{lo},\, \text{hi}\,]$, sodass seine obere Kante
        eingeschlossen ist. Falls \texttt{hi} f\"ur das letzte
        Fenster \texttt{None}/\texttt{inf} ist, ist es nach oben
        offen.
\end{itemize}
Damit wird garantiert, dass jede Mode in genau einem Fenster z\"ahlt,
und Doppelz\"ahlungen von $B$-Faktor-Beitr\"agen und
$\Lambda_X(\omega_\text{cut})$-Werten f\"ur Grenz-Modes
verhindert (\S\ref{sec:migration-v104}). Vor v1.0.4 wurden an beiden
Enden geschlossene Intervalle verwendet, die exakt auf der Grenze
liegende Modes doppelt z\"ahlten.
\end{infobox}

\subsection{Thermik}

\begin{longtable}{p{0.27\textwidth} p{0.16\textwidth} p{0.50\textwidth}}
\toprule
\textbf{Feld} & \textbf{Default} & \textbf{Beschreibung} \\
\midrule
\endhead
\texttt{temp\_k}       & \texttt{None} & Probentemperatur in Kelvin
  f\"ur die thermische Skalierung der Eigenvektoren. \texttt{None}
  $\to$ Klassik-Modus mit Amplitude $=$ \texttt{amplitude}. Typische
  NRVS-Werte: 40, 80, 100. \\
\texttt{temp\_k\_nis}  & \texttt{[40]} & Liste der Temperaturen f\"ur
  NIS-Spektrum-Simulation. Pro Temperatur wird eine Spalte im NIS-Sheet
  geschrieben. \\
\texttt{amplitude}      & \texttt{1.0}  & Fallback-Amplitude in
  \AA{} bei \texttt{temp\_k = None}. Default 1.0 ist die historische
  Konvention. \\
\bottomrule
\end{longtable}

\subsection{Cluster-Erkennung}

\begin{longtable}{p{0.27\textwidth} p{0.16\textwidth} p{0.50\textwidth}}
\toprule
\textbf{Feld} & \textbf{Default} & \textbf{Beschreibung} \\
\midrule
\endhead
\texttt{fe\_s\_cutoff}  & \texttt{3.0} & Maximaler Fe-S-Abstand (\AA{}),
  innerhalb dessen S-Atome zum Cluster z\"ahlen. \\
\texttt{fe\_fe\_cutoff} & \texttt{3.5} & Maximaler Fe-Fe-Abstand (\AA{})
  f\"ur die Cluster-Identifikation. \\
\texttt{cluster\_index} & \texttt{0}   & Bei Multi-Cluster-Systemen
  (z.\,B.\ Glutaredoxin-Dimer): Index des zu analysierenden Clusters.
  Default 0 = engstes Fe-Fe-Paar. Pro Cluster ein separater Lauf. \\
\texttt{strict\_cluster} & \texttt{False} & Wenn \texttt{True},
  bricht das Programm mit Fehler ab, wenn der Cluster nicht
  eindeutig identifiziert werden kann (kein Fallback auf nahegelegene
  Atome). \\
\bottomrule
\end{longtable}

\subsection{Fe-Liganden-Erkennung (PDB-basiert)}

\begin{longtable}{p{0.27\textwidth} p{0.16\textwidth} p{0.50\textwidth}}
\toprule
\textbf{Feld} & \textbf{Default} & \textbf{Beschreibung} \\
\midrule
\endhead
\texttt{fe\_coord\_cutoff\_n} & \texttt{2.50} & Maximaler Fe-N-Abstand
  (\AA{}) f\"ur Liganden-N (His). \\
\texttt{fe\_coord\_cutoff\_s} & \texttt{2.70} & Maximaler Fe-S-Abstand
  f\"ur Liganden-S (Cys). \\
\texttt{fe\_coord\_cutoff\_o} & \texttt{2.50} & Maximaler Fe-O-Abstand
  f\"ur Liganden-O (Asp/Glu, selten f\"ur {[}2Fe-2S{]}). \\
\texttt{include\_hn\_vibration} & \texttt{True} & N-H-Streckschwingungs-
  Analyse f\"ur protonierte His-Liganden (Sheet \texttt{His\_HN}).
  Erfordert \texttt{include\_hydrogen=True}. \\
\bottomrule
\end{longtable}

\subsection{OOP/INP-Klassifikation}

\begin{longtable}{p{0.34\textwidth} p{0.16\textwidth} p{0.43\textwidth}}
\toprule
\textbf{Feld} & \textbf{Default} & \textbf{Beschreibung} \\
\midrule
\endhead
\texttt{mode\_type\_threshold} & \texttt{60.0} & Bin\"are OOP-Schwelle in
  Prozent (\S\ref{sec:oop-inp}). \\
\texttt{mode\_type\_detail\_thresholds} & \texttt{(60, 75, 90)} & Drei
  Schwellen f\"ur die feine 7-Stufen-Klassifikation. Symmetrisch zu 50\,\%
  angewandt. \\
\bottomrule
\end{longtable}

\subsection{SCSD und Sekund\"arstruktur}

\begin{longtable}{p{0.27\textwidth} p{0.16\textwidth} p{0.50\textwidth}}
\toprule
\textbf{Feld} & \textbf{Default} & \textbf{Beschreibung} \\
\midrule
\endhead
\texttt{analyze\_scsd}  & \texttt{True} & Symmetry-Coordinate
  Structural Decomposition aktivieren (erfordert \texttt{scsdpy}). \\
\texttt{analyze\_sse}    & \texttt{True} & Sekund\"arstruktur-Analyse
  aktivieren (Helices, Sheets aus PDB). \\
\texttt{sse\_chain}      & \texttt{""}   & Chain-Filter f\"ur
  SSE-Erkennung. Leer = wie \texttt{pdb\_chain}. \\
\bottomrule
\end{longtable}

\subsection{PCET / Marcus-Hush-Reorganisation}

\begin{longtable}{p{0.34\textwidth} p{0.16\textwidth} p{0.43\textwidth}}
\toprule
\textbf{Feld} & \textbf{Default} & \textbf{Beschreibung} \\
\midrule
\endhead
\texttt{pcet\_enabled} & \texttt{True} & Marcus-Hush-Reorg-Pipeline
  aktivieren. Wenn \texttt{False}, werden $\Delta r_X$ und $\lambda_X$
  pro Mode nicht berechnet, alle Reorg-Sheets entfallen. \\
\texttt{pcet\_hbond\_cutoff\_a} & \texttt{4.0} & H...Akzeptor-Maximal-
  distanz (\AA{}) f\"ur die Akzeptor-Suche. Gemessen vom H-Atom. \\
\texttt{pcet\_acceptor\_r0\_a} & \texttt{2.8} & Optimal-Abstand
  (\AA{}) f\"ur die Gauss-Gewichtung der Akzeptor-Auswahl;
  entspricht typischer N--H$\cdots$O/N-Bindungsl\"ange. \\
\texttt{pcet\_acceptor\_sigma\_a} & \texttt{0.4} & Gauss-Standard-
  abweichung (\AA{}) der Akzeptor-Gewichtung. Bei
  $|r - r_0| > 2\sigma$ tr\"agt der Akzeptor faktisch nichts mehr bei. \\
\texttt{reorg\_spectrum\_sigma\_cm1} & \texttt{5.0} & Gauss-Verbreiterung
  der Modulations-Spektren $M_X(\omega)$. 5\,cm$^{-1}$ entspricht
  typischer DFT-Frequenz-Unsch\"arfe + NRVS-Aufl\"osung. \\
\texttt{reorg\_spectrum\_step\_cm1} & \texttt{0.5} & Frequenz-Raster-
  Schrittweite f\"ur $M_X(\omega)$ und $\Lambda_X(\omega)$. 0.5\,cm$^{-1}$
  ergibt feine Banden bei vertretbarer Datei-Gr\"o{\ss}e. \\
\bottomrule
\end{longtable}

\subsection{Numerische Parameter}

\begin{longtable}{p{0.27\textwidth} p{0.16\textwidth} p{0.50\textwidth}}
\toprule
\textbf{Feld} & \textbf{Default} & \textbf{Beschreibung} \\
\midrule
\endhead
\texttt{include\_hydrogen} & \texttt{True} & H-Atome in Eigenvektor-
  Reads ber\"ucksichtigen. F\"ur PCET zwingend. \\
\texttt{sigma\_eigvec}    & \texttt{5e-4}  & Numerisches Rauschen
  in Eigenvektor-Komponenten (Bernoulli-Var.). Kontrolliert Schwellen
  f\"ur ``unbedeutende'' Beitr\"age. \\
\texttt{sigma\_coord}     & \texttt{1e-3} & Koordinatenrauschen
  in \AA{}. \\
\texttt{amplitude\_threshold} & \texttt{0.001} & Schwelle f\"ur
  signifikante Atomverschiebung in einem Mode. \\
\texttt{coord\_match\_tol} & \texttt{1.0}   & Toleranz (\AA{}) f\"ur
  PDB$\leftrightarrow$Gaussian-Atom-Matching nach Kabsch-Alignment. \\
\texttt{interp\_step}     & \texttt{0.05}  & Schrittweite (cm$^{-1}$)
  fuer interpolierten pDOS-Export. \\
\texttt{interp\_boundary\_mode} & \texttt{"context"} & Verhalten am
  Fenster-Rand: \texttt{"context"} (mit Pufferzone), \texttt{"hard"}
  (scharfer Schnitt), \texttt{"extend"} (linear extrapoliert). \\
\texttt{interp\_edge\_extend} & \texttt{0.5} & Pufferzone-Breite (cm$^{-1}$
  bzw.\ Anteil) im \texttt{"context"}-Modus. \\
\texttt{scan\_chunk\_mb}  & \texttt{16}    & Block-Gr\"o\ss e f\"ur
  Logfile-Scan in MB. F\"ur sehr gro\ss e Logfiles ($>$1 GB) kann
  Erh\"ohen das Lesen beschleunigen. \\
\texttt{use\_cache}       & \texttt{True}  & Aktiviert Datei-Caching f\"ur
  Block-Offsets im Logfile. Halbiert die Laufzeit bei wiederholten Calls. \\
\bottomrule
\end{longtable}

\subsection{Embeddings (UMAP)}

\begin{longtable}{p{0.27\textwidth} p{0.16\textwidth} p{0.50\textwidth}}
\toprule
\textbf{Feld} & \textbf{Default} & \textbf{Beschreibung} \\
\midrule
\endhead
\texttt{umap\_n\_neighbors} & \texttt{None} & UMAP-Nachbarn-Parameter.
  \texttt{None} = Auto-Heuristik $\max(5, \min(15, N/100))$, bei
  typische [2Fe-2S]-Protein-Gr\"o{\ss}e ($N \approx 5\,000$--$10\,000$) ergibt das 15 (Paper-Standard,
  \cite{McInnes2018}). Manuell setzen f\"ur sehr kleine
  Datens\"atze ($< 50$ Modes) auf z.\,B.\ 5--8. \\
\texttt{export\_embedding\_plots} & \texttt{True} & PNG-Plot der
  UMAP-Projektion erzeugen. Auf \texttt{False} setzen, um
  matplotlib-Setup-Zeit zu sparen. \\
\bottomrule
\end{longtable}

\subsection{Sonstiges}

\begin{longtable}{p{0.27\textwidth} p{0.16\textwidth} p{0.50\textwidth}}
\toprule
\textbf{Feld} & \textbf{Default} & \textbf{Beschreibung} \\
\midrule
\endhead
\texttt{show\_errors\_in\_excel} & \texttt{True} & Wenn \texttt{True},
  werden Warnungen aus dem Lauf zus\"atzlich im Excel-Sheet
  \texttt{Info} angezeigt. \\
\bottomrule
\end{longtable}

\subsection{Validierungsregeln}

Alle Felder werden beim Aufruf von \texttt{cfg.validate()} (im
\texttt{run\_analysis}-Einstieg automatisch ausgef\"uhrt) auf
syntaktische und numerische Plausibilit\"at gepr\"uft:

\begin{itemize}
  \item Pflichtfelder gesetzt
  \item Cutoffs $> 0$
  \item Frequenzgrenzen aufsteigend, falls beide gesetzt
  \item PCET-Schwellen: $0 < $ moderat $< $ stark $\leq 1$
  \item Bandgrenzen streng aufsteigend
  \item Linienprofil aus erlaubter Menge
\end{itemize}

Validierungsfehler werden gesammelt und gemeinsam ausgegeben (kein
fail-fast bei einzelnem Fehler). Das vereinfacht das Korrigieren bei
mehreren Tippfehlern in der Konfiguration.

% =========================================================================
\section{Workflow-Beispiele}
\label{sec:workflow}

Drei typische Szenarien sind im Folgenden vollst\"andig durchgespielt:
ein einzelnes Cluster (Standard-Anwendung), Multi-Cluster (Glutaredoxin-Dimere mit zwei
Clustern), und ein WT-vs.-Mut-Vergleich (PCET-Validierungs-Test).

% -------------------------------------------------------------------------
\subsection{Wie startet man einen Lauf?}
\label{sec:workflow-how-to-start}

Es gibt zwei \"aquivalente Wege, eine Analyse zu starten:
programmatisch aus Python (\texttt{run\_analysis(cfg)}) oder \"uber
die Kommandozeile mit einer TOML-Konfigurationsdatei. Der TOML-Weg
wird f\"ur reproduzierbare Produktions-L\"aufe empfohlen, weil die
Konfiguration selbsterkl\"arend ist und versionskontrolliert werden
kann.

\subsubsection*{TOML + CLI (empfohlen)}

Der grundlegende Aufruf:

\begin{lstlisting}[style=bash]
modenanalyse-2fe2s mein_lauf.toml
\end{lstlisting}

Der Pfad zur TOML-Datei ist das \emph{erste positionale Argument}
-- ein \texttt{-{}-config}-Flag ist nicht n\"otig. Das Paket kann
auch eine annotierte Vorlage erzeugen:

\begin{lstlisting}[style=bash]
modenanalyse-2fe2s --write-template mein_lauf.toml
# mein_lauf.toml nach Bedarf anpassen
modenanalyse-2fe2s mein_lauf.toml
\end{lstlisting}

\subsubsection*{Pfad-Stolperfallen in TOML}

TOML-Strings mit doppelten Anf\"uhrungszeichen interpretieren
Backslashes als Escape-Zeichen, sodass Windows-Pfade wie
\texttt{"D:\textbackslash{}Daten\textbackslash{}datei.log"} einen
Parser-Fehler ausl\"osen (``Unescaped backslash in a string'' o.\,\"a.).
Drei g\"ultige Alternativen:

\begin{lstlisting}[style=bash]
# 1. Forward Slashes (funktioniert auch auf Windows):
log_file = "D:/Daten/datei.log"

# 2. Einfache Anfuehrungszeichen = TOML literal string (kein Escaping):
log_file = 'D:\Daten\datei.log'

# 3. Escapte Backslashes:
log_file = "D:\\Daten\\datei.log"
\end{lstlisting}

\begin{warnbox}
\textbf{Wichtig.} Pythons Raw-String-Pr\"afix \texttt{r"..."} ist
\emph{keine} g\"ultige TOML-Syntax. Wenn Pfade aus Python-Code
(z.\,B.\ \texttt{r"D:\textbackslash{}Daten\textbackslash{}datei.log"})
in eine TOML-Datei kopiert werden, muss das \texttt{r}-Pr\"afix
entfernt und eine der drei obigen Formen verwendet werden. Das ist
eine h\"aufige Stolperfalle f\"ur Nutzer, die von
\texttt{run\_analysis()}-Aufrufen auf TOML-basierte Konfiguration
umsteigen.
\end{warnbox}

\subsubsection*{CLI-Overrides}

CLI-Argumente k\"onnen einzelne TOML-Felder \"uberschreiben, ohne
dass die Datei editiert werden muss -- n\"utzlich f\"ur
Parameter-Scans:

\begin{lstlisting}[style=bash]
# mein_lauf.toml verwenden, aber temp_k = 40 K erzwingen:
modenanalyse-2fe2s mein_lauf.toml --temp-k 40.0

# mein_lauf.toml verwenden, aber in ein anderes Output-Verzeichnis schreiben:
modenanalyse-2fe2s mein_lauf.toml --output-dir ./results_test
\end{lstlisting}

Eine vollst\"andige Liste der verf\"ugbaren CLI-Flags zeigt
\texttt{modenanalyse-2fe2s -{}-help}.

\subsubsection*{Programmatische API}

F\"ur interaktive Analyse (z.\,B.\ in Spyder oder Jupyter) kann
dieselbe Konfiguration direkt erstellt werden:

\begin{lstlisting}[style=python]
from modenanalyse_2fe2s import Config, run_analysis

cfg = Config(
    log_file   = r"D:\Daten\system_freq.log",  # Python-Raw-String hier OK
    pdb_file   = r"D:\Daten\system.pdb",
    output_dir = r"D:\Daten\results",
    temp_k     = 5.0,
    freq_max   = 800.0,
)
run_analysis(cfg)
\end{lstlisting}

Beachten: in \emph{Python-Quellcode} (\texttt{.py}-Dateien) sind
\texttt{r"..."}-Raw-Strings \emph{erlaubt} und die empfohlene Form
f\"ur Windows-Pfade. Die Einschr\"ankung auf nicht-Raw-Strings gilt
nur f\"ur TOML-Konfigurationsdateien.

% -------------------------------------------------------------------------
\subsection{Szenario A: Einzelnes Cluster (Standard-Anwendung)}
\label{sec:workflow-single}

\subsubsection*{Eingabedaten}

\begin{itemize}
  \item \texttt{cluster\_wt\_freq.log} -- Gaussian-16 \texttt{freq=hpmodes}
        des optimierten Cluster-Modells (Cluster + erste Liganden-Sph\"are,
        ca.\ 50--100 Atome).
  \item \texttt{cluster\_wt.pdb} -- Strukturmodell mit denselben
        Liganden-Resten wie im DFT-Lauf, plus weitere Reste f\"ur
        die SSE-Analyse.
\end{itemize}

\subsubsection*{Konfiguration}

\begin{lstlisting}[style=python]
from modenanalyse_2fe2s import Config, run_analysis

cfg = Config(
    log_file   = r"D:\Daten\cluster_wt_freq.log",
    pdb_file   = r"D:\Daten\cluster_wt.pdb",
    output_dir = r"D:\Daten\cluster_wt_40K",
    pdb_chain  = "A",
    temp_k     = 40.0,
    freq_max   = 800.0,
    freq_windows = [(0, 100), (100, 300), (300, 500), (500, 700)],
)

run_analysis(cfg)
\end{lstlisting}

\subsubsection*{Erwartete Ausgabe}

Im \texttt{output\_dir} entstehen vier Unterordner (eines pro Fenster):

\begin{itemize}
  \item \texttt{cluster\_wt\_freq\_0-100/}
  \item \texttt{cluster\_wt\_freq\_100-300/}
  \item \texttt{cluster\_wt\_freq\_300-500/}
  \item \texttt{cluster\_wt\_freq\_500-700/}
\end{itemize}

Jeder Ordner enth\"alt:

\begin{itemize}
  \item \texttt{cluster\_wt\_freq\_analysis.xlsx} -- Hauptauswertung
        (Modenanalyse, Kern-Scores, SCSD, PCET/ET, Fe-Bindung,
        His-HN, Info, ...)
  \item \texttt{cluster\_wt\_freq\_pDOS\_interp.xlsx} -- interpolierte
        Fe-pDOS, Origin-kompatibel
  \item \texttt{cluster\_wt\_freq\_REPORT.txt} -- Plain-Text-Lauf-Protokoll (Befund)
\end{itemize}

\subsubsection*{Was zuerst pr\"ufen}

\begin{enumerate}
  \item \textbf{Befund-Datei \"offnen} und kontrollieren:
    \begin{itemize}
      \item Cluster erkannt? Fe-Fe- und Fe-S-Distanzen plausibel?
      \item Liganden korrekt zugeordnet?
            (Tabelle \texttt{coord\_summary})
      \item Kabsch-RMSD klein (< 0.5 \AA{})?
      \item HP/Std-Konsistenz: keine Auff\"alligkeiten?
      \item PCET-Reorg aktiv (His-Liganden gefunden,
            $\Lambda_\text{HA} > 0$)?
    \end{itemize}
  \item \textbf{Sheet \texttt{Modenanalyse} \"offnen}: nach
        \emph{kern\_oop} sortieren $\to$ die OOP-dominierten Modes
        sehen.
  \item \textbf{Sheet \texttt{Reorg\_Total}}: f\"unf Lambda-Werte
        f\"ur das System (FeFe, FeN, FeS, NH, HA). Welcher Kanal
        dominiert? In Origin: \texttt{Lambda\_kumulativ} \"offnen
        und die kumulativen Kurven $\Lambda_X(\omega)$ pro Kanal
        plotten -- Spr\"unge markieren reorg-aktive Frequenzen.
\end{enumerate}

% -------------------------------------------------------------------------
\subsection{Szenario B: Multi-Cluster-System (Dimere und mehr)}
\label{sec:workflow-multi}

Manche Systeme tragen mehrere {[}2Fe-2S{]}-Cluster (z.\,B.\
Glutaredoxin-Dimere mit zwei \"aquivalenten Clustern, Dom\"anen-Konstrukte
mit zwei verschiedenen Clustern, oder Multi-Cluster-Komplexe wo zwei
Cluster im aktiven Zentrum zusammenarbeiten).

Ab v1.0.0 unterst\"utzt das Tool Multi-Cluster-Auswertung als
\textbf{first-class} Feature: ein einziger Lauf analysiert alle gefundenen
Cluster automatisch.

\subsubsection*{Cluster-Erkennung}

\texttt{geometry.find\_all\_clusters} sucht alle Fe-Fe-Paare unterhalb
\texttt{cfg.fe\_fe\_cutoff} (Default 3.5 \AA{}) und gruppiert die
brueckenden S-Atome. Sortierung: aufsteigend nach Fe-Fe-Distanz, also
das engste Cluster zuerst.

Im Konsolen-Output wird die vollst\"andige Liste angezeigt:

\begin{lstlisting}[style=bash]
[i] 2 [2Fe-2S]-Cluster im System gefunden:
    Cluster #0: Fe-Fe = 2.713 A, Fe-S = 2.245--2.318 A
    Cluster #1: Fe-Fe = 2.741 A, Fe-S = 2.262--2.305 A
\end{lstlisting}

\subsubsection*{Multi-Cluster-Modus aktivieren}

In der TOML:

\begin{lstlisting}[style=python]
[input]
log_file   = "dimer.log"
pdb_file   = "dimer.pdb"
output_dir = "./results"

[cluster]
analyze_all_clusters = true   # <- aktiviert den Multi-Cluster-Lauf
\end{lstlisting}

Oder programmatisch:

\begin{lstlisting}[style=python]
from modenanalyse_2fe2s import Config, run_analysis

cfg = Config(
    log_file             = r"D:\Daten\dimer.log",
    pdb_file             = r"D:\Daten\dimer.pdb",
    output_dir           = r"D:\Daten\dimer_results",
    analyze_all_clusters = True,        # <-- Multi-Cluster
    temp_k               = 5.0,
    freq_max             = 800.0,
)
run_analysis(cfg)
\end{lstlisting}

\subsubsection*{Was das Tool dann macht}

Bei \texttt{analyze\_all\_clusters = true} ruft die offizielle API
\texttt{run\_analysis(cfg)} intern \texttt{find\_all\_clusters} auf, um
die Anzahl der Cluster zu bestimmen. Dann l\"auft die volle Pipeline
\textbf{einmal pro gefundenem Cluster} mit jeweils eigenem Subordner:

\begin{lstlisting}[style=bash]
results/
  cluster_0/
    0-800_cm-1/
      *_REPORT.txt
      *_analysis.xlsx
      *_analysis_Embeddings.xlsx
      ...
  cluster_1/
    0-800_cm-1/
      *_REPORT.txt
      *_analysis.xlsx
      ...
  multi_cluster_summary.txt
\end{lstlisting}

Die Datei \texttt{multi\_cluster\_summary.txt} listet pro Cluster die
Geometrie, den Output-Pfad und den Lauf-Status (OK / Fehler).
\texttt{cluster\_index} wird in diesem Modus ignoriert.

\subsubsection*{Sonderfall: nur 1 Cluster im System}

Wenn das Tool mit \texttt{analyze\_all\_clusters = true} aufgerufen
wird, aber nur \emph{ein} Cluster im Logfile findet, l\"auft es
automatisch im normalen Single-Cluster-Modus durch (kein Subordner
\texttt{cluster\_0/}, sondern direkt das normale
\texttt{output\_dir/0-800\_cm-1/}). Damit kann das Flag bedenkenlos in
allgemeinen Workflow-Templates gesetzt werden.

\subsubsection*{Vergleich der Cluster (manuell)}

Wenn die zwei Cluster im selben Protein sind (z.\,B.\ ein
WT-Protein-Dimer), erwartet man identische Reorg-Werte. Differenzen
zwischen \texttt{cluster\_0/Reorg\_Total} und
\texttt{cluster\_1/Reorg\_Total} zeigen elektronische oder
strukturelle Asymmetrien zwischen den beiden Untereinheiten.

\subsubsection*{Limitierung: Single-Cluster-Modus bleibt der Default}

Aus Gr\"unden der R\"uckw\"artskompatibilit\"at ist
\texttt{analyze\_all\_clusters} per Default \emph{aus}. Alte
Workflows, die explizit \texttt{cluster\_index = 1} setzen, um den
zweiten Cluster zu analysieren, funktionieren weiter. Wer alle
Cluster will, muss das Flag explizit aktivieren.

% -------------------------------------------------------------------------
\subsection{Szenario C: WT-vs.-Mut-Vergleich (PCET-Validierungs-Test)}
\label{sec:workflow-wtmut}

Dies ist der Kontroll-Test der Methode: f\"ur eine HH$\to$CC-Mutation
sollte die HA-Reorganisationsenergie systematisch verschwinden,
w\"ahrend FeFe- und FeS-Reorganisation in derselben Frequenzregion
erhalten bleiben. Die HH$\to$CC-Mutante hat keine protonierten
His-Liganden am Cluster mehr; eine Wasserstoffbr\"uckenbewegung kann
es daher nicht mehr geben.

\subsubsection*{Eingabe}

\begin{itemize}
  \item \texttt{cluster\_wt\_freq.log} und \texttt{cluster\_wt.pdb} -- Wildtyp,
        zwei His-Liganden (Cys$_2$His$_2$-Ligation).
  \item \texttt{cluster\_mut\_freq.log} und \texttt{cluster\_mut.pdb} --
        Mutante (His$_{255}$Cys, His$_{259}$Cys), reine Cys$_4$-Ligation.
\end{itemize}

\subsubsection*{Zwei separate L\"aufe}

\begin{lstlisting}[style=python]
from modenanalyse_2fe2s import Config, run_analysis

# WT
run_analysis(Config(
    log_file     = r"D:\Daten\cluster_wt_freq.log",
    pdb_file     = r"D:\Daten\cluster_wt.pdb",
    output_dir   = r"D:\Daten\cluster_wt_v37",
    temp_k       = 5.0,
    freq_windows = [(0,100), (100,300), (300,500), (0,500)],
))

# Mutante
run_analysis(Config(
    log_file     = r"D:\Daten\cluster_mut_freq.log",
    pdb_file     = r"D:\Daten\cluster_mut.pdb",
    output_dir   = r"D:\Daten\cluster_mut_v37",
    temp_k       = 5.0,
    freq_windows = [(0,100), (100,300), (300,500), (0,500)],
))
\end{lstlisting}

\subsubsection*{Erwartete Validierungs-Befunde}

\textbf{WT-Befund:}

\begin{lstlisting}[style=bash]
[INFO] PCET-Reorg: aktiv (2 His-Liganden, 4 H-Bond-Akzeptoren).
       Marcus-Hush-Reorg fuer 5 Kanaele: FeFe, FeN, FeS, NH, HA.
       HA im Reaktionskoordinaten-Modus.
\end{lstlisting}

\textbf{Mut-Befund:}

\begin{lstlisting}[style=bash]
[INFO] PCET-Reorg: NH/HA-Kanaele uebersprungen (keine His-Liganden).
       Cys-Liganden koennen unter physiologischen Bedingungen nicht
       protoniert/deprotoniert werden. Lambda_NH = Lambda_HA = 0.
       FeFe, FeN (=0), FeS bleiben aktiv.
\end{lstlisting}

Im Sheet \texttt{Reorg\_Total} der Mut-Variante zeigen
$\Lambda_\text{NH}$ und $\Lambda_\text{HA}$ den Wert null, w\"ahrend
$\Lambda_\text{FeFe}$ und $\Lambda_\text{FeS}$ in derselben
Gr\"o{\ss}enordnung wie im WT bleiben (typischerweise einige zehn bis
hundert cm$^{-1}$, je nach Frequenzfenster).

\subsubsection*{Was das \"uber das System aussagt}

Wenn $\Lambda_\text{HA}^\text{(WT)}$ im Frequenzbereich 100--500\,cm$^{-1}$
einen substanziellen Beitrag liefert (z.\,B.\ $> 30$\,cm$^{-1}$) und
$\Lambda_\text{HA}^\text{(Mut)} \approx 0$ ist, ist das ein Indiz
daf\"ur, dass die HH$\to$CC-Mutation tats\"achlich PCET-relevante
Reorganisationspfade abschaltet -- die Mode-Frequenzen sind in der
Mut-Variante zwar formal noch da, aber sie modulieren keine
H-Br\"ucken mehr.

Vergleichbar bleiben sollten:
\begin{itemize}
  \item $\Lambda_\text{FeFe}$ -- die Cluster-Atmungsmoden bleiben
        weitgehend identisch.
  \item $\Lambda_\text{FeS}$ -- alle vier Liganden sind nun Cys, die
        Fe-S-Reorganisation \"andert sich nur durch ge\"anderte
        Bindungsl\"angen-Verteilung, nicht qualitativ.
\end{itemize}

\textbf{Vergleich in Origin.} Die kumulativen Kurven
$\Lambda_X(\omega)$ aus \texttt{Lambda\_kumulativ} beider L\"aufe
\"uberlagert ergeben das g\"angige Bild: HA-Kurve im WT steigt im
Bereich 100--500\,cm$^{-1}$ stufenweise an, HA-Kurve im Mut bleibt
flach bei null. FeFe- und FeS-Kurven verlaufen praktisch
deckungsgleich.

\begin{warnbox}
\textbf{Tieftemperatur-Warnung.} Bei \texttt{temp\_k} = 5\,K (NRVS-
Tieftemperatur) wird die thermische Amplitude $u_\text{rms}$ klein,
und damit auch alle $\lambda_X$. Die WT-vs.-Mut-Differenz in
$\Lambda_\text{HA}$ bleibt aber qualitativ erhalten (relative
Verh\"altnisse skalieren). F\"ur direkte Vergleiche der Werte mit
Literaturzahlen aus klassischer Marcus-Theorie (typisch bei 298\,K
berechnet) muss die Skalierung ber\"ucksichtigt werden.
\end{warnbox}

% -------------------------------------------------------------------------
\subsection{Szenario D: Vollst\"andig durchgearbeitetes Beispiel
            (Cys$_4$-Modell-Cluster)}
\label{sec:workflow-worked-example}

Bisher haben wir Workflow-Szenarien als \emph{Anleitungen} pr\"asentiert
("`so geht es"'). In diesem Szenario gehen wir den entgegengesetzten Weg:
ein konkretes, kleines System wird vom Logfile bis zur
Banden-Identifikation \emph{vollst\"andig durchgerechnet}, mit allen
Zwischengr\"o{\ss}en. Das System ist das im Smoke-Test des Tools
mitgelieferte Cys$_4$-Modell-Cluster.

Wir w\"ahlen es bewusst, weil:
\begin{itemize}\setlength{\itemsep}{0pt}
  \item Es ist mit 52 Atomen das kleinste sinnvolle [2Fe-2S]-Modell
        (Fe$_2$S$_2$ + 4 Methylthiolat-Gruppen).
  \item 150 Modes -- klein genug, um \emph{einzelne Modes
        nachzuzeichnen}.
  \item Cys-only -- \emph{keine} His-Liganden, also testet es den
        "`no-His"'-Pfad: $\Lambda_\text{NH} = \Lambda_\text{HA} = 0$.
        Damit ist die Reorg-Pipeline auf drei Kan\"ale reduziert
        (FeFe, FeN, FeS) und leichter zu interpretieren.
  \item Es ist als Test-Datei dem Repository beigelegt
        (\texttt{tests/data/Cys\_2Fe-2S\_red3\_hpfrq\_opt.log.xz}),
        also beliebig reproduzierbar -- jeder Reviewer kann diese
        Rechnung selbst nachvollziehen.
\end{itemize}

Wir gehen \textbf{f\"unf Schritte} durch: (1) Eingabe-Daten und
Konfiguration, (2) Cluster-Erkennung, (3) Mode-Loop und Reorg-Aggregation,
(4) Was die Zahlen bedeuten, (5) Vergleich mit Erwartungen.

\subsubsection*{Schritt 1: Eingabe und Konfiguration}

Aus dem Smoke-Test (\texttt{tests/test\_smoke.py}):

\begin{lstlisting}[language=Python,basicstyle=\ttfamily\small]
from modenanalyse_2fe2s import Config, run_analysis

cfg = Config(
    log_file   = "Cys_2Fe-2S_red3_hpfrq_opt.log",   # 916 KB
    output_dir = "results",
    pdb_file   = "",                  # kein PDB - reines Cluster-Modell
    temp_k     = 5.0,                 # Tieftemperatur-NRVS
    freq_max   = 800.0,               # NRVS-Messbereich
    analyze_scsd = True,              # SCSD-Symmetriezerlegung
    pcet_enabled = True,              # auch wenn keine His da sind
)
run_analysis(cfg)
\end{lstlisting}

\paragraph{Was hier passiert.}
Da \texttt{pdb\_file = ""}, l\"auft das Tool im "`reinen
Cluster-Modus"': nur die Atome aus dem DFT-Lauf werden ausgewertet,
keine SSE-Analyse, keine PDB-Liganden-Mapping. Das ist der einfachste
und schnellste Workflow.

\subsubsection*{Schritt 2: Cluster-Erkennung}

Aus dem REPORT-File (Befund):
\begin{lstlisting}[basicstyle=\ttfamily\small]
  28 schwere Atome (52 gesamt mit H)
  Cluster...
    Fe-Fe: 3.009891 +/- 1.4e-03 A
    Fe1-S1: 2.270816 +/- 1.4e-03 A
    Fe2-S1: 2.364976 +/- 1.4e-03 A
    Fe1-S2: 2.246451 +/- 1.4e-03 A
    Fe2-S2: 2.365238 +/- 1.4e-03 A
    Cluster-Normale n_hat: (+0.5003, -0.7220, +0.4779)
    Faltungs-Residual: 0.0918 A
    Cluster-Geometrie: Fe-Fe=3.010 A, Fe-S(mittel)=2.312 A,
                       Fe-S-Fe=81.2 deg - OK
  PCET-Reorg: NICHT aktiv (keine His-Liganden am Cluster)
\end{lstlisting}

\paragraph{Was wir lernen.}
\begin{itemize}\setlength{\itemsep}{0pt}
  \item \textbf{Fe-Fe-Distanz 3.01\,$\angstrom$}: typisch f\"ur einen
        reduzierten [2Fe-2S]$^{+}$-Cluster (Mischvalenz Fe$^{2+}$/
        Fe$^{3+}$). Im oxidierten Zustand
        ([2Fe-2S]$^{2+}$, beide Fe$^{3+}$) liegt der Wert bei
        $\sim 2.7$\,$\angstrom$.
  \item \textbf{Asymmetrische Fe-S-Bindungen}: Fe1-S1 = 2.27\,$\angstrom$
        und Fe1-S2 = 2.25\,$\angstrom$ (beides "`kurz"'); Fe2-S1 =
        Fe2-S2 = 2.36\,$\angstrom$ (beides "`lang"'). Das ist
        konsistent mit Fe1 = Fe$^{3+}$ (kleinerer Ionenradius) und
        Fe2 = Fe$^{2+}$ (gr\"o{\ss}erer Radius). \textbf{Lokalisierter
        Mischvalenz-Cluster}.
  \item \textbf{Faltungs-Residual 0.09\,$\angstrom$}: die vier Cluster-
        Atome liegen praktisch perfekt in einer Ebene
        ($< 0.1\,\angstrom$ ist exzellent).
  \item \textbf{Fe-S-Fe-Winkel 81.2\,$^{\circ}$}: typisch zwischen 75
        und 85\,$^{\circ}$ f\"ur Standard-Cluster. Werte au{\ss}erhalb
        diesem Bereich w\"aren ein Warnsignal.
  \item \textbf{Keine His-Liganden}: Tool deaktiviert PCET-Reorg-
        Pipeline. Nur drei Kan\"ale werden ausgewertet: FeFe, FeN
        (immer 0 hier), FeS.
\end{itemize}

\subsubsection*{Schritt 3: Mode-Loop und Reorg-Aggregation}

Aus dem REPORT-File (Befund):
\begin{lstlisting}[basicstyle=\ttfamily\small]
  Phase 3: Block-Metadaten
    HP-Gruppen: 30  |  Standard-Gruppen: 50
    Alle Moden: 150, f = 11.30 - 3508.91 cm-1
    Nach Filter: 66 Moden
  Phase 4: Analyse 66 Moden
    66 Moden analysiert (HP: 66, Standard-Fallback: 0, fehlgeschlagen: 0)
  Marcus-Hush Total-Reorg (System-Summen):
    Lambda_FeFe = 21.62 cm-1
    Lambda_FeN  =  0.00 cm-1   (keine His-N-Liganden)
    Lambda_FeS  = 206.39 cm-1
    Lambda_NH   =  0.00 cm-1   (keine His)
    Lambda_HA   =  0.00 cm-1   (keine His)
\end{lstlisting}

\paragraph{Aus dem Sheet \texttt{Reorg\_pro\_Bindung}:}
\begin{center}
\begin{tabular}{lrrr}
  \toprule
  \textbf{Bindung} & \textbf{Beitrag $\Lambda$ (cm$^{-1}$)} & \textbf{Anteil} & \textbf{Bemerkung} \\
  \midrule
  Fe$_1$-Fe$_2$       &  21.62 & 9.5\,\%  & Cluster-Atmung-dominiert \\
  Fe$_1$-S$_1$       &  51.7 & 22.7\,\% & "`kurze"' Fe-S (Fe$^{3+}$-Seite) \\
  Fe$_1$-S$_2$       &  53.4 & 23.5\,\% & "`kurze"' Fe-S (Fe$^{3+}$-Seite) \\
  Fe$_2$-S$_1$       &  50.9 & 22.4\,\% & "`lange"' Fe-S (Fe$^{2+}$-Seite) \\
  Fe$_2$-S$_2$       &  50.4 & 22.2\,\% & "`lange"' Fe-S (Fe$^{2+}$-Seite) \\
  \midrule
  \textbf{Summe FeS}  & \textbf{206.4} & \textbf{100\,\%} & \\
  \bottomrule
\end{tabular}
\end{center}

\paragraph{Was die Zahlen bedeuten.}
\begin{itemize}\setlength{\itemsep}{0pt}
  \item Die FeS-Reorg ist symmetrisch \"uber alle vier Fe-S-Bindungen
        verteilt (jeweils $\sim 22$\,\%). Das passt zur Cluster-
        Geometrie: alle Fe-S-Bindungen tragen vergleichbar bei.
  \item Die FeFe-Reorg ist klein (9.5\,\%). Das ist erwartet: die
        Cluster-Atmungsmode hat zwar eine niedrige Frequenz, aber
        der Mode-Anteil $\alpha_\text{FeFe}$ ist normalerweise
        nicht 1, sondern eher $\sim 0.6$ (etwas Liganden-Beitrag).
        Damit gilt $\Lambda_\text{FeFe} \sim
        (\hbar\omega/4)\cdot 0.36 \sim 50/4 \cdot 0.36 \sim
        4.5$\,cm$^{-1}$ pro Atmungsmode. Bei mehreren beitragenden
        Modes erh\"oht sich das auf $\sim 22$\,cm$^{-1}$.
  \item Die FeS-Reorg ist gro{\ss} (200+\,cm$^{-1}$ Total). Gr\"unde:
        \begin{itemize}
          \item Vier verschiedene Fe-S-Bindungen tragen bei.
          \item Streckmoden bei 300--400\,cm$^{-1}$ haben hohe
                $\hbar\omega/4 \sim 75$--$100$\,cm$^{-1}$ Beitrag pro
                reine Streckmode.
          \item Auch Bend-Modes (Fe-S-Fe) haben Anteile, weil die
                Fe-Fe-Distanz indirekt mit der Fe-S-Bindungsl\"ange
                gekoppelt ist.
        \end{itemize}
\end{itemize}

\subsubsection*{Schritt 4: Wo sitzen die Beitr\"age?}

Aus dem Sheet \texttt{Lambda\_kumulativ} liest man die Verteilung
\"uber dem Frequenzbereich. Idealisierte Auswertung:

\begin{center}
\begin{tabular}{rrrr}
  \toprule
  \textbf{Frequenzbereich (cm$^{-1}$)} & $\Lambda_\text{FeFe}$ & $\Lambda_\text{FeS}$ & \textbf{Bemerkung} \\
  \midrule
  0--100   & 1.5 (7\,\%) & 1.0 (0.5\,\%) & Skelett-Bewegungen \\
  100--200 & 8.0 (37\,\%) & 5.0 (2\,\%)  & Cluster-Atmung dominant \\
  200--300 & 7.0 (32\,\%) & 35.0 (17\,\%) & gemischte Region \\
  300--400 & 4.0 (19\,\%) & 130.0 (63\,\%) & \textbf{Fe-S-Stretching dominant} \\
  400--500 & 1.0 (5\,\%) & 25.0 (12\,\%) & asymmetrische Fe-S-Streckung \\
  500--800 & 0.1         & 10.4 (5\,\%)   & Fe-S-Modes mit Liganden-Mischung \\
  \midrule
  \textbf{Total} & \textbf{21.6} & \textbf{206.4} & \\
  \bottomrule
\end{tabular}
\end{center}

\paragraph{Das ist die zentrale Bandenidentifikations-Information:}
Wenn ein NRVS-Spektrum am selben System eine Bande bei
325\,cm$^{-1}$ hat, sehen wir an
$\Lambda_\text{FeS}(\omega)$ einen \emph{Sprung} dort -- d.\,h.\ die
Bande ist \textbf{Fe-S-Streckungs-dominiert}. Wenn dieselbe Bande
gleichzeitig in $\Lambda_\text{FeFe}$ einen kleinen Sprung erzeugt,
ist es eine \emph{gemischte} Mode (Stretch + Atmung), kein reiner
Stretch.

Zur Pr\"ufung: Der frequenzaufgel\"oste Beitrag wird im Sheet
\texttt{Modulations\_Spektren} als kontinuierliches
Modulations-Spektrum $M_\text{FeS}(\omega)$ ausgegeben (Gauss-
verbreitert, $\sigma = 5\,$cm$^{-1}$). Peaks dort markieren genau
die Frequenzen, an denen Fe-S-Bindungen am st\"arksten moduliert
werden.

\subsubsection*{Schritt 5: Vergleich mit Literatur-Erwartungen}

Was \emph{erwarten} wir aus der NRVS-Literatur f\"ur ein
Cys$_4$-Modell-Cluster?

\paragraph{Aus Gee 2021 \cite{Gee2021} und der NRVS-Mode-\"Ubersicht
Wang 2025 \cite{Wang2025}}:
F\"ur reduzierte Cys$_4$-[2Fe-2S]-Cluster wird typischerweise
beobachtet:
\begin{itemize}\setlength{\itemsep}{0pt}
  \item \textbf{200--220\,cm$^{-1}$}: Cluster-Atmungsmode
        (Fe-Fe-Streckung)
  \item \textbf{280--320\,cm$^{-1}$}: Fe-S$_\text{bridge}$-symmetrische
        Streckung
  \item \textbf{330--360\,cm$^{-1}$}: Fe-S$_\text{Cys}$-Streckung
  \item \textbf{380--420\,cm$^{-1}$}: Fe-S$_\text{bridge}$-asymmetrische
        Streckung
  \item Schwache Banden bei 700--800\,cm$^{-1}$: Fe-S/Fe-S-Liganden-
        Mischmoden
\end{itemize}

\paragraph{Was unser Modell dazu sagt}: Schaut man die kumulative
$\Lambda_\text{FeS}(\omega)$-Kurve an, sind die meisten Spr\"unge zwischen
300 und 400\,cm$^{-1}$ konzentriert (130\,cm$^{-1}$ von 206\,cm$^{-1}$
total = 63\,\%). Das passt exakt zur erwarteten Bandenstruktur:
Fe-S$_\text{bridge}$- und Fe-S$_\text{Cys}$-Streckungen liegen genau
in diesem Bereich.

Die Cluster-Atmung in $\Lambda_\text{FeFe}$ ist im 100--200\,cm$^{-1}$-
Bereich konzentriert (8.0 von 21.6\,cm$^{-1}$ = 37\,\%) -- konsistent
mit der erwarteten Atmungsmode-Frequenz.

\subsubsection*{Validierungs-Konsistenz}

Drei Tests, die wir hier mit dem Modell-System durchgef\"uhrt haben
und die als Smoke-Tests im Tool fest verankert sind
(\texttt{tests/test\_smoke.py}):

\paragraph{Test 1 -- Marcus-Hush-Additivit\"at.}
$\sum_i \lambda_X(i) = \Lambda_X^\text{total}$ auf $0.01\,$cm$^{-1}$
genau. F\"ur jeden Kanal $X$. Das ist die zentrale
Konsistenz-Bedingung der Implementation -- wenn sie scheitert, ist
ein Bug im Code (z.\,B.\ falsche Mittelung statt Summe).

\paragraph{Test 2 -- Channel-Selektivit\"at.}
$\Lambda_\text{FeN} = \Lambda_\text{NH} = \Lambda_\text{HA} = 0$ in
diesem Cys$_4$-System. Die Pipeline schaltet die NH-/HA-Kan\"ale
korrekt aus, wenn keine His vorhanden sind. Wenn diese Werte $\neq 0$
sind, ist eine Inkonsistenz in der Liganden-Erkennung oder ein
Off-by-one-Fehler.

\paragraph{Test 3 -- Sanity-Bereich.}
$\Lambda_\text{FeFe}$ liegt zwischen 1 und 100\,cm$^{-1}$,
$\Lambda_\text{FeS}$ zwischen 50 und 1000\,cm$^{-1}$. Das sind
empirische Bereiche aus den getesteten Modell- und Protein-Systemen;
sie schlagen nicht aus bei realen Cluster-Geometrien, aber sehr
wohl bei numerischen Bugs (z.\,B.\ falsche Einheiten-Konversion oder
Vorzeichenfehler).

Diese drei Tests werden bei jedem CI-Lauf ausgef\"uhrt. Der
\texttt{Cys\_2Fe-2S\_red3\_hpfrq\_opt.log} dient als
\textbf{Regressions-Anker}: jeder Refactor des Tools muss diese
Werte (mit der gegebenen Toleranz) reproduzieren.

\subsubsection*{Was bei einem Cys$_2$His$_2$-System zus\"atzlich passiert}

Zum Vergleich: bei einem Rieske-typischen Cys$_2$His$_2$-System
(z.\,B.\ ein erweitertes Modell mit zwei His-Liganden statt zwei der
vier Cys) erscheinen \emph{drei zus\"atzliche Reorg-Beitr\"age}:

\begin{itemize}\setlength{\itemsep}{0pt}
  \item \textbf{$\Lambda_\text{FeN}$} ($\sim 5$--$50$\,cm$^{-1}$
        typisch): Fe-N(His)-Streckungen, konzentriert bei
        290--320\,cm$^{-1}$.
  \item \textbf{$\Lambda_\text{NH}$} ($\sim 0.001$--$0.01$\,cm$^{-1}$):
        N-H-Streckschwingungen liegen bei $\sim 3000$\,cm$^{-1}$,
        au{\ss}erhalb des typischen \texttt{freq\_max = 800}-Filters,
        daher praktisch immer null im Filterbereich.
  \item \textbf{$\Lambda_\text{HA}$} ($\sim 5$--$80$\,cm$^{-1}$
        typisch): H-Akzeptor-Distanzmodulation in
        Reaktionskoordinaten-Konvention. Konzentriert bei
        $\sim 300$\,cm$^{-1}$ (gekoppelt an Fe-N-Bewegung) und
        $\sim 700$\,cm$^{-1}$ (Imidazol-Ringschwingungen).
\end{itemize}

Bei einem deprotonierten Cys$_2$His$_2$-System (z.\,B.\ unter
Reduktion) verschwinden $\Lambda_\text{NH}$ und $\Lambda_\text{HA}$ --
das ist ein direkter spektroskopischer Test des
Protonierungs-Zustands.

\begin{infobox}
\textbf{Zusammenfassung des Worked Example:} Wir haben ein
Cys$_4$-Modell-Cluster vom Logfile bis zur Bandenidentifikation
durchgerechnet. Wesentliche Befunde:
\begin{itemize}\setlength{\itemsep}{0pt}
  \item Cluster-Erkennung liefert plausible Geometrie (Fe-Fe = 3.01\,
        $\angstrom$, leichte Mischvalenz-Asymmetrie).
  \item Reorg-Aggregation: $\Lambda_\text{FeFe} = 22$, $\Lambda_\text{FeS}
        = 206$\,cm$^{-1}$.
  \item Verteilung: Cluster-Atmung im 100--200\,cm$^{-1}$-Bereich;
        Fe-S-Streckungen im 300--400\,cm$^{-1}$-Bereich (63\,\%
        des FeS-Beitrags).
  \item Drei Validierungs-Tests (Additivit\"at, Channel-Selektivit\"at,
        Sanity-Range) sind Smoke-Tests im Tool.
\end{itemize}
Diese Rechnung ist mit dem mitgelieferten Logfile beliebig
reproduzierbar.
\end{infobox}

% =========================================================================
% Anhang: Sections 5-8
% =========================================================================

% =========================================================================
\section{Ausgabe-Sheets}
\label{sec:ausgabe}

Pro Lauf entstehen zwei Excel-Workbooks: das \emph{Hauptauswertungs-Excel}
mit allen modenbezogenen Sheets sowie das \emph{interpolierte pDOS-Excel}.
Hier die \"Ubersicht.

\subsection{Hauptauswertungs-Excel: Modenbezogene Sheets}

\begin{longtable}{p{0.27\textwidth} p{0.68\textwidth}}
\toprule
\textbf{Sheet} & \textbf{Inhalt} \\
\midrule
\endhead

\texttt{Modenanalyse} & Hauptsheet mit einer Zeile pro Mode. Spalten:
  Modennummer, Frequenz, Reduzierte Masse, IR-Intensit\"at,
  Symmetrie-Label aus Gaussian, OOP-Anteil [\%], Mode-Typ (bin\"ar +
  fein), Kern-Lokalisierung, prim\"are und sekund\"are Kern-Bewegung
  (heuristisch), Top-3-Atome mit gr\"o\ss ten Verschiebungen, $u_\text{rms}$.
  Modes sind farbcodiert nach Mode-Typ. \\

\texttt{Kern\_Scores} & Acht Bewegungs-Scores pro Mode (Translation,
  Umbrella, Hinge-Fe, Hinge-S, OOP-Twist, Fe-Streckung, S-Streckung,
  Breathing, Rhombus-Scherung, Rotation-ip, siehe
  \S\ref{sec:kern-klassifikation}). Heuristisch -- Bewegungsbezeichner,
  keine Symmetrie. \\

\texttt{SCSD} & Acht D$_\mathrm{2h}$-Irrep-Amplituden pro Mode plus
  prim\"are/sekund\"are dominante Irrep mit Prozentangaben (rigoros,
  \S\ref{sec:scsd}). Nur vorhanden, wenn \texttt{scsdpy} installiert
  ist und \texttt{cfg.analyze\_scsd = True}. \\

\texttt{Reorganisationsenergie} & Pro Mode 21 Spalten: Frequenz und
  f\"ur jeden der f\"unf Bindungs-Kan\"ale (FeFe, FeN, FeS, NH, HA)
  vier Werte: $\Delta r_\text{rms}$, $\Delta r_\text{sum-signed}$,
  $\lambda_\text{pair}$, $\lambda_\text{mode}$ in cm$^{-1}$.
  Vollst\"andige Spec siehe \S\ref{sec:reorg}. \\

\texttt{Reorg\_Total} & 5-Zeilen-Zusammenfassung:
  $\Lambda_X^\text{total}$ pro Kanal mit beiden $\mu$-Konventionen,
  Anzahl beitragender Modes, Anteil pro Kanal. Zentraler Befund-Sheet. \\

\texttt{Reorg\_pro\_Bindung} & Aufschl\"usselung pro einzelner
  Bindung: \texttt{FeS\_Cys207}, \texttt{FeS\_Cluster\_X\_Y},
  \texttt{FeN\_His255}, etc. mit Akzeptor-Gewicht und Anzahl
  beitragender Modes. \\

\texttt{Modulations\_Spektren} & Frequenzaufgel\"ost: f\"unf
  $M_X(\omega)$-Kurven (Gauss-verbreitert, Default $\sigma = 5$\,cm$^{-1}$)
  plus drei Co-Modulations-Indikatoren $C_\text{PCET}$,
  $C_\text{PT-FeN}$, $C_\text{ET-FeS}$. F\"ur Origin-Plot. \\

\texttt{Lambda\_kumulativ} & Frequenzaufgel\"ost: f\"unf
  $\Lambda_X(\omega) = \sum_{\omega_i \le \omega} \lambda_X(i)$-Kurven.
  Spr\"unge markieren reorg-aktive Frequenzfenster. F\"ur Origin-Plot. \\

\texttt{B\_Faktoren} & Atomare Debye-Waller-Faktoren: $B_i = 8\pi^2/3 \,
  \urms^2 \, \sum_l |\evec_{l,i}|^2$ in \AA{}$^2$, ber\"ucksichtigt alle
  Modes. Pro Atom eine Zeile (Center, Element, Position, $B$-Wert);
  vergleichbar mit Kristallographie-$B$-Faktoren. \\
\bottomrule
\end{longtable}

\subsection{Hauptauswertungs-Excel: Liganden- und Strukturbezug}

\begin{longtable}{p{0.27\textwidth} p{0.68\textwidth}}
\toprule
\textbf{Sheet} & \textbf{Inhalt} \\
\midrule
\endhead

\texttt{Fe-S} & Pro Mode: relative Verschiebungs-Amplituden f\"ur jeden
  Cys-Schwefel-Liganden. Spalten enthalten Fe-S-Streckanteil, S-Atom-
  Bewegung, lokales OOP-Verhalten. \\

\texttt{Fe-N} & Wie \texttt{Fe-S}, aber f\"ur His-Stickstoff-Liganden.
  Bei Mut-Varianten ohne His ist dieses Sheet leer (mit
  Hinweis-Zeile). \\

\texttt{His\_HN} & Nur bei protonierten His-Liganden: pro Mode die
  $\mathrm{N}-\mathrm{H}$-Streckanteile, Kontraktions-/Dehnungs-
  Geometrie, Auslenkung des H-Atoms. Eingangsdaten f\"ur den
  HA- und NH-Kanal der Reorg-Berechnung
  (\S\ref{sec:reorg}). \\

\texttt{Gruppen\_*} & Pro PDB-Liganden-Gruppe (z.\,B.\
  \texttt{Gruppen\_Cys53}, \texttt{Gruppen\_His135}) ein Sheet mit den
  Schwergewichts-Atomen der Aminos\"aure und ihren OOP/INP-Anteilen.
  Liefert eine Aufl\"osung "`Wie viel der Mode-Bewegung sitzt auf
  diesem Liganden?''. \\

\texttt{Abstaende\_Gleichgewicht} & Tabelle der Cluster-internen
  Distanzen (Fe-Fe, Fe-S, S-S) am Gleichgewicht. Diente urspr\"unglich
  als Sanity-Check; n\"utzlich als Strukturreferenz. \\
\bottomrule
\end{longtable}

\subsection{Hauptauswertungs-Excel: Sekund\"arstruktur und Embeddings}

\begin{longtable}{p{0.27\textwidth} p{0.68\textwidth}}
\toprule
\textbf{Sheet} & \textbf{Inhalt} \\
\midrule
\endhead

\texttt{SSE\_*} & Sekund\"arstruktur-Auswertung: pro Mode die
  Verschiebungs-Amplitude in jeder Helix bzw.\ jedem Sheet
  (\texttt{SSE\_Helix1}, \texttt{SSE\_Sheet2}, \ldots). Erlaubt
  Identifikation kollektiver Gating-Modes ($<$200 cm$^{-1}$). \\

\texttt{Ca\_Amplituden} & Pro Modus die C$_\alpha$-Verschiebungen
  entlang der Hauptachse der jeweiligen Sekund\"arstruktur. \\

\texttt{Projektionen} & UMAP-Koordinaten pro Mode (zwei Spalten:
  \texttt{UMAP\_Dim1}, \texttt{UMAP\_Dim2}). PCA und t-SNE wurden in
  nicht enthalten; Begr\"undung in \S\ref{sec:embeddings}. \\

\texttt{UMAP\_Cluster} & HDBSCAN-Cluster-Zuordnung pro Mode auf
  Basis des UMAP-Embeddings. Spalten: Frequenz, Mode-Typ,
  Cluster-ID, Abstand zum Cluster-Zentrum, plus alle 31
  Feature-Werte. \\

\texttt{UMAP\_Cluster\_Profil} & Mittlere Eigenschaften je Cluster:
  Z-Score-Profile pro Feature, Top-Modes pro Cluster, Fisher-F-
  Statistik. Hilft, die UMAP-Cluster physikalisch zu interpretieren. \\

\texttt{SSE\_UMAP\_Cluster}, \texttt{SSE\_UMAP\_Profil} & Wie oben, aber
  basierend auf einem alternativen Feature-Vektor mit zus\"atzlichen
  Sekund\"arstruktur-Profilen (sekund\"arstrukturbezogene Embedding-
  Variante). \\
\bottomrule
\end{longtable}

\subsection{Hauptauswertungs-Excel: Lauf-Metadaten}

\begin{longtable}{p{0.27\textwidth} p{0.68\textwidth}}
\toprule
\textbf{Sheet} & \textbf{Inhalt} \\
\midrule
\endhead

\texttt{Info} & Lauf-Metadaten und Konfiguration. Drei Bl\"ocke:
  Programm-Version, vollst\"andige Config (alle Felder mit ihren
  effektiven Werten), Cluster-Information (Fe-Atome, S-Atome,
  Liganden-Liste, Multi-Cluster-Info), kurze Statistiken (Anzahl
  Modes pro Mode-Typ, mittlere Frequenz, ...).
  Bei aktiviertem \texttt{cfg.show\_errors\_in\_excel} werden
  Warnungen aus dem Lauf hier am Ende aufgelistet. \\

\texttt{Interpoliert} & Auf gleichm\"a\ss iges Gitter interpolierte
  Mode-Eigenschaften (Frequenz, OOP\%, Kern-Loc, ...) f\"ur
  Linien-Diagramme in Origin. Schrittweite via
  \texttt{cfg.interp\_step}. \\
\bottomrule
\end{longtable}

\subsection{Plain-Text-Befund}
\label{sec:befund}

Neben den Excel-Dateien wird im Output-Verzeichnis ein
Plain-Text-Befund (\texttt{*\_REPORT.txt}) geschrieben. Inhalt:

\begin{itemize}
  \item Programm-Version, Aufruf-Datum, vollst\"andige Konfiguration
  \item Cluster-Erkennung (gefundene Cluster, ausgew\"ahltes Cluster
        bei Multi-Cluster)
  \item Liganden-Erkennungs-Tabelle (alle Fe-Liganden, ihre PDB-Reste,
        Bindungsl\"angen, Protonierungs-Status)
  \item Kabsch-Alignment-RMSD
  \item Sekund\"arstruktur-Statistik
  \item Anzahl Modes pro Frequenzfenster
  \item Marcus-Hush-Reorg-Diagnose (aktive Kan\"ale, integrale
        $\Lambda_X$-Werte)
  \item Vollst\"andige Liste aller Warnungen und Fehler aus dem Lauf
\end{itemize}

\begin{infobox}
\textbf{Empfehlung.} Den Befund \emph{immer} nach einem Lauf \"offnen,
bevor die Excel-Sheets interpretiert werden. Der Befund deckt alle
nicht-fatalen Probleme auf, die das Programm w\"ahrend der Analyse
stillschweigend abfedert (z.\,B.\ ein nicht-protoniertes His, das im
Lauf-Verlauf als ``unprotoniert'' geloggt wird, dessen PCET-Beitrag
deshalb null bleibt -- ohne Befund-Lekt\"ure leicht zu \"ubersehen).
\end{infobox}

% =========================================================================
\section{Modul-\"Ubersicht}
\label{sec:moduluebersicht}

Das Paket ist in zehn Module gegliedert. Tabelle nach
funktionalen Bereichen sortiert.

\begin{longtable}{p{0.18\textwidth} p{0.07\textwidth} p{0.70\textwidth}}
\toprule
\textbf{Modul} & \textbf{LOC} & \textbf{Aufgabe} \\
\midrule
\endhead

\texttt{config} & 684 & Konfigurations-Datenklasse \texttt{Config}
  mit allen 47 Feldern, Validierungslogik (\texttt{cfg.validate()}),
  Default-Werte. Siehe \S\ref{sec:konfig}. \\

\texttt{logio} & 1441 & Gaussian-Logfile-Reader (Block-Scan mit Cache,
  Eigenvektor-Parser, HP/Std-Konsistenzpr\"ufung), PDB-Reader,
  Sekund\"arstruktur-Erkennung (HELIX/SHEET, DSSP-Fallback, $\varphi/\psi$-
  Fallback), Lauf-Logger \texttt{RunLog}. \\

\texttt{geometry} & 1483 & Cluster-Erkennung
  (\texttt{find\_cluster}, \texttt{find\_all\_clusters}),
  Cluster-Normale via SVD, Kabsch-Alignment PDB$\leftrightarrow$Gauss,
  Liganden-Erkennung mit Element-spezifischen Cutoffs,
  His-Protonierungs-Test. \\

\texttt{core} & 1993 & Per-Mode-Analyse (\texttt{analyze\_mode\_with\_fallback}),
  thermische Skalierung, OOP/INP-Klassifikation
  (\texttt{classify\_oop\_inp}, bin\"ar + 7-Stufen),
  heuristische Kern-Klassifikation, SCSD-Adapter
  (\texttt{extract\_dominant\_scsd\_irreps}). Ruft
  \texttt{reorganization} fuer die Marcus-Hush-Modulationen auf. \\

\texttt{reorganization} & 1096 & Marcus-Hush-Reorg-Pipeline.
  Pro Mode geometrische Modulationen $\Delta r_X$ und
  $\lambda_X = \tfrac{1}{2}\mu\omega^2(\Delta r)^2$ fuer FeFe, FeN,
  FeS, NH, HA. HA im Reaktionskoordinaten-Modus, ohne
  Gauss-Daempfung in der Aggregation.
  System-Aggregation: Total-Reorg, Modulations-Spektren $M_X(\omega)$,
  kumulative Lambdas $\Lambda_X(\omega)$. Siehe \S\ref{sec:reorg}. \\

\texttt{pcet\_et} & 311 & Geometrie-Adapter fuer die
  Reorg-Pipeline. Atomgruppen-Bau (\texttt{build\_pcet\_info},
  \texttt{build\_et\_info}), H-Bond-Akzeptor-Suche. \\

\texttt{embedding} & 603 & UMAP-Embedding der 31-dim
  Feature-Matrix; HDBSCAN-Clusterung; SSE-bezogene
  Embedding-Variante (siehe \S\ref{sec:embeddings}). \\

\texttt{export} & 2063 & Excel-Schreiber fuer 25 Sheet-Typen,
  Stilformatierung (Mode-Typ-F\"arbung), Plot-Generierung. \\

\texttt{runner} & 1341 & Hauptablauf:
  \texttt{run\_analysis(cfg)}-Funktion, orchestriert Logfile-Scan
  $\to$ Cluster $\to$ Liganden $\to$ Mode-Loop $\to$
  Reorg-Aggregation $\to$ Embedding $\to$ Export. Multi-Window-
  Modus mit Reorg-Aggregaten pro Fenster. \\

\texttt{orca\_io} & 319 & ORCA-\texttt{.hess}-Reader (eigenst\"andiger
  Parser). \\

\texttt{cli} & 251 & argparse-basierter Kommandozeilen-Wrapper.
  Befehl \texttt{modenanalyse-2fe2s} ist hier registriert. \\

\bottomrule
\end{longtable}

\subsection{Datenfluss}

\begin{enumerate}
  \item \texttt{config.Config} wird vom User erstellt und an
        \texttt{runner.run\_analysis} \"ubergeben.
  \item \texttt{logio} scannt das Logfile, liest Atome (mit/ohne H),
        identifiziert HP- und Standard-Bl\"ocke, pr\"uft
        Konsistenz.
  \item \texttt{geometry} findet das Cluster, das Kabsch-Alignment
        zur PDB, die Liganden.
  \item \texttt{pcet\_et} baut die PCET-Atomlisten einmal vor (His-H
        plus Akzeptor-Paare); \texttt{reorganization} berechnet pro
        Mode $\Delta r_X$ und $\lambda_X$ f\"ur alle Bindungs-Kan\"ale.
  \item Mode-Loop in \texttt{runner}: pro Mode
        \texttt{core.analyze\_mode\_with\_fallback} $\to$ Result-Dict
        mit OOP/INP-Klassifikation, Kern-Scores, SCSD-Amplituden und
        Reorg-Beitr\"agen.
  \item \texttt{embedding} berechnet UMAP und HDBSCAN-Cluster auf der
        31-dim Feature-Matrix.
  \item \texttt{export} schreibt das Haupt-Excel mit allen Sheets und
        den Befund.
\end{enumerate}

% =========================================================================
\section{Validierung und UMAP -- Anwendungshinweise}

Die theoretischen Grundlagen des UMAP-Embeddings (Algorithmus,
Feature-Matrix, $n_\text{neighbors}$-Auto-Heuristik, HDBSCAN-Cluster,
volle Feature-Liste) sind in \S\ref{sec:embeddings} ausf\"uhrlich
beschrieben. Dieser Abschnitt sammelt Anwendungshinweise zur
Verwendung der UMAP-Sheets in der praktischen Auswertung.

\subsection{Reproduzierbarkeit (Seed-Test)}

UMAP ist nicht deterministisch -- es h\"angt von einem
Initialisierungs-Seed ab. Mit demselben Seed liefert ein wiederholter
Lauf exakt dasselbe Embedding. Mit unterschiedlichen Seeds \"andert
sich die Geometrie, aber \emph{nicht die Cluster-Topologie}: Modes,
die in einem Lauf zusammen sind, sollten in den meisten anderen
L\"aufen ebenfalls zusammen sein.

\texttt{modenanalyse\_2fe2s} fixiert den UMAP-Seed auf 42, damit das
Sheet bei wiederholten L\"aufen reproduzierbar bleibt. F\"ur eine
formale Robustheits-Studie m\"usste man manuell mehrere Seeds laufen
lassen und Cluster-Konsistenz pr\"ufen (Adjusted Rand Index, NMI).

\subsection{Z-Scores, Fisher-F und der Diskriminanz-Score}
\label{sec:zscore-fisher}

Das Sheet \texttt{UMAP\_Cluster\_Profil} zeigt nicht die rohen
Feature-Werte pro Cluster (die w\"aren \"uber Features mit
unterschiedlichen physikalischen Einheiten schwer interpretierbar),
sondern drei abgeleitete Statistiken: den \emph{Z-Score}, die
\emph{Fisher-F-Statistik} und den \emph{Diskriminanz-Score}. Sie
beantworten drei verschiedene Fragen, und die Kombination aller
drei ist es, die die Cluster-Labels physikalisch sinnvoll macht.
Wir definieren sie nacheinander.

\subsubsection*{Z-Score: wie ungew\"ohnlich ist dieser Cluster bei diesem Feature?}

F\"ur ein Feature $f$ (z.\,B.\ \texttt{kern\_OOP\%}) und ein Cluster
$C_k$:
\begin{equation}
  z_{k,f} = \frac{\bar{x}_{k,f} - \bar{x}_f}{\sigma_f}
\end{equation}
wobei $\bar{x}_{k,f}$ der Mittelwert von Feature $f$ innerhalb von
Cluster $C_k$ ist, $\bar{x}_f$ der globale Mittelwert von $f$
\"uber \emph{alle} Modes, und $\sigma_f$ die globale
Standardabweichung. Der Z-Score ist dimensionslos: $z = +2{,}0$
bedeutet, der Cluster liegt zwei Standardabweichungen \"uber dem
globalen Mittel f\"ur dieses Feature, $z = -1{,}5$ bedeutet, er
liegt 1,5 Standardabweichungen darunter.

Z-Scores sind also ein Cluster-Fingerabdruck: ein Cluster mit
$z_{\text{kern\_OOP}\%} = +2{,}5$,
$z_{\text{Umbrella}} = +2{,}1$ und
$z_{\text{Breathing}} = -0{,}1$ ist klar ein Out-of-Plane-Umbrella-Cluster,
unabh\"angig davon, was die absoluten Werte der zugrundeliegenden
Features sind. Deshalb ist das Z-Score-Profil der nat\"urliche Weg,
einem Cluster einen physikalischen Charakter zuzuweisen.

\paragraph{Warum global standardisiert wird, nicht pro Cluster.}
Wir benutzen das globale $\sigma_f$ im Nenner, nicht das
clusterinterne $\sigma$. Das globale $\sigma$ ist eine feste
Eigenschaft des Systems, sodass $z_{k,f}$-Werte direkt zwischen
Features und zwischen Clustern verglichen werden k\"onnen. Ein
clusterinternes $\sigma$ w\"urde den Score von Cluster-Gr\"o\ss e
und der Geometrie des Clusters selbst abh\"angig machen, was die
Frage, die wir beantworten wollen, verzerren w\"urde.

\subsubsection*{Fisher-F-Statistik: wie cluster-spezifisch ist dieses Feature?}

Ein Feature kann in einem Cluster einen gro\ss en Z-Score allein
durch Zufall haben, besonders bei kleinen Clustern. Die
Fisher-F-Statistik misst, ob das Feature \emph{tats\"achlich}
zwischen Clustern diskriminiert, indem sie die Varianz zwischen
Clustern mit der Varianz innerhalb von Clustern vergleicht:
\begin{equation}
  F_f = \frac{
    \sum_k n_k\,(\bar{x}_{k,f} - \bar{x}_f)^2 / (K - 1)
  }{
    \sum_k \sum_{i \in C_k} (x_{i,f} - \bar{x}_{k,f})^2 / (N - K)
  }
\end{equation}
wobei $K$ die Anzahl der Cluster ist, $N$ die Gesamtzahl der Modes
und $n_k$ die Gr\"o\ss e von Cluster $C_k$. Der Z\"ahler ist die
Varianz der Cluster-Mittel um das globale Mittel (Signal), der
Nenner ist die mittlere Varianz innerhalb der Cluster (Rauschen).
Berechnet \"uber \texttt{scipy.stats.f\_oneway}, ist dies die
Standard-Einweg-ANOVA-F-Statistik.

Interpretation:
\begin{itemize}\setlength{\itemsep}{0pt}
  \item $F_f \gg 1$: das Feature variiert deutlich mehr
        \emph{zwischen} Clustern als \emph{innerhalb} -- es ist
        ein diskriminierendes Feature, und Z-Scores auf diesem
        Feature sind statistisch sinnvoll.
  \item $F_f \approx 1$: zwischen-Cluster-Varianz ist \"ahnlich
        zur intra-Cluster-Varianz -- das Feature trennt die
        Cluster nicht; Z-Scores f\"ur dieses Feature sollten
        ignoriert werden, auch wenn sie gro\ss\ aussehen.
  \item $F_f \ll 1$ oder $F_f = 0$: das Feature ist im Wesentlichen
        konstant oder Rauschen (z.\,B.\ Null-Amplitude auf allen
        Modes f\"ur Sekund\"arstruktur-Features in einem System
        ohne dieses SSE-Element).
\end{itemize}

Die Fisher-F wird einmal pro Feature berechnet, unabh\"angig vom
Cluster -- sie ist eine globale Eigenschaft der Partition.

\subsubsection*{Diskriminanz-Score: welche Features definieren einen Cluster?}

Z-Score und Fisher-F beantworten unterschiedliche Fragen: der
Z-Score sagt uns, \emph{in welche Richtung} ein Cluster
ungew\"ohnlich ist, und der Fisher-F sagt uns, \emph{ob} diese
Richtung verl\"asslich ist. Ihre Kombination ist der
Diskriminanz-Score
\begin{equation}
  d_{k,f} = |z_{k,f}|\,\sqrt{F_f},
  \label{eq:discriminant}
\end{equation}
wobei der Betrag dient, damit ``ungew\"ohnlich hoch'' und
``ungew\"ohnlich niedrig'' gleich gewichtet werden. Die Wurzel auf
$F_f$ rechnet die F-Statistik von einer Varianz-Ratio in eine
Standardfehler-Skala um, sodass $|z|$ und $\sqrt{F}$ vergleichbare
Einheiten haben.

F\"ur jeden Cluster sortiert das Sheet \texttt{UMAP\_Cluster\_Profil}
die Features absteigend nach $d_{k,f}$ und listet die ersten als
charakteristische Eigenschaften des Clusters. Ein Feature steht
oben in der Liste, wenn es \emph{sowohl} einen gro\ss en Z-Score
hat (der Cluster ist auf diesem Feature ungew\"ohnlich) \emph{als
auch} eine gro\ss e Fisher-F (das Feature diskriminiert
tats\"achlich zwischen Clustern).

\paragraph{Beispielrechnung.}
Cluster~3 habe
$z_{\text{kern\_OOP}\%} = +2{,}5$ mit $F_{\text{kern\_OOP}\%} = 81$,
und $z_{\text{His255\_amplitude}} = +1{,}9$ mit
$F_{\text{His255\_amplitude}} = 4$.
Die Diskriminanz-Scores sind
$d_{3,\text{kern\_OOP}\%} = 2{,}5\,\sqrt{81} = 22{,}5$ und
$d_{3,\text{His255\_amplitude}} = 1{,}9\,\sqrt{4} = 3{,}8$.
Obwohl die Z-Scores \"ahnlich sind, rangiert \texttt{kern\_OOP\%}
sechsfach h\"oher in der Diskriminanz-Reihenfolge und wird als
dominanter Charakter von Cluster~3 ausgegeben. Die His255-Amplitude
hat in genau diesem Cluster den gr\"o\ss eren Z-Score, ist aber
\"uber alle Cluster hinweg im Wesentlichen Rauschen ($F = 4$ ist
grenzwertig), und sollte daher nicht als definierend gelesen
werden.

\subsubsection*{Was man in der Praxis mit den Zahlen tut}

\begin{itemize}\setlength{\itemsep}{2pt}
  \item Sheet \texttt{UMAP\_Cluster\_Profil} \"offnen.
  \item Pro Cluster die \emph{Top-3-Features} nach
        Diskriminanz-Score betrachten.
  \item Wenn $d > 5$ f\"ur das Top-Feature: der Cluster hat einen
        klaren physikalischen Charakter; ein Label vergeben
        (z.\,B.\ ``OOP-Umbrella'', ``Cluster-Breathing'',
        ``His-N-Stretch'').
  \item Wenn $d < 2$ f\"ur alle Features in einem Cluster: der
        Cluster ist statistisch schwach; als undifferenzierter
        Rausch-Cluster behandeln, nicht als eigenst\"andige
        Mode-Familie.
  \item \"Uber Cluster hinweg: schauen, \emph{welche Features
        immer wieder oben auftauchen}. Das sind die Richtungen,
        entlang derer UMAP die Modes tats\"achlich organisiert.
        Features, die in keinem Cluster oben auftauchen, tragen
        zum Embedding nichts bei.
\end{itemize}

F\"ur typische [2Fe-2S]-L\"aufe bei 5\,K mit dem 31-Feature-Standard
sind die dominierenden diskriminierenden Features in der Regel
\texttt{kern\_OOP\%}, \texttt{Breathing}, \texttt{Umbrella} und die
\texttt{His\_*\_Amplitude}-Spalten -- konsistent damit, dass die
Cluster-Modes den gr\"o\ss ten Teil der spektroskopischen
Information tragen.

\subsection{Sinnvolle Cluster?}

Nach dem Lauf das Sheet \texttt{UMAP\_Cluster\_Profil} \"offnen
(siehe \S\ref{sec:zscore-fisher} f\"ur die Bedeutung der Spalten).
Pro Cluster sollten die Z-Score-Profile physikalisch sinnvoll sein:
hoch korrelierte Features (z.\,B.\ \texttt{kern\_OOP\%} mit
\texttt{Umbrella}, oder \texttt{His 255\_INP} mit
\texttt{His 255\_bend\_inp}) landen typischerweise in denselben
Cluster-Profilen. Wenn ein Cluster im Profil keinen klaren
Charakter zeigt (alle Z-Scores nahe 0, alle Diskriminanz-Scores
unter 2), ist er statistisch unsicher und sollte ignoriert werden.

\subsection{Vergleich verschiedener Systeme im Embedding}

F\"ur direkten WT-vs.-Mut-Vergleich gibt es zwei sinnvolle
Strategien:

\begin{enumerate}
  \item \textbf{Getrennte Embeddings.} WT und Mut werden separat in
        UMAP eingebettet. Die Cluster-Topologien werden visuell
        verglichen. Vorteil: jedes System f\"ur sich aussagekr\"aftig.
        Nachteil: kein direkter Punkt-zu-Punkt-Vergleich.
  \item \textbf{Gemeinsames Embedding.} Beide Systeme zusammen in
        UMAP einbetten, mit System-Label als Farbcode. Vorteil:
        homologe Modes liegen im selben Embedding-Punkt. Nachteil:
        die Hyperparameter m\"ussen f\"ur die kombinierte Datenmenge
        getuned werden, und die einzelne Cluster-Struktur kann sich
        durch das andere System verschieben.
\end{enumerate}

\texttt{modenanalyse\_2fe2s} bietet aktuell nur die getrennte Variante.
F\"ur das gemeinsame Embedding m\"usste man die Feature-Matrizen
der zwei L\"aufe au\ss erhalb des Programms zusammenf\"uhren.


\section{Fehlerbehebung}
\label{sec:fehlerbehebung}

Diese Sektion enth\"alt die wichtigsten Fehlerbilder mit ihren
Diagnosen.

\subsection{Installation}

\subsubsection*{``ModuleNotFoundError'' nach erfolgter Installation}

Der Python-Kernel l\"auft noch mit dem Stand von \emph{vor} der
Installation. \emph{L\"osung}: Kernel neu starten (Spyder:
``Consoles $\to$ Restart kernel'', Strg+. ).

\subsubsection*{Mehrere Python-Installationen, falsche Umgebung}

Wenn z.\,B.\ \texttt{pip install} im System-Python installiert, aber
Spyder mit Anaconda-Python l\"auft. \emph{Pr\"ufung}:

\begin{lstlisting}[style=python]
import sys
print(sys.executable)
import modenanalyse_2fe2s
print(modenanalyse_2fe2s.__file__)
\end{lstlisting}

Die zwei Pfade m\"ussen zur gleichen Python-Installation geh\"oren.

\subsection{Lauf-Fehler}

\subsubsection*{Cluster nicht erkannt}

Symptom: Im Befund steht ``Kein {[}2Fe-2S{]}-Cluster gefunden''. M\"ogliche
Ursachen:

\begin{itemize}
  \item Fe-Fe-Abstand $>$ \texttt{fe\_fe\_cutoff} (Default 3.5 \AA{}).
        Pr\"ufen via PDB-Visualisierung; bei sehr stark verzerrten
        Geometrien Cutoff hochsetzen.
  \item S-Atome nicht br\"uckend (z.\,B.\ ein {[}2Fe-2S{]}-Modell mit
        nur einem br\"uckenden S). Programm erfordert vier Cluster-Atome.
  \item PDB enth\"alt das Cluster nicht, oder Fe-Atome haben nicht
        Element-Symbol \texttt{FE}. Manuell pr\"ufen.
\end{itemize}

\subsubsection*{Liganden-Erkennung scheitert (alle Liganden zeigen ``unknown'')}

\begin{itemize}
  \item Kabsch-Alignment fehlgeschlagen (RMSD $>$ 2 \AA{}). PDB und
        Gaussian-Cluster sind zu stark unterschiedlich; Liganden-
        Mapping per Distanz wird unzuverl\"assig.
  \item PDB-Kette stimmt nicht mit \texttt{cfg.pdb\_chain}
        \"uberein. Pr\"ufen mit einem PDB-Editor.
\end{itemize}

\subsubsection*{HP/Standard-Frequenz-Diskrepanz}

Im Befund steht eine Tabelle ``HP/Std-Frequenzdifferenzen''. Wenn:

\begin{itemize}
  \item 1--2 Modes mit $\Delta\omega \leq 0.05$ cm$^{-1}$:
        Numerische Reduzierung in Gaussian-Output, harmlos.
  \item Viele Modes mit systematischem Offset:
        Block-Pairing falsch. Logfile pr\"ufen, ggf.\ Job neu starten.
  \item Mode mit $\Delta\omega > 1$ cm$^{-1}$:
        Logfile m\"oglicherweise korrupt; Manuell pr\"ufen.
\end{itemize}

\subsubsection*{NIS-Spektrum sieht falsch aus}

\begin{itemize}
  \item \textbf{Zu schmale Peaks}: \texttt{nis\_fwhm\_gauss} zu klein
        gesetzt. Default 8.07 cm$^{-1}$ entspricht 1 meV instrumenteller
        Breite -- typisch f\"ur APS-Beamline ID-9.
  \item \textbf{Falsche Peaklage}: Logfile-Frequenzen sind nicht
        skaliert. Wenn experimentell harmonische Skalierungsfaktoren
        anliegen, m\"ussen sie \emph{vor} dem Lauf angewendet werden
        (in der Logfile-Bearbeitung). \texttt{modenanalyse\_2fe2s}
        skaliert die Frequenzen nicht selbst.
  \item \textbf{pDOS-Konsistenzcheck warnt vor $|1 - \int g|>10^{-3}$}:
        Frequenzgitter zu schmal. \texttt{nis\_freq\_max} hochsetzen,
        oder \texttt{nis\_n\_points} erh\"ohen.
\end{itemize}

\subsubsection*{Reorg-Sheets zeigen keine HA/NH-Beitr\"age}

Wenn $\Lambda_\text{HA}^\text{total}$ und $\Lambda_\text{NH}^\text{total}$
beide null sind:

\begin{enumerate}
  \item Im Befund nach ``PCET-Reorg: NH/HA-Kanaele uebersprungen''
        suchen. Wenn dort ``keine His-Liganden'' steht, ist das
        physikalisch korrekt (z.\,B.\ Cys$_4$-Cluster ohne His).
  \item Wenn ``His-Liganden gefunden, aber keine H-Bond-Akzeptoren'',
        \texttt{cfg.pcet\_hbond\_cutoff\_a} hochsetzen (z.\,B.\ auf
        4.5 oder 5.0 \AA{}).
  \item Wenn nach Erh\"ohen immer noch null: PDB enth\"alt keine
        Wassermolek\"ule oder andere Akzeptoren in der Cluster-Umgebung.
        Strukturmodell pr\"ufen, ggf.\ explizite Wasser-Molek\"ule
        ergaenzen.
\end{enumerate}

\subsection{Output und Excel}

\subsubsection*{Excel-Datei l\"asst sich nicht \"offnen}

\begin{itemize}
  \item Mehrere L\"aufe in dasselbe \texttt{output\_dir}: die alte
        Excel-Datei ist noch in Excel ge\"offnet. Schlie\ss en und neu
        anlegen.
  \item \texttt{openpyxl}-Version zu alt: $>=$ 3.1 erforderlich.
\end{itemize}

\subsubsection*{Sheet \texttt{Gruppen\_*} fehlt oder leer}

\begin{itemize}
  \item Keine Liganden erkannt (siehe oben).
  \item PDB-Datei nicht angegeben (\texttt{pdb\_file = ""}).
  \item PDB-Kette stimmt nicht.
\end{itemize}

\subsubsection*{Sheet \texttt{SCSD} fehlt}

\texttt{scsdpy} nicht installiert. \texttt{pip install scsdpy} f\"ur
das optionale Sheet, oder ignorieren -- alle anderen Sheets sind
vollst\"andig auch ohne SCSD.

\subsection{Performance}

\subsubsection*{Lauf braucht $>$10 Minuten}

\begin{itemize}
  \item Sehr gro\ss es Logfile ($>$1 GB): \texttt{cfg.scan\_chunk\_mb}
        auf 64 oder 128 setzen.
  \item Wiederholter Lauf desselben Logfiles: \texttt{cfg.use\_cache =
        True} aktivieren (Default). Block-Offsets werden auf Disk
        gecacht, sp\"atere L\"aufe halbieren ihre Logfile-Lese-Zeit.
  \item Lauf mit aktiviertem PCET ist ca.\ 30\,\% langsamer als ohne
        (zus\"atzlicher Eigenvektor-Read mit H-Atomen pro Mode). Wenn
        nicht ben\"otigt: \texttt{pcet\_enabled = False}.
\end{itemize}

\subsubsection*{Sehr gro\ss er Speicher-Verbrauch}

\texttt{modenanalyse\_2fe2s} h\"alt alle Mode-Resultate im RAM. Bei
sehr gro\ss en Systemen ($>$2000 Modes) und vielen Liganden kann der
Speicher 2--4 GB erreichen. \emph{L\"osung}: \texttt{freq\_max}
herabsetzen oder Multi-Window-Modus deaktivieren.


% =========================================================================
% Lizenz
% =========================================================================

\section{Lizenz und Zitierung}
\label{sec:lizenz}

\subsection{Lizenz}

\texttt{modenanalyse\_2fe2s} ist freie Software unter der
\textbf{GNU General Public License Version 3 oder sp\"ater (GPL-3.0-or-later)}.

Was die Lizenz erlaubt:

\begin{itemize}
  \item Frei verwenden -- f\"ur Forschung, Lehre und auch kommerziell.
  \item Modifizieren -- den Quellcode \"andern, anpassen, eigene
        Erweiterungen schreiben.
  \item Verteilen -- modifizierte oder unmodifizierte Versionen
        weitergeben (auch verkaufen).
\end{itemize}

Was die Lizenz \emph{verlangt}:

\begin{itemize}
  \item Modifikationen, die ver\"offentlicht oder weitergegeben werden,
        m\"ussen ebenfalls unter GPL-3.0-or-later stehen
        (\emph{Copyleft-Klausel}). Damit bleibt die Software
        \emph{open source}, auch in modifizierter Form.
  \item Bei Weiterverteilung muss der Quellcode mitgeliefert werden
        (oder ein klarer Hinweis, wie er bezogen werden kann).
  \item Der Copyright-Hinweis und die GPL-3.0-Klauseln d\"urfen nicht
        entfernt werden.
\end{itemize}

Der vollst\"andige Lizenztext liegt der Distribution als Datei
\texttt{LICENSE} bei. Jede Quelltextdatei tr\"agt im Header einen
SPDX-License-Identifier:

\begin{lstlisting}[style=python]
# Copyright (C) 2026 Lukas Knauer, AG Sch\"unemann, RPTU Kaiserslautern-Landau
# SPDX-License-Identifier: GPL-3.0-or-later
# Teil von modenanalyse_2fe2s -- siehe LICENSE im Wurzelverzeichnis.
\end{lstlisting}

\subsection{Zitierung}

Wenn \texttt{modenanalyse\_2fe2s} in einer wissenschaftlichen
Publikation verwendet wird, sollten:

\begin{itemize}
  \item Die Software selbst zitiert werden (Programm-Name, Version,
        Autor, Jahr, evtl.\ DOI/Zenodo-Eintrag, falls vergeben).
  \item Die zugrundeliegenden Methoden zitiert werden, sofern sie f\"ur
        die Auswertung relevant sind:
        \begin{itemize}
          \item NIS/NRVS-Theorie: Paulsen 1999 \cite{Paulsen1999};
                Sturhahn 2004 \cite{Sturhahn2004}.
          \item Maradudin-Bessel-Mehrphonon: Maradudin \& Flinn 1963
                \cite{MaradudinFlinn1963}.
          \item SCSD-Methode: Kingsbury \& Senge 2024
                \cite{Kingsbury2024}.
          \item PCET-Theorie: Hammes-Schiffer \& Stuchebrukhov 2010
                \cite{HammesSchiffer2010};
                Layfield \& Hammes-Schiffer 2014
                \cite{LayfieldHammesSchiffer2014}.
          \item Embedding-Methoden je nach Anwendung: PCA
                \cite{Pearson1901}; t-SNE \cite{vanDerMaaten2008};
                UMAP \cite{McInnes2018}; HDBSCAN \cite{Campello2013}.
        \end{itemize}
\end{itemize}

\subsection{Beitrag und Issue-Tracker}

Bug-Reports, Feature-Requests und Pull Requests sind willkommen \"uber
das GitHub-Repository des Projekts (Adresse siehe
\texttt{pyproject.toml}). Beitr\"age werden unter denselben Lizenz-
Bedingungen (GPL-3.0-or-later) integriert.

% =========================================================================
% Validierung
% =========================================================================

\section{Validierung an [2Fe-2S]-Modell-Clustern und Proteinen}
\label{sec:validierung}

Die Methodik der Marcus-Hush-Reorganisations-Pipeline wird auf zwei
Ebenen gepr\"uft: an einer Matrix von kleinen Modell-Clustern (rein
DFT-Optimierungen ohne Protein-Umgebung) und an einem
QM/MM-System mit voller Protein-Umgebung (mitoNEET-H87C-Mutante).

Beide Ebenen pr\"ufen unabh\"angige Eigenschaften des Tools:

\begin{itemize}
  \item \textbf{Modell-Matrix}: Sensitivit\"at der Cluster-Reorg
        ($\Lambda_\text{FeFe}$, $\Lambda_\text{FeS}$) auf Oxidationsstufe und
        Imidazol-Protonierung. Pr\"uft die DFT-Pipeline und die
        Marcus-Hush-Aggregation, ohne von der PDB-Liganden-
        Erkennung abzuh\"angen.
  \item \textbf{Protein-System}: Vollst\"andige Pipeline inklusive
        Kabsch-Alignment, PDB-Liganden-Identifikation und
        PCET-Pipeline. Pr\"uft die Tool-Logik, die nur mit
        Protein-Kontext relevant wird.
\end{itemize}

\subsection{Ebene 1: Modell-Matrix}
\label{sec:val-modell-matrix}

F\"unf [2Fe-2S]-Modell-Cluster wurden mit unterschiedlichen
Liganden-Topologien, Oxidationsstufen und Protonierungs-Zust\"anden
gerechnet:

\begin{table}[h]
\centering
\small
\begin{tabular}{|l|c|c|c|c|c|}
\hline
\textbf{System} & \textbf{Ligation} & \textbf{Oxid.} &
  \textbf{N-H?} & \textbf{Eigenvec} & \textbf{Modes} \\
\hline
A: Cys$_4$ red       & Cys$_4$    & Fe(II)+Fe(III) & --   & HP   & 66 \\
B: FeIII prot HP   & Cys$_2$His$_2$ & Fe(III)$_2$    & ja   & HP   & 81 \\
C: FeIII deprot      & Cys$_2$His$_2$ & Fe(III)$_2$    & nein & Std  & 78 \\
D: FeII prot         & Cys$_2$His$_2$ & Fe(II)$_2$     & ja   & Std  & 81 \\
E: FeII deprot       & Cys$_2$His$_2$ & Fe(II)$_2$     & nein & Std  & 79 \\
\hline
\end{tabular}
\caption{Validierungs-Matrix: f\"unf [2Fe-2S]-Modell-Cluster mit
unterschiedlichen Liganden, Oxidationsstufen und Protonierungs-
Zust\"anden. Eigenvec: HP = \texttt{freq=hpmodes} (5 Dezimalstellen
in Gaussian-Logfile), Std = Standard (4 Dezimalstellen, mit dem
v1.0.0 Standard-Fallback verarbeitet). Modes = nach
\texttt{freq\_max=800}-Filterung verbleibende Schwingungsmoden.}
\label{tab:validierung-matrix}
\end{table}

\paragraph{Geometrie-Befunde der Modell-L\"aufe}

\begin{table}[h]
\centering
\small
\begin{tabular}{|l|c|c|c|l|}
\hline
\textbf{System} & \textbf{Fe-Fe} & \textbf{$\langle$Fe-S$\rangle$} &
  \textbf{Fe-S-Fe} & \textbf{Bemerkung} \\
\hline
A: Cys$_4$ red    & 3.010 & 2.312 & 81.2$^{\circ}$ & Mischvalenz, leicht asymmetrisch \\
B: FeIII prot HP  & 2.856 & 2.217 & 80.1$^{\circ}$ & oxidiert, kompaktes Cluster \\
C: FeIII deprot   & 2.773 & 2.181 & 78.9$^{\circ}$ & deprot.\ Imidazol $\to$ k\"urzeres Fe-N \\
D: FeII prot      & 2.833 & 2.282 & 76.7$^{\circ}$ & reduziert, l\"angere Fe-S \\
E: FeII deprot    & 2.967 & 2.328 & 79.1$^{\circ}$ & reduziert + deprot, gr\"o{\ss}tes Cluster \\
\hline
\end{tabular}
\caption{Cluster-Geometrien der f\"unf Modell-Systeme (v1.0.0-L\"aufe).
Distanzen in $\angstrom$. Alle Geometrien wurden vom Tool als
"`Cluster-Geometrie OK"' klassifiziert.}
\label{tab:val-geometrie}
\end{table}

\paragraph{Reorganisations-Energien der Modell-L\"aufe}

\begin{table}[h]
\centering
\small
\begin{tabular}{|l|c|c|c|c|c|}
\hline
\textbf{System} & $\Lambda_\text{FeFe}$ & $\Lambda_\text{FeN}$ &
  $\Lambda_\text{FeS}$ & $\Lambda_\text{NH}$ & $\Lambda_\text{HA}$ \\
\hline
A: Cys$_4$ red       & 21.6 & 0.0 & 206.4 & 0.0 & 0.0 \\
B: FeIII prot HP   & 29.8 & 0.0 & 283.0 & 0.0 & 0.0 \\
C: FeIII deprot      & 34.4 & 0.0 & 170.6 & 0.0 & 0.0 \\
D: FeII prot         & 14.0 & 0.0 & 130.1 & 0.0 & 0.0 \\
E: FeII deprot       & 20.8 & 0.0 & 183.7 & 0.0 & 0.0 \\
\hline
\end{tabular}
\caption{Reorganisations-Energien $\Lambda_X$ in cm$^{-1}$
($T=5$\,K, \texttt{freq\_max=800}). Reine Cluster-L\"aufe
\emph{ohne PDB-File}; daher sind die Liganden-Kan\"ale (FeN, NH, HA)
identisch null. Die Cluster-internen Kan\"ale (FeFe, FeS) zeigen
chemisch konsistente Trends.}
\label{tab:val-reorg-modell}
\end{table}

\paragraph{Was die Modell-Matrix zeigt}
\begin{itemize}\setlength{\itemsep}{0pt}
  \item \textbf{Geometrie-Trends sind chemisch plausibel.}
        Fe-Fe-Distanz korreliert mit Oxidationsstufe (Fe(III)$_2$
        kompakt, Fe(II)$_2$ aufgespreizt). Fe-S-Mittelwert
        korreliert mit Fe-Ionenradius (Fe(III) klein, Fe(II)
        gr\"o{\ss}er). Imidazol-Deprotonierung verk\"urzt Fe-S leicht
        (st\"arkerer $\sigma$-Donor).
  \item \textbf{$\Lambda_\text{FeS}$ liegt in derselben Gr\"o{\ss}enordnung}
        \"uber alle f\"unf Modelle (130--283\,cm$^{-1}$), konsistent
        mit dem erwarteten Bereich
        $(\hbar\omega/4)\cdot\alpha_X^2$ f\"ur Mode-Anteile zwischen
        30 und 60\,\% in der 300--400\,cm$^{-1}$-Region.
  \item \textbf{$\Lambda_\text{FeFe}$ ist klein und stabil}
        (14--34\,cm$^{-1}$), wie f\"ur Cluster-Atmungsmodes
        erwartet.
  \item \textbf{Wichtig zu wissen}: Da die Modell-L\"aufe \emph{ohne PDB}
        gerechnet wurden, sind die Liganden-Kan\"ale (FeN, NH, HA)
        identisch null -- nicht weil das Tool falsch rechnet,
        sondern weil ohne PDB die Liganden-Reste nicht
        identifiziert werden k\"onnen. F\"ur die volle PCET-Pipeline
        muss das Tool mit PDB-Datei laufen (siehe Ebene 2).
\end{itemize}

\subsection{Ebene 2: Vollst\"andiges Protein-System
            (mitoNEET H87C, QM/MM)}
\label{sec:val-mitoneet-h87c}

Als Validierungs-Beispiel mit voller Pipeline wurde eine
QM/MM-Frequenzrechnung der \textbf{mitoNEET-H87C-Mutante}
durchgef\"uhrt (ORCA-6, B3LYP/def2-TZVP-Niveau, QM/MM-Setup).

\paragraph{Warum H87C?}
mitoNEET-Wildtyp ist Cys$_3$His$_1$-koordiniert (Cys72, Cys74,
Cys83, His87). Die H87C-Mutante \emph{ersetzt} His87 durch ein
zus\"atzliches Cystein -- damit ist der Cluster Cys$_4$-koordiniert,
und der PCET-Pfad \"uber Imidazol-Deprotonierung ist
\emph{strukturell deaktiviert}. Das ist ein idealer
\textbf{Negativ-Test} f\"ur die PCET-Pipeline: das Tool muss
$\Lambda_\text{NH} = \Lambda_\text{HA} = 0$ vorhersagen, ohne dass
chemische "`Lecks"' (z.\,B.\ residuelle H-Bond-Akzeptoren von
benachbarten Resten) f\"alschlich Beitr\"age erzeugen.

\paragraph{System-Eigenschaften}
\begin{itemize}\setlength{\itemsep}{0pt}
  \item \textbf{422 Atome}, 1266 Schwingungsmoden in der ORCA-Hessian
  \item Nach \texttt{freq\_max=800}-Filter: \textbf{736 Moden} im
        relevanten NRVS-Bereich
  \item Eigenvektoren durchgehend in voller Genauigkeit
        (ORCA \texttt{.hess}, kein HP/Std-Unterschied)
  \item PDB-Datei mit korrekter HIS/CYS-Annotation der Liganden-
        Reste
\end{itemize}

\paragraph{Vom Tool erkannte Cluster-Geometrie und Liganden}

\begin{verbatim}
Cluster: Fe-Fe = 2.779 A, Fe-S(mittel) = 2.261 A,
         Fe-S-Fe = 75.9 deg - OK

Cluster-Liganden:
  Fe1 <- Cys 72 (S, SG, d=2.347 A)
  Fe1 <- Cys 74 (S, SG, d=2.311 A)
  Fe2 <- Cys 83 (S, SG, d=2.288 A)
  Fe2 <- Cys 87 (S, SG, d=2.327 A)

PCET-Reorg: NICHT aktiv (keine His-Liganden am Cluster).
ET-Reorg-Liganden: 4 Atome erkannt.
\end{verbatim}

\paragraph{Vier Cys-Liganden korrekt identifiziert}: Das Tool hat
trotz fehlender HIS-Beteiligung am Cluster die vier Cys-Liganden
\"uber Kabsch-Alignment + PDB-Mapping korrekt zugeordnet
(Distanzen 2.29--2.35\,$\angstrom$, alle plausibel). Damit ist die
Liganden-Erkennungs-Pipeline auf einem realen Protein-System
erfolgreich getestet.

\paragraph{Ergebnis-Reorg}

\begin{table}[h]
\centering
\begin{tabular}{|l|c|c|}
\hline
\textbf{Kanal $X$} & \textbf{Wert (cm$^{-1}$)} & \textbf{Bewertung} \\
\hline
$\Lambda_\text{FeFe}$ & 28.7  & \checkmark\ Cluster-Atmung im Sanity-Bereich \\
$\Lambda_\text{FeN}$  & 0.0   & \checkmark\ Negativ -- kein His-N \\
$\Lambda_\text{FeS}$  & 337.9 & \checkmark\ Cys$_4$-Cluster, hoher Beitrag \\
$\Lambda_\text{NH}$   & 0.0   & \checkmark\ Negativ -- kein Imidazol-NH \\
$\Lambda_\text{HA}$   & 0.0   & \checkmark\ Negativ -- kein H-Akzeptor-Pfad \\
\hline
\end{tabular}
\caption{Reorganisations-Energien des mitoNEET-H87C-Systems
($T=5$\,K, \texttt{freq\_max=800}). Drei der f\"unf Kan\"ale sind
exakt null: das ist \emph{nicht} ein null-Lauf, sondern die
\emph{korrekte Vorhersage} f\"ur ein System ohne His-Liganden. Der
Cluster-FeFe- und Fe-S-Cys-Beitrag liegt im erwarteten Bereich.}
\label{tab:val-reorg-mNT-H87C}
\end{table}

\paragraph{Mode-Klassifikation in mNT-H87C}
\begin{itemize}\setlength{\itemsep}{0pt}
  \item In-plane Modes:    \textbf{509} (69\,\%)
  \item Out-of-plane Modes: \textbf{91} (12\,\%)
  \item Torsional Modes:   \textbf{136} (18\,\%)
\end{itemize}
Die OOP/INP-Verteilung ist konsistent mit der Erwartung f\"ur ein
voll solvatisiertes Protein-System: viele In-plane-Skelett-Modes
des Backbones, einige OOP-Cluster-Modes, der Rest gemischt-Charakter.

\paragraph{Frequenzaufgel\"oste Verteilung der Reorg-Beitr\"age}

Die Total-Werte $\Lambda_X^\text{total}$ verteilen sich \"uber den
NRVS-relevanten Frequenzbereich wie folgt
(aus \texttt{Lambda\_kumulativ}-Sheet):

\begin{table}[h]
\centering
\small
\begin{tabular}{|r|r|r|l|}
\hline
\textbf{Bereich (cm$^{-1}$)} & $\Delta\Lambda_\text{FeFe}$ &
  $\Delta\Lambda_\text{FeS}$ & \textbf{Bandenzuordnung} \\
\hline
0--100   & 0.02   & 0.10    & Skelett-Modes, kaum Cluster-Aktivit\"at \\
100--200 & 3.74   & 1.49    & Niederfrequenz-Cluster-Bewegungen \\
200--300 & 0.37   & 12.22   & Cluster-Atmung-Region \\
300--400 & 1.86   & 122.52  & \textbf{Fe-S symm.\ Streckung dominant} \\
400--500 & 22.65  & 198.39  & \textbf{Fe-S asymm.\ + Fe-Fe gekoppelt} \\
500--800 & 0.04   & 3.17    & Liganden-Mix-Modes \\
\hline
\textbf{Total} & \textbf{28.68} & \textbf{337.89} & \\
\hline
\end{tabular}
\caption{Frequenzaufgel\"oste Verteilung der Reorg-Beitr\"age in
mNT-H87C. Die FeS-Pipeline konzentriert ihren Beitrag bei
300--500\,cm$^{-1}$ (95\,\% des Total-Werts). Die FeFe-Pipeline
ist stark im 400--500\,cm$^{-1}$-Bereich (79\,\% des Total-Werts) --
das deutet auf gekoppelte Cluster-Atmungs- und Fe-S-asymm.-Streck-Modes
hin.}
\label{tab:val-mNT-frequenz}
\end{table}

Diese Verteilung ist mit den literaturbekannten NRVS-Banden f\"ur
[2Fe-2S]-Cluster konsistent
\cite{Gee2021} und die NRVS-Mode-\"Ubersicht Wang 2025 \cite{Wang2025} abgleichen:
\begin{itemize}\setlength{\itemsep}{0pt}
  \item 300--400\,cm$^{-1}$: typischer Bereich f\"ur Fe-S$_b$
        symmetrische Streckschwingungen (324, 358\,cm$^{-1}$ in
        mitoNEET-WT \cite{Gee2021}).
  \item 400--500\,cm$^{-1}$: typischer Bereich f\"ur Fe-S$_b$
        asymmetrische Streckschwingungen (395, 413\,cm$^{-1}$
        in mitoNEET-WT \cite{Gee2021}). Die starke Kopplung an
        Cluster-Atmung wird hier durch den hohen FeFe-Beitrag
        sichtbar.
\end{itemize}

\paragraph{Was der mNT-H87C-Lauf validiert}

\begin{enumerate}
  \item \textbf{ORCA-Reader}: Liest die volle 1266-Mode-Hessian
        sauber, alle 736 Modes im Filter-Bereich werden mit
        HP-Genauigkeit verarbeitet. Pipeline-Tests aus
        \S\ref{sec:workflow-worked-example} (Marcus-Hush-Additivit\"at,
        Sanity-Range, Channel-Selektivit\"at) sind alle erf\"ullt.
  \item \textbf{Kabsch-Alignment}: PDB- und ORCA-Geometrie werden
        \"uber alle 4 Fe + 2 S-Atome korrekt aligniert (RMSD
        $\approx 0.000$\,$\angstrom$, da QM/MM-PDB direkt aus der
        ORCA-Optimierung stammt).
  \item \textbf{Liganden-Identifikation}: Alle vier Cys-Liganden
        (72, 74, 83, 87) werden korrekt aus dem PDB ausgelesen
        und den Cluster-Fe-Atomen zugeordnet.
  \item \textbf{PCET-Negativ-Test}: $\Lambda_\text{FeN}, \Lambda_\text{NH},
        \Lambda_\text{HA}$ sind exakt null -- die Pipeline schaltet
        diese Kan\"ale korrekt ab, wenn keine His-Liganden vorhanden
        sind. \emph{Keine residuellen Beitr\"age aus benachbarten
        Resten}, kein numerisches Rauschen.
  \item \textbf{Sekund\"arstruktur-Auto-Erkennung}: Da das PDB keine
        HELIX/SHEET-Records enth\"alt, wurden 3 SSE-Elemente per
        DSSP-H-Br\"ucken-Analyse rekonstruiert (Fallback-Pfad
        funktioniert).
\end{enumerate}

\subsection{Ausblick: Cys$_3$His-Wildtyp-Validierung}
\label{sec:val-cys3his-wt}

F\"ur die volle Validierung der \emph{positiven} PCET-Pipeline
(d.\,h.\ $\Lambda_\text{HA} > 0$ in einem realen His-haltigen Cluster)
ist eine analoge Rechnung am \textbf{mitoNEET-Wildtyp}
(Cys$_3$His$_1$) der n\"achste logische Schritt. Vergleich
zwischen WT und H87C-Mutante:

\begin{itemize}\setlength{\itemsep}{0pt}
  \item Bei \textbf{WT prot}: erwartet $\Lambda_\text{FeN} > 0$
        (Fe-N(His87)-Streckung bei $\sim 295$\,cm$^{-1}$
        \cite{Gee2021}); $\Lambda_\text{HA} > 0$ aus
        H-Akzeptor-Modulationen mit benachbarten Resten;
        $\Lambda_\text{NH} \approx 0$, weil N-H-Streckschwingungen
        bei $\sim 3300$\,cm$^{-1}$ liegen, au{\ss}erhalb
        \texttt{freq\_max=800}.
  \item Bei \textbf{WT deprot}: erwartet $\Lambda_\text{FeN} > 0$
        (Fe-N$^-$-Bindung bleibt aktiv); $\Lambda_\text{HA} \to 0$
        (kein N-H mehr verf\"ugbar).
  \item Bei \textbf{H87C} (dieser Lauf): bereits gezeigt, alle drei
        Liganden-Kan\"ale = 0.
\end{itemize}

Diese Drei-Punkt-Validierungs-Strategie (WT prot $>$ WT deprot
$>$ H87C) erlaubt eine quantitative Trennung zwischen
\emph{Liganden-induzierter} (FeN aus Fe-N$_\delta$) und
\emph{Reaktions-spezifischer} (HA aus N-H-Akzeptor-Geometrie)
Reorg-Beitr\"age. Gee, Pelmenschikov \emph{et al.}\ (2021)
\cite{Gee2021} liefern die experimentellen NRVS-Banden f\"ur den
direkten Vergleich.

\subsection{Erfolgskriterien des Validierungs-Tests}
\label{sec:val-erfolgskriterien}

\begin{enumerate}
  \item \textbf{Geometrie-Konsistenz} (alle f\"unf Modell-Systeme +
        mNT-H87C). \emph{Erf\"ullt}.
  \item \textbf{Cluster-Reorg-Trends} \"uber Oxidationsstufe und
        Protonierung. \emph{Erf\"ullt} f\"ur Modell-Matrix.
  \item \textbf{Marcus-Hush-Additivit\"at}: $\sum_i \lambda_X(i) =
        \Lambda_X^\text{total}$ exakt (auf 0.01\,cm$^{-1}$).
        \emph{Erf\"ullt} (Smoke-Test 4 in \texttt{tests/test\_smoke.py}).
  \item \textbf{Channel-Selektivit\"at}: Cys$_4$-Systeme haben
        $\Lambda_\text{FeN} = \Lambda_\text{NH} = \Lambda_\text{HA} = 0$.
        \emph{Erf\"ullt} f\"ur Modell A und mNT-H87C.
  \item \textbf{Sanity-Range}: $\Lambda_\text{FeFe} \in [1, 100]$,
        $\Lambda_\text{FeS} \in [50, 1000]$\,cm$^{-1}$. \emph{Erf\"ullt}
        f\"ur alle sechs Systeme.
  \item \textbf{ORCA-Pipeline}: Volle .hess-Datei wird gelesen,
        Kabsch-Alignment, Liganden-Erkennung. \emph{Erf\"ullt} f\"ur
        mNT-H87C (3 Bugs der ORCA-Pipeline in v1.0.0 behoben, siehe Changelog).
  \item \textbf{NRVS-Bandenzuordnung} \"uber kumulative
        $\Lambda_X(\omega)$-Kurven \cite{Gee2021,Wang2025}.
        \emph{Steht aus} -- erfordert mNT-Wildtyp-Lauf zum
        direkten Vergleich.
\end{enumerate}

\begin{infobox}
\textbf{Status der Validierung} ($\today$, v1.0.0).
Sechs of sieben Validierungs-Kriterien sind erf\"ullt: f\"unf
Modell-Cluster-L\"aufe (Geometrie + Cluster-Reorg) plus eine
volle Protein-Pipeline (mNT-H87C, Negativ-Test). Das Tool
funktioniert auf realen QM/MM-Frequenzrechnungen, identifiziert
Liganden korrekt, und schaltet die PCET-Kan\"ale strukturell ab,
wenn keine His-Liganden vorhanden sind.

Die noch fehlende positive PCET-Validierung (mNT-WT mit Cys$_3$His)
ist ein logischer n\"achster Schritt -- die Methode ist hier nicht
mehr methodisch zu validieren, sondern an einem konkreten
Forschungs-System anzuwenden.
\end{infobox}

Wenn alle sieben Kriterien erf\"ullt sind, ist die Marcus-Hush-
Reorg-Methodik publikations-validiert und kann auf neue Systeme
(unbekannte Cys$_2$His$_2$-Cluster, Doppel-Mutanten,
Multi-Cluster-Systeme) angewandt werden.

% =========================================================================
% ANHAENGE
% =========================================================================
\appendix

\part{Anh\"ange}

% -------------------------------------------------------------------------
\section{Mathematische Notation und Konventionen}
\label{app:notation}

Dieser Anhang fasst die im gesamten Handbuch verwendete
mathematische Notation, Einheiten-Konventionen und Symbol-
Bedeutungen kompakt zusammen. Er ist als Nachschlagereferenz
gedacht: wer in einem Kapitel \"uber ein Symbol oder eine
Einheiten-Konvention stolpert, findet hier die endgueltige
Definition.

\subsection{Skalare, Vektoren, Matrizen, Tensoren}

\begin{itemize}\setlength{\itemsep}{2pt}
  \item \textbf{Skalare}: kursiv, einfacher Buchstabe -- $a$, $b$,
        $\omega$, $T$. Im Fall vorzeichenfreier physikalischer
        Gr\"o{\ss}en wie Massen oder Frequenzen wird das positive
        Vorzeichen impliziert.
  \item \textbf{Vektoren}: fettgedruckt, einfacher Buchstabe --
        $\mathbf{r}$, $\mathbf{e}$, $\mathbf{a}$. Komponenten als
        $r_x, r_y, r_z$ oder $r_1, r_2, r_3$ je nach Kontext.
  \item \textbf{Matrizen}: fettgedruckt und gro{\ss}geschrieben --
        $\mathbf{R}$, $\mathbf{H}$, $\mathbf{F}$. Eintraege als
        $R_{ij}$ oder $(\mathbf{R})_{ij}$.
  \item \textbf{Mengen, Tupel}: kursiv, gro{\ss}geschrieben --
        $\mathcal{C}$ (Cluster-Atome), $\mathcal{L}$
        (Liganden-Mengen), $\mathcal{M}$ (Modes).
\end{itemize}

\subsection{Indizes}

\begin{itemize}\setlength{\itemsep}{2pt}
  \item $i, j, k$: laufen \"uber Schwingungsmoden ($i = 1, \ldots,
        N_\text{modes}$)
  \item $a, b, c$: laufen \"uber Atome ($a = 1, \ldots, N_\text{atoms}$)
  \item $\alpha, \beta$: laufen \"uber Kartesische Komponenten
        ($x, y, z$)
  \item $X, Y, Z$: stehen f\"ur Bindungs-Kan\"ale
        ($X \in \{$FeFe, FeN, FeS, NH, HA$\}$)
  \item $p, q$: laufen \"uber Kanal-Sub-Indizes (z.\,B.\ verschiedene
        Fe-N-Bindungen innerhalb des FeN-Kanals)
\end{itemize}

\subsection{Einheiten-Konventionen}

Das gesamte Handbuch und die Software verwenden konsequent
folgende Einheiten:

\paragraph{L\"angen.}
$\angstrom$ (1\,$\angstrom$ = $10^{-10}$\,m). PDB-Standard. Cluster-
Geometrien und Bindungsabstaende werden in $\angstrom$ angegeben.
Im SI-internen Code-Pfad wird zur Umrechnung
$r_\text{SI} = r_{\angstrom} \cdot 10^{-10}$\,m verwendet.

\paragraph{Frequenzen.}
cm$^{-1}$ (Wellenzahlen, NRVS-Konvention). Wandlung in SI:
\begin{equation}
  \omega_\text{SI} = 2\pi c\,\omega_{\text{cm}^{-1}}
  \quad\text{mit}\quad c = 2.99792458 \times 10^{10}\,\text{cm/s}
\end{equation}
und in Energie:
\begin{equation}
  E_\text{cm}^{-1} = \frac{E_\text{J}}{h c}
  \quad\text{mit}\quad h c = 1.98644586 \times 10^{-23}\,\text{J\,cm}
\end{equation}

\paragraph{Massen.}
amu (atomic mass units; 1\,amu = $1.66053906660 \times 10^{-27}$\,kg).
Reduzierte Massen $\mu$ werden \emph{immer} in amu angegeben, auch
wenn sie aus Mode-Eigenvektoren ($\mu_\text{mode}$) berechnet sind.

\paragraph{Energien.}
cm$^{-1}$ in den Reorg-Sheets, J im SI-internen Code-Pfad. Wandlung
am Ende der Berechnung im Modul \texttt{reorganization} via
\verb|_HC_JCM = h*c|.

\paragraph{Temperaturen.}
K (Kelvin). Tieftemperatur-NRVS: 5\,K (Standard im hier
beschriebenen Tool); Standard-NRVS: 40\,K; Raumtemperatur: 300\,K.

\paragraph{Auslenkungen, Modulationen.}
Die thermisch skalierten Mode-Auslenkungen
$\mathbf{e}_a^\text{thermal} = \mathbf{e}_a^\text{normalized} \cdot
u_\text{rms}$ haben Einheit $\angstrom$. Die daraus berechneten
Bindungs-Modulationen $\Delta r_X$ ebenfalls $\angstrom$.

\subsection{Vorzeichen-Konventionen}

\paragraph{Bindungs-Modulationen $\Delta r_X$.}
Vorzeichen-behaftet: positiver Wert bedeutet \emph{Streckung} der
Bindung in der gegebenen Mode-Auslenkung; negativer Wert bedeutet
Stauchung. F\"ur Reorg-Beitr\"age $\lambda_X = \tfrac{1}{2}\mu\omega^2
(\Delta r_X)^2$ ist nur das Quadrat relevant; bei der Aggregation
\"uber Sub-Kan\"ale (z.\,B.\ FeN\_His1 + FeN\_His2 $\to$ FeN) wird
\emph{vorzeichenfrei} aggregiert (RSS, siehe \S\ref{app:notation}).

\paragraph{Kabsch-Alignment-Vorzeichen.}
Das Kabsch-Alignment liefert eine Rotation, die die PDB-Geometrie
auf die Gauss-Geometrie abbildet. Bei links-/rechts-h\"andigen
Spiegelungen wird das Vorzeichen der z-Komponente korrigiert
(siehe \texttt{geometry.kabsch\_align}).

\paragraph{Cluster-Normale $\hat{\mathbf{n}}$.}
Per Konvention zeigt $\hat{\mathbf{n}}$ in die Richtung mit positiver
$z$-Komponente in der Cluster-Lokalachsen-Konvention. Diese
Wahl ist beliebig, aber konsistent in allen Sheets gepflegt.
Eine Mode mit OOP\% > 60 hat den \"uberwiegenden Teil ihrer
quadrierten Auslenkung in $\hat{\mathbf{n}}$-Richtung; ob diese
Auslenkung "`nach oben"' oder "`nach unten"' zeigt, ist
spektroskopisch egal.

\subsection{Summen-, Produkt- und Erwartungswert-Konventionen}

\paragraph{Summe \"uber Mode-Indizes.}
$\sum_i$ ohne weitere Spezifikation l\"auft \emph{immer} \"uber alle
nach Filterung im Frequenzfenster verbleibenden Schwingungsmoden,
nicht \"uber alle $3N_\text{atoms} - 6$ Moden des
Gesamtsystems.

\paragraph{Quantum-Mechanischer Harmonischer Oszillator (QHO).}
$\langle x^2 \rangle = \frac{\hbar}{2 m \omega}\coth(\hbar\omega/2k_BT)$
ist der Erwartungswert des Auslenkungs-Quadrats im thermischen
Gleichgewicht. Bei $T \to 0$ wird das zu $\hbar/(2m\omega)$
(Nullpunkts-Energie). Im Tool wird daraus
\begin{equation}
  u_\text{rms} = \sqrt{\langle x^2\rangle}
\end{equation}
abgeleitet (siehe \texttt{core.compute\_thermal\_amplitude}).

\paragraph{Reduzierte Massen-Konvention.}
F\"ur eine Bindung zwischen den Atomen $a$ und $b$:
\begin{equation}
  \mu_\text{pair}(a, b) = \frac{m_a m_b}{m_a + m_b}
\end{equation}
Diese ist \emph{eindeutig} pro Bindung, unabh\"angig von Mode-Index.
Im Gegensatz dazu ist $\mu_\text{mode}(i)$ die effektive reduzierte
Masse der Mode $i$, wie sie Gaussian/ORCA in der Frequenzrechnung
ausgibt -- diese ist mode-spezifisch.

\subsection{Notation f\"ur algorithmische Schritte}

Im sp\"ater zu erg\"anzenden Pseudocodes-Anhang werden folgende Konventionen
verwendet:

\begin{itemize}\setlength{\itemsep}{2pt}
  \item $a \gets b$: Zuweisung
  \item $a == b$: Gleichheits-Test
  \item $\textsc{Function}(a, b)$: Funktions-Aufruf in Kapitel\"alchen
  \item \texttt{procedure}: Pseudocode-Block (kein expliziter Output)
  \item \texttt{function}: Pseudocode-Block (mit Return-Wert)
\end{itemize}

% -------------------------------------------------------------------------
\section{Symbolverzeichnis}
\label{app:symbols}

Sortierte Liste aller im Handbuch und in der Software verwendeten
Symbole. Die Reihenfolge folgt der Konvention: lateinische
Buchstaben, dann griechische, dann Sonderzeichen.

\subsection{Lateinische Buchstaben}

\begin{tabular}{p{1.5cm} p{6cm} p{4cm}}
  \toprule
  \textbf{Symbol} & \textbf{Bedeutung} & \textbf{Einheit / Wertebereich} \\
  \midrule
  $a, b, c$       & Atom-Indizes        & $1, \ldots, N_\text{atoms}$ \\
  $B$             & Debye-Waller-Tensor / B-Faktor & $\angstrom^2$ \\
  $c$             & Lichtgeschwindigkeit & $2.998 \times 10^{10}$\,cm/s \\
  $\mathcal{C}$   & Menge der Cluster-Atome (Fe + S$_\text{bridge}$) & $|\mathcal{C}| = 4$ \\
  $C_\text{clus}$ & Cluster-Beteiligung einer Mode & $[0, 1]$ \\
  $C_\text{PCET}$ & Co-Modulations-Indikator & $[0, 1]$ \\
  $\Delta r_X$    & Bindungs-Modulation des Kanals $X$ & $\angstrom$ \\
  $\mathbf{e}_a$  & Eigenvektor-Komponente an Atom $a$ & $\angstrom$ (thermisch skaliert) \\
  $E$             & Energie & J oder cm$^{-1}$ \\
  $f_\text{LM}$   & Lamb-M\"o{\ss}bauer-Faktor & $[0, 1]$ \\
  $\mathbf{F}$    & Frequenz-Vektor & cm$^{-1}$ \\
  $h$             & Plancksches Wirkungsquantum & $6.626 \times 10^{-34}$\,Js \\
  $\hbar$         & Reduzierte Planck-Konstante & $h/(2\pi)$ \\
  $H$             & Hesse-Matrix der Kraftkonstanten & $\text{a.u.}$ in Gaussian \\
  $\hat{H}$       & Hamilton-Operator & J \\
  $i, j, k$       & Mode-Indizes        & $1, \ldots, N_\text{modes}$ \\
  $k_B$           & Boltzmann-Konstante & $1.381 \times 10^{-23}$\,J/K \\
  $\mathcal{L}$   & Menge der Liganden-Atome (Fe-koordiniert) & $|\mathcal{L}| \approx 4$--$8$ \\
  $m_a$           & Masse des Atoms $a$ & amu \\
  $\mathbf{n}$    & Cluster-Normale (durch Best-Fit-Ebene) & dimensionsloser Einheitsvektor \\
  $N_\text{atoms}$ & Anzahl der Atome im System & $\approx 50$--$3000$ \\
  $N_\text{modes}$ & Anzahl der Schwingungsmoden im Filterbereich & $\approx 100$--$10\,000$ \\
  $\mathbf{r}_a$  & Position des Atoms $a$ in Kartesischen Koord. & $\angstrom$ \\
  $r_\text{eq}$   & Gleichgewichts-Bindungsabstand & $\angstrom$ \\
  $S$             & Huang-Rhys-Faktor & dimensionslos \\
  $T$             & Temperatur & K \\
  $u_\text{rms}$  & RMS-Amplitude einer Mode (QHO-Erwartungswert) & $\angstrom$ \\
  $X, Y, Z$       & Bindungs-Kanal-Index & $\in \{\text{FeFe, FeN, FeS, NH, HA}\}$ \\
  \bottomrule
\end{tabular}

\subsection{Griechische Buchstaben}

\begin{tabular}{p{1.5cm} p{6cm} p{4cm}}
  \toprule
  \textbf{Symbol} & \textbf{Bedeutung} & \textbf{Einheit / Wertebereich} \\
  \midrule
  $\alpha$        & In SCSD: D$_\text{2h}$-Irrep-Koeffizient & $\angstrom$ \\
  $\alpha_X$      & Mode-Anteil auf Bindung $X$ & $[0, 1]$ \\
  $\Delta Q$      & Statische Mode-Verschiebung (Marcus-Hush) & amu$^{1/2}\angstrom$ \\
  $\Delta r_X$    & Bindungs-Modulation & $\angstrom$ \\
  $\hat{\mathbf{n}}$ & Cluster-Normale & dimensionsloser Einheitsvektor \\
  $\lambda$       & Marcus-Hush-Reorganisationsenergie & cm$^{-1}$ \\
  $\lambda_X(i)$  & Pro-Mode-Beitrag, Kanal $X$, Mode $i$ & cm$^{-1}$ \\
  $\Lambda_X^\text{total}$ & System-Total-Reorg, Kanal $X$ & cm$^{-1}$ \\
  $\Lambda_X(\omega)$ & Kumulative Reorg-Kurve & cm$^{-1}$ \\
  $\mu_\text{pair}$ & Reduzierte Masse einer Bindung & amu \\
  $\mu_\text{mode}$ & Reduzierte Masse einer Mode (aus DFT) & amu \\
  $\sigma$        & Gauss-Verbreiterung der Spektren & cm$^{-1}$ \\
  $\sigma_X$      & Statistische Unsicherheit von $X$ & wie $X$ \\
  $\Phi$          & Phase einer Mode-Auslenkung (intern) & rad \\
  $\omega$        & Mode-Frequenz & cm$^{-1}$ oder rad/s \\
  $\omega_i$      & Frequenz der Mode $i$ & cm$^{-1}$ \\
  $\mathbf{\Omega}$ & Vektor aller Mode-Frequenzen & cm$^{-1}$ \\
  \bottomrule
\end{tabular}

\subsection{Sonderzeichen und Abk\"urzungen}

\begin{tabular}{p{2cm} p{8.5cm}}
  \toprule
  \textbf{Symbol} & \textbf{Bedeutung} \\
  \midrule
  $\angstrom$    & \AA{}ngstrom-Einheit \\
  $\langle\cdot\rangle$ & Quantenmechanischer/thermischer Erwartungswert \\
  $|\cdot|$      & Betrag (Skalar) bzw. Norm (Vektor) \\
  $\|\cdot\|$    & L2-Norm eines Vektors \\
  $\hat{\cdot}$  & Einheitsvektor (z.\,B.\ $\hat{\mathbf{n}}$) \\
  $\cdot^*$      & Komplex konjugiert \\
  $\sum_i$       & Summe \"uber Mode-Indizes (in Filterbereich) \\
  $\sum_a$       & Summe \"uber Atom-Indizes \\
  $\sum_X$       & Summe \"uber Kanal-Indizes (FeFe, FeN, FeS, NH, HA) \\
  RSS            & Root Sum of Squares: $\sqrt{\sum_i x_i^2}$ \\
  RMS            & Root Mean Square: $\sqrt{\frac{1}{n}\sum_i x_i^2}$ \\
  pDOS           & Projizierte Zustandsdichte (projected density of states) \\
  pCH            & Pyrrol-C-H (in NRVS-Literatur) \\
  HP             & High-Precision (Gaussian-Eigenvektoren) \\
  Std            & Standard-Genauigkeit (Gaussian-Eigenvektoren) \\
  PCET           & Proton-Coupled Electron Transfer \\
  ET             & Electron Transfer \\
  PT             & Proton Transfer \\
  NRVS           & Nuclear Resonance Vibrational Spectroscopy \\
  SCSD           & Symmetry-Coordinate Structural Decomposition \\
  OOP            & Out-of-Plane (senkrecht zur Cluster-Ebene) \\
  INP            & In-Plane (in der Cluster-Ebene) \\
  QHO            & Quantum-Mechanischer Harmonischer Oszillator \\
  WT             & Wild-Typ \\
  Mut            & Mutante \\
  prot           & Protoniert (Imidazol-NH vorhanden) \\
  deprot         & Deprotoniert (Imidazol-NH abgespalten) \\
  D$_{2h}$       & Punktgruppe D$_{2h}$ (rhombischer Cluster) \\
  $A_g$, $B_{1g}$, $\ldots$ & Irreps der D$_{2h}$-Symmetrie \\
  \bottomrule
\end{tabular}

% -------------------------------------------------------------------------
\section{Glossar}
\label{app:glossary}

Begriffe und Konzepte, die in der wissenschaftlichen Literatur und
diesem Handbuch verwendet werden, mit kurzer Definition. Sortiert
alphabetisch.

\paragraph{Akzeptor (H-Bridge).}
Atom mit freiem Elektronenpaar, das eine Wasserstoffbr\"ucke
empf\"angt. In [2Fe-2S]-Proteinen typisch: Carbonyl-Sauerstoff
(C=O) eines Backbones, Hydroxyl-Sauerstoff einer Tyr/Ser/Thr,
oder Wasser-Molek\"ule. Der "`Hauptakzeptor"' eines His-N-H ist
der mit max. Gauss-Gewicht (Distanz n\"aher zum Optimum
$r_0 = 2.8\,\angstrom$).

\paragraph{B-Faktor (Debye-Waller-Faktor).}
Mittlere quadratische Auslenkung eines Atoms im thermischen
Gleichgewicht: $B_a = 8\pi^2\langle u_a^2\rangle/3$. Einheit
$\angstrom^2$. Aus den Modes berechnet:
$B_a = 8\pi^2/3 \cdot \sum_i u_\text{rms}^2(i) |e_a(i)|^2$.

\paragraph{Bindungs-Kanal.}
Ein Bindungs-Kanal ist eine Klasse aller chemisch \"aquivalenten
Bindungen, deren Modulation gemeinsam aggregiert wird. F\"unf
Kan\"ale werden im Tool berechnet: FeFe (Cluster-Atmung), FeN
(Fe-N$_\text{His}$), FeS (Fe-S$_\text{Cys}$ + Fe-S$_\text{bridge}$),
NH (Imidazol-Stretching), HA (H-Akzeptor; Reaktionskoordinaten-Modus).

\paragraph{Cluster-Atmung.}
Symmetrische Mode, in der sich Fe-Fe-Distanz \"andert
(Streckung/Stauchung des Cluster-Rhombus). $A_g$-Symmetrie in
D$_{2h}$. Frequenz $\sim 200$\,cm$^{-1}$ in
typischen [2Fe-2S]-Proteinen.

\paragraph{Cluster-Normale $\hat{\mathbf{n}}$.}
Einheitsvektor senkrecht zur durch die vier Cluster-Atome (2 Fe + 2 S$_\mu$)
am besten gefitteten Ebene. Wird via SVD der zentrierten
Atom-Koordinaten berechnet. Definiert die OOP/INP-Trennung.

\paragraph{Co-Modulation.}
Gleichzeitige Modulation zweier Bindungs-Kan\"ale durch dieselbe
Mode. Indikator: $C_\text{PCET}(\omega) = (M_\text{NH}(\omega)
\cdot M_\text{HA}(\omega))^{1/2}$. Hohe Werte bei reaktions-
relevanten PCET-Modes.

\paragraph{Cys$_2$His$_2$, Cys$_3$His$_1$, Cys$_4$.}
Liganden-Schemata des [2Fe-2S]-Clusters:
\begin{itemize}\setlength{\itemsep}{0pt}
  \item Cys$_4$: alle vier Fe-koordinierenden Liganden sind
        Cysteinate (S$^-$). Standard-Ferredoxine.
  \item Cys$_2$His$_2$: zwei Cys + zwei His. Rieske-Typ. PCET m\"oglich.
  \item Cys$_3$His$_1$: drei Cys + ein His. mitoNEET-Typ.
\end{itemize}

\paragraph{D$_{2h}$-Irreps.}
Irreduzible Darstellungen der Punktgruppe D$_{2h}$ (rhombischer
Cluster mit drei orthogonalen $C_2$-Achsen + Inversionszentrum).
Acht Irreps: $A_g, B_{1g}, B_{2g}, B_{3g}, A_u, B_{1u}, B_{2u},
B_{3u}$. SCSD zerlegt jede Mode-Auslenkung in genau eine
Linearkombination dieser acht Irreps.

\paragraph{DFT (Density Functional Theory).}
Quantenchemische Methode zur Berechnung elektronischer Strukturen
und Schwingungsfrequenzen. Im Tool werden DFT-Frequenzen aus
Gaussian-16 oder ORCA-6 erwartet (typisch B3LYP/6-31G* mit
Konvergenz auf $10^{-8}$\,Hartree).

\paragraph{Eigenvektor.}
Massengewichtete Auslenkungs-Komponenten einer Schwingungsmode
relativ zur Gleichgewichtsgeometrie. F\"ur Atom $a$ in Mode $i$:
$\mathbf{e}_a^{(i)}$. Normierung in Gaussian: $\sum_a m_a
|\mathbf{e}_a^{(i)}|^2 = 1$ (massengewichtet).

\paragraph{Embedding.}
Niedrigdimensionale Darstellung hochdimensionaler Mode-Feature-
Vektoren. Im Tool: 31-dim. Feature-Matrix wird via UMAP auf 2D
abgebildet.

\paragraph{Feature-Matrix.}
$N_\text{modes} \times 31$-dim. Matrix mit pro Mode 31 Skalaren
(Frequenz, OOP\%, Cluster-Beteiligung, vier Liganden-OOP,
SCSD-Irrep-Anteile, etc.). Eingabe f\"ur das UMAP-Embedding.

\paragraph{Frequenz-Fenster (Window).}
Frequenz-Bereich \texttt{(lo, hi)} in cm$^{-1}$, der separat
ausgewertet werden soll. Multi-Window-Modus: mehrere Fenster
parallel, jeweils mit eigenen Reorg-Aggregaten und
Embeddings.

\paragraph{Gauss-Eigenvektor (HP / Std).}
Gaussian schreibt Schwingungsmoden-Eigenvektoren in zwei Bl\"ocken:
\begin{itemize}\setlength{\itemsep}{0pt}
  \item \textbf{Standard (Std)}: 4 Dezimalstellen, kompakt, immer
        vorhanden.
  \item \textbf{High Precision (HP)}: 5 Dezimalstellen, durch
        \texttt{freq=hpmodes} aktiviert; ben\"otigt mehr Speicher
        in der Logdatei. Empfohlen f\"ur dieses Tool.
\end{itemize}

\paragraph{HDBSCAN.}
Hierarchisches Density-Based Spatial Clustering of Applications
with Noise. Cluster-Algorithmus, der UMAP-Punkte zu Clustern
zusammenfasst und Rauschen separat markiert.

\paragraph{HPmodes.}
Gaussian-Schl\"usselwort (\texttt{freq=hpmodes}), das die HP-
Eigenvektoren zus\"atzlich in die Logdatei schreibt.

\paragraph{Huang-Rhys-Faktor.}
$S = \frac{\mu\omega(\Delta q)^2}{2\hbar}$. Dimensionslose
Verschiebung einer harmonischen Mode zwischen elektronischen
Zust\"anden, in Einheiten der Nullpunkts-Auslenkung. Verwandt mit
unserem $\lambda_X$, aber in einer anderen Konvention
(Verschiebungs- vs.\ Auslenkungs-Quadrat).

\paragraph{Imidazol.}
Aromatischer 5-Ring mit zwei Stickstoffen (N$_\delta$ und N$_\epsilon$)
in Histidin. Eines der Stickstoffe (typisch N$_\epsilon$) bindet
das Fe; der andere kann je nach Tautomer protoniert oder nicht
sein.

\paragraph{INP (In-Plane).}
Mode-Auslenkung \"uberwiegend in der Cluster-Ebene (senkrecht zu
$\hat{\mathbf{n}}$). Klassifiziert mit OOP\% < 40 (binaere
Schwelle bei 60).

\paragraph{Kabsch-Alignment.}
Algorithmus zum optimalen Aufeinanderlegen zweier
Atomkoordinaten-S\"atze (PDB $\leftrightarrow$ Gaussian).
Liefert Rotations-Matrix und RMSD.

\paragraph{Kern-Klassifikation.}
Heuristische Mode-Typ-Zuweisung: Cluster-Atmung,
Cluster-Hinge, Cluster-Twist, Liganden-Stretch, etc.
Implementiert in \texttt{core.py:classify\_kern}. Nicht orthogonal
mit SCSD-Irreps; mehrere Kern-Typen koennen demselben Irrep
entsprechen.

\paragraph{Lamb-M\"o{\ss}bauer-Faktor.}
$f_\text{LM} = \exp(-2W)$, mit $W$ dem Debye-Waller-Faktor des
M\"o{\ss}bauer-aktiven Kerns ($^{57}$Fe). Misst, welcher Anteil
der Resonanzabsorption ohne Phononen-Anregung erfolgt
(elastischer Anteil im NRVS-Spektrum).

\paragraph{Marcus-Hush-Theorie.}
Klassische Theorie des Elektronentransfers (Marcus 1956,
Hush 1958), die die Reaktionsrate \"uber eine harmonische
Approximation der Reaktanten- und Produkt-Energiefl\"achen mit
einer einzigen Reorganisations-Koordinate zusammenfasst.

\paragraph{mitoNEET.}
Cys$_3$His$_1$-koordiniertes [2Fe-2S]-Protein. Bekannt f\"ur
seine spektroskopisch besonders gut untersuchten NRVS-Bandenstruktur
(Gee 2021, Bergner 2017).

\paragraph{Modulations-Spektrum $M_X(\omega)$.}
Gauss-verbreitertes Spektrum der Bindungs-Modulationen pro Kanal:
$M_X(\omega) = \sum_i |\Delta r_X(i)| \cdot G(\omega - \omega_i,
\sigma)$. Zeigt, in welchem Frequenzbereich Bindung $X$ stark
moduliert wird.

\paragraph{NRVS (Nuclear Resonance Vibrational Spectroscopy).}
Synchrotron-Spektroskopie an M\"o{\ss}bauer-aktiven Kernen
(typisch $^{57}$Fe). Misst $^{57}$Fe-projizierte
Vibrations-Zustandsdichte. Komplement\"ar zu Raman/IR
(elementselektiv, nicht symmetrieselektiv). Im Tool wird
nicht das NRVS-Spektrum selbst berechnet, sondern die
Reorg-Beitr\"age, die zur Bandenidentifikation genutzt
werden koennen.

\paragraph{OOP (Out-of-Plane).}
Mode-Auslenkung \"uberwiegend senkrecht zur Cluster-Ebene
(parallel zu $\hat{\mathbf{n}}$). Klassifiziert mit OOP\% > 60.

\paragraph{OOP\%.}
Quadrat-Anteil der Mode-Auslenkung in $\hat{\mathbf{n}}$-Richtung,
relativ zur Gesamt-Auslenkung. $\text{OOP\%} = 100 \cdot \sum_a
(\mathbf{e}_a \cdot \hat{\mathbf{n}})^2 / \sum_a |\mathbf{e}_a|^2$.

\paragraph{PCET (Proton-Coupled Electron Transfer).}
Konzertierter Transfer eines Elektrons und eines Protons. In
[2Fe-2S]-Proteinen typischerweise: Reduktion des Clusters
($\text{Fe}^{3+} + e^- \to \text{Fe}^{2+}$) gekoppelt mit
Imidazol-NH-Deprotonierung (Rieske: $pK_a$ verschiebt um 4-5
Einheiten zwischen oxidiertem und reduziertem Zustand).

\paragraph{Pseudocode.}
Algorithmus-Beschreibung in einer informellen, aber pr\"azisen
Notation ohne spezifische Programmiersprache. Wird im
Anhang \ref{app:pseudocodes} f\"ur die zentralen Algorithmen
verwendet.

\paragraph{Reaktionskoordinaten-Modus (HA).}
Definition der HA-Modulation: Anstelle der Modulation der
H-Akzeptor-Distanz wird die Bewegung des H-Atoms relativ zum
Donor-N entlang der N$\to$Akzeptor-Achse berechnet. Klassisches
Schreibweise: $\Delta r_\text{HA} = (\mathbf{e}_H - \mathbf{e}_N)
\cdot \hat{\mathbf{n}}_{N\to\text{Akz}}$. Misst echte
Protonenbewegung statt unspezifischer Skelett-Modes.

\paragraph{Reorganisations-Energie $\lambda_X$.}
Im Tool: thermische Bindungs-L\"angen-Fluktuations-Energie pro
Mode entlang Bindung $X$, berechnet in Marcus-Hush-Form mit
thermisch skalierten Mode-Auslenkungen. Verwandt mit, aber nicht
identisch zur klassischen Marcus-Hush-Reorganisationsenergie
(siehe \S\ref{sec:reorg}).

\paragraph{Rieske-Protein.}
Cys$_2$His$_2$-koordiniertes [2Fe-2S]-Protein. Komponente in
Cytochrom-bc$_1$-Komplexen, beteiligt am Q-Cycle. Klassisches
PCET-System (Imidazol-NH-Deprotonierung gekoppelt an
Cluster-Reduktion).

\paragraph{RMS (Root Mean Square).}
$\sqrt{\frac{1}{n}\sum_i x_i^2}$. Mittelwert der Quadrate, dann
Wurzel.

\paragraph{RSS (Root Sum of Squares).}
$\sqrt{\sum_i x_i^2}$. Summe der Quadrate, dann Wurzel.
\textbf{Nicht} mit RMS verwechseln (durch $n$ geteilt).

\paragraph{SCSD (Symmetry-Coordinate Structural Decomposition).}
Methode von Kingsbury \& Senge (2024), die strukturelle
Verzerrungen in einer beliebigen Geometrie als Linearkombination
von Symmetriekoordinaten einer kanonischen Referenzgeometrie
darstellt. Im Tool: D$_{2h}$-Zerlegung mit kanonischer
Referenz-[2Fe-2S]-Cluster (Fe-Fe = 2.73\,$\angstrom$,
Fe-S = 2.20\,$\angstrom$).

\paragraph{Sekund\"arstruktur (SSE).}
Lokale 3D-Struktur eines Proteins (Helix, Beta-Faltblatt, Loop).
Im Tool aus dem PDB-File abgeleitet (HELIX/SHEET-Records oder
DSSP-Fallback oder $\varphi/\psi$-Heuristik).

\paragraph{TOML.}
Konfigurationsformat (\emph{Tom's Obvious, Minimal Language}).
Im Tool: \texttt{config.toml} oder \texttt{run.toml} f\"ur
reproduzierbare L\"aufe.

\paragraph{Thermisch skalierte Eigenvektoren.}
$\mathbf{e}_a^\text{thermal}(i) = \mathbf{e}_a^\text{normalized}(i)
\cdot u_\text{rms}(i)$. Eigenvektoren multipliziert mit der
Mode-spezifischen RMS-Amplitude (QHO-Erwartungswert bei der
gegebenen Temperatur). Direkt verwendbar zur Berechnung
thermischer Bindungs-Modulationen.

\paragraph{UMAP (Uniform Manifold Approximation and Projection).}
Topologie-erhaltende Dimensions-Reduktions-Methode (McInnes 2018).
Lokale + globale Strukturen werden gleichzeitig erhalten,
deterministisch reproduzierbar mit Seed.

\paragraph{Wasserstoffbr\"ucke (H-Bridge).}
Schwache, aber gerichtete Wechselwirkung zwischen einem
Proton-Donor (X-H mit X = N, O) und einem Akzeptor (Y mit
freiem Elektronenpaar, Y = O, N). Energie typisch 5-30
kJ/mol. In Proteinen verantwortlich fuer Sekund\"arstruktur und
PCET-Reaktionskan\"ale.

% -------------------------------------------------------------------------
\section{Algorithmus-Pseudocodes}
\label{app:pseudocodes}

Die zentralen Algorithmen des Tools werden hier als Pseudocode
dokumentiert. Sprachebene: zwischen formaler Mathematik und Python --
keine Typdetails, aber alle nicht-trivialen Kontrollfl\"usse
explizit. Variablen-Namen folgen den Modulnamen im Code.

Die Pseudocodes sind \emph{nicht} 1:1-Implementierungen, sondern
Skizzen der Logik. Der echte Code in \texttt{src/modenanalyse\_2fe2s/}
hat zus\"atzlich Performance-Optimierungen, Caching, Logging und
Fehlerbehandlung, die hier zur Lesbarkeit weggelassen sind.

\subsection{Cluster-Erkennung}
\label{app:pseudo-cluster}

Zentrale Aufgabe: aus einer Liste von Atomen finde alle [2Fe-2S]-
Cluster. Bei Multi-Cluster-Systemen (Glutaredoxin-Dimere etc.) m\"ussen
mehrere Cluster eindeutig gruppiert und sortiert werden.

\paragraph{Algorithmus.}

\begin{algorithm}[H]
\caption{\textsc{FindAllClusters}}
\begin{algorithmic}[1]
\Require Atomliste $A = \{a_1, \ldots, a_N\}$ mit Element + Kartesischer Position;
         Cutoffs $c_\text{FeFe}, c_\text{FeS}$ in $\angstrom$
\Ensure Liste von Clustern $\mathcal{C}_1, \mathcal{C}_2, \ldots$,
        sortiert nach Fe-Fe-Distanz aufsteigend
\State $\text{Fe} \gets \{a \in A : \text{element}(a) = \text{Fe}\}$
\State $\text{S}  \gets \{a \in A : \text{element}(a) = \text{S}\}$
\If{$|\text{Fe}| < 2$ \textbf{oder} $|\text{S}| < 2$}
  \State \Return $\emptyset$  \Comment{Kein Cluster m\"oglich}
\EndIf
\State Bilde alle Fe-Paare $(f_i, f_j)$ mit
       $\|\mathbf{r}(f_i) - \mathbf{r}(f_j)\| \leq c_\text{FeFe}$
\State Sortiere die Fe-Paare nach Distanz aufsteigend
\State $\text{verwendet} \gets \emptyset$  \Comment{Set bereits zugeordneter Atome}
\State $\text{cluster\_list} \gets \emptyset$
\ForAll{Fe-Paare $(f_i, f_j)$ in sortierter Reihenfolge}
  \If{$f_i \in \text{verwendet}$ \textbf{oder} $f_j \in \text{verwendet}$}
    \State \textbf{weiter}  \Comment{Atom schon einem Cluster zugeordnet}
  \EndIf
  \State $\mathcal{S}_\mu \gets \{s \in \text{S} :$ S \"uberbr\"uckt $f_i, f_j$
         mit $\|\mathbf{r}(s) - \mathbf{r}(f_i)\| \leq c_\text{FeS}$
         \textbf{und} $\|\mathbf{r}(s) - \mathbf{r}(f_j)\| \leq c_\text{FeS}\}$
  \If{$|\mathcal{S}_\mu| \geq 2$}  \Comment{Genug Br\"uckenschwefel?}
    \State $(s_1, s_2) \gets$ die zwei n\"achsten S aus $\mathcal{S}_\mu$
    \State $\text{cluster\_list} \gets \text{cluster\_list} \cup
           \{\{f_i, f_j, s_1, s_2\}\}$
    \State $\text{verwendet} \gets \text{verwendet} \cup \{f_i, f_j, s_1, s_2\}$
  \EndIf
\EndFor
\State \Return $\text{cluster\_list}$
\end{algorithmic}
\end{algorithm}

\paragraph{Subtilit\"aten in der Implementation.}
\begin{itemize}\setlength{\itemsep}{0pt}
  \item Die Sortierung der Fe-Paare nach Distanz ist
        \emph{deterministisch} und stellt sicher, dass das engste
        Cluster zuerst belegt wird. Damit ist
        \texttt{cluster\_index = 0} immer das engste, ohne dass der
        Nutzer die Atom-Reihenfolge im Logfile ber\"ucksichtigen muss.
  \item Die Cutoffs $c_\text{FeFe}$ und $c_\text{FeS}$ sind
        Config-Parameter (Default $3.5$ und $3.0\,\angstrom$). F\"ur
        sehr ausgedehnte Cluster (z.\,B.\ [4Fe-4S] oder Mo-Fe-Cofaktor)
        m\"ussen diese ggf.\ erh\"oht werden -- wobei das Tool dann nur
        noch das [2Fe-2S]-Substrukturteilst\"uck liest.
  \item Wird ein Atom durch ein engeres Cluster bereits "`verbraucht"',
        so wird es nicht nochmal in einem weiteren Cluster verwendet.
        Damit findet der Algorithmus auch in eng gepackten
        Multi-Cluster-Systemen die wahre Anzahl von Clustern.
\end{itemize}

\subsection{Cluster-Normale via SVD}
\label{app:pseudo-svd}

Aus den vier Cluster-Atomen wird per Singul\"arwertzerlegung die
Best-Fit-Ebene und damit die Cluster-Normale $\hat{\mathbf{n}}$ bestimmt.

\begin{algorithm}[H]
\caption{\textsc{ClusterNormale}}
\begin{algorithmic}[1]
\Require Cluster-Atome $\mathcal{C} = \{a_1, a_2, a_3, a_4\}$
        mit Positionen $\mathbf{r}_1, \ldots, \mathbf{r}_4 \in \mathbb{R}^3$
\Ensure Einheitsvektor $\hat{\mathbf{n}}$ senkrecht zur Best-Fit-Ebene
\State $\bar{\mathbf{r}} \gets \frac{1}{4} \sum_{i=1}^{4} \mathbf{r}_i$
       \Comment{Schwerpunkt}
\State Konstruiere zentrierte Matrix
       $\tilde{\mathbf{R}} \in \mathbb{R}^{4\times 3}$:
       $\tilde{\mathbf{R}}_{i,:} = \mathbf{r}_i - \bar{\mathbf{r}}$
\State $\mathbf{U}, \boldsymbol{\Sigma}, \mathbf{V}^\top \gets
       \textsc{SVD}(\tilde{\mathbf{R}})$
       \Comment{$\mathbf{V}$ ist orthogonale $3\times 3$-Matrix}
\State $\hat{\mathbf{n}} \gets \mathbf{V}_{:,3}$
       \Comment{dritte Spalte = Richtung kleinster Varianz}
\If{$\hat{\mathbf{n}}_z < 0$}  \Comment{Vorzeichen-Konvention}
  \State $\hat{\mathbf{n}} \gets -\hat{\mathbf{n}}$
\EndIf
\State \textbf{return} $\hat{\mathbf{n}}$
\end{algorithmic}
\end{algorithm}

Bemerkung: das Faltungs-Residual ist $\Sigma_{33}$ (kleinster Singul\"arwert).
Werte $> 0.1\,\angstrom$ deuten auf einen stark nicht-planaren Cluster hin --
ungew\"ohnlich f\"ur native [2Fe-2S]-Cluster, aber m\"oglich bei stark
verzerrten Mutanten.

\subsection{Kabsch-Alignment (PDB $\leftrightarrow$ Gauss)}
\label{app:pseudo-kabsch}

Wenn ein PDB-File mitgegeben wird, m\"ussen die PDB-Koordinaten auf
die Gaussian-Koordinaten gemappt werden -- die Geometrien sind in der
Regel rotationsverschoben. Der Kabsch-Algorithmus
\cite{Kabsch1976} liefert die optimale Rotation.

\begin{algorithm}[H]
\caption{\textsc{KabschAlign}}
\begin{algorithmic}[1]
\Require Punktmengen $\mathbf{P}, \mathbf{Q} \in \mathbb{R}^{n\times 3}$,
         korrespondierend (gleicher Atom-Reihenfolge)
\Ensure Rotation $\mathbf{R}$, Translation $\mathbf{t}$, RMSD
\State $\bar{\mathbf{p}} \gets \frac{1}{n}\sum_i \mathbf{P}_{i,:}$,
       $\bar{\mathbf{q}} \gets \frac{1}{n}\sum_i \mathbf{Q}_{i,:}$
\State $\mathbf{P}' \gets \mathbf{P} - \bar{\mathbf{p}}$,
       $\mathbf{Q}' \gets \mathbf{Q} - \bar{\mathbf{q}}$
       \Comment{Schwerpunkt jeweils auf Ursprung}
\State $\mathbf{H} \gets \mathbf{P}'^\top \mathbf{Q}'$
       \Comment{$3\times 3$-Kreuzkorrelation}
\State $\mathbf{U}, \boldsymbol{\Sigma}, \mathbf{V}^\top \gets
       \textsc{SVD}(\mathbf{H})$
\State $d \gets \text{sign}(\det(\mathbf{V} \mathbf{U}^\top))$
       \Comment{Reflektions-Vorzeichen}
\State $\mathbf{D} \gets \text{diag}(1, 1, d)$
\State $\mathbf{R} \gets \mathbf{V} \mathbf{D} \mathbf{U}^\top$
       \Comment{Optimale Rotation}
\State $\mathbf{t} \gets \bar{\mathbf{q}} - \mathbf{R} \bar{\mathbf{p}}$
\State $\text{RMSD} \gets \sqrt{\frac{1}{n} \|\mathbf{R}\mathbf{P}' -
       \mathbf{Q}'\|_F^2}$
\State \textbf{return} $(\mathbf{R}, \mathbf{t}, \text{RMSD})$
\end{algorithmic}
\end{algorithm}

\paragraph{Was die $\mathbf{D}$-Matrix tut.}
Die "`naive"' L\"osung $\mathbf{R} = \mathbf{V}\mathbf{U}^\top$ kann
eine Reflektion sein (Determinante $-1$), wenn die zwei Punktmengen
chiral verschieden sind (z.\,B.\ versehentlich gespiegelt). Die
$\mathbf{D}$-Matrix korrigiert das Vorzeichen der dritten Singul\"arvektor-
Achse, sodass eine echte Rotation (Determinante $+1$) zur\"uckgegeben
wird. Bei korrekt korrespondierenden Punktmengen ist $d = +1$, also
$\mathbf{D} = \mathbf{I}$.

\subsection{Mode-Loop mit thermischer Skalierung}
\label{app:pseudo-mode-loop}

Der zentrale Loop des Tools: f\"ur jede Schwingungsmode wird der
Eigenvektor gelesen, thermisch skaliert, und in die OOP/INP-,
Reorg- und SCSD-Pipelines gespeist.

\begin{algorithm}[H]
\caption{\textsc{AnalyzeMode}}
\begin{algorithmic}[1]
\Require Mode-Index $i$, Frequenz $\omega_i$, reduzierte Masse $\mu_i$,
         Temperatur $T$, Eigenvektor $\mathbf{e}^{(i)}$ (mass-gewichtet,
         normalisiert), Cluster-Atome $\mathcal{C}$,
         Liganden $\mathcal{L}$
\Ensure Result-Dict $r$ mit OOP\%, $\Delta r_X$, $\lambda_X$,
        SCSD-Anteilen, Kern-Klassifikation
\State \Comment{1. Thermische RMS-Amplitude (Quanten-Erwartungswert)}
\State $u_\text{rms}(i) \gets \sqrt{\frac{\hbar}{2\mu_i\omega_i}
       \coth\!\left(\frac{\hbar\omega_i}{2 k_B T}\right)}$
\State $\mathbf{e}^\text{therm}_a \gets \mathbf{e}_a^{(i)} \cdot u_\text{rms}(i)
       \quad \forall a \in \mathcal{C} \cup \mathcal{L}$
\State \Comment{2. OOP/INP-Klassifikation}
\State $\mathbf{e}_a^\text{OOP} \gets (\mathbf{e}^\text{therm}_a \cdot
       \hat{\mathbf{n}})\,\hat{\mathbf{n}}$
\State $\alpha_\text{OOP} \gets 100 \cdot
       \frac{\sum_{a\in\mathcal{C}} \|\mathbf{e}_a^\text{OOP}\|^2}
            {\sum_{a\in\mathcal{C}} \|\mathbf{e}^\text{therm}_a\|^2}$
\State $r.\text{oop\_pct} \gets \alpha_\text{OOP}$
\State $r.\text{type} \gets \textsc{ClassifyOOP}(\alpha_\text{OOP},
       \theta_\text{OOP}=60)$
\State \Comment{3. Bindungs-Modulationen}
\ForAll{Bindungs-Kan\"ale $X \in \{\text{FeFe, FeN, FeS, NH, HA}\}$}
  \State $\Delta r_X(i) \gets \textsc{ComputeBondModulation}(X,
         \mathbf{e}^\text{therm}, \text{Geometrie})$
       \Comment{siehe \S\ref{sec:reorg}}
  \State $\lambda_X(i) \gets \tfrac{1}{2}\mu_i \omega_i^2 \,(\Delta r_X(i))^2$
\EndFor
\State \Comment{4. SCSD-Symmetriezerlegung (optional)}
\If{\texttt{cfg.analyze\_scsd}}
  \State $\textsc{Verschiebung} \gets \mathbf{e}^\text{therm}|_\mathcal{C}$
       \Comment{nur Cluster-Atome}
  \State $\textsc{Irrep-Anteile} \gets \textsc{ProjektOnD2hIrreps}(
         \textsc{Verschiebung}, \text{Referenzgeometrie})$
  \State $r.\text{scsd} \gets \textsc{Irrep-Anteile}$
\EndIf
\State \Comment{5. Heuristische Kern-Klassifikation}
\State $r.\text{kern\_type} \gets \textsc{ClassifyKern}(\mathbf{e}^\text{therm},
       \mathcal{C}, \alpha_\text{OOP})$
\State \textbf{return} $r$
\end{algorithmic}
\end{algorithm}

\paragraph{Subtilit\"aten:}
\begin{itemize}\setlength{\itemsep}{0pt}
  \item Bei \texttt{temp\_k = None} entf\"allt die thermische Skalierung;
        $u_\text{rms}$ wird durch \texttt{cfg.amplitude} (Default 1.0)
        ersetzt. Das ist der "`klassische"' Modus, in dem alle Modes
        gleich gewichtet sind.
  \item Bei niedriger Frequenz und hoher Temperatur wird das
        Coth-Argument klein und $u_\text{rms}$ gro{\ss} -- f\"ur Modes
        unter $\sim 30\,$cm$^{-1}$ kann das in einem typischen
        Filter-Bereich Anteile $> 50\,\%$ der Total-Reorg ergeben.
        Daher die Empfehlung: \texttt{freq\_min} immer auf 50 oder 100
        setzen, wenn Niederfrequenz-Skelett-Bewegungen unterdr\"uckt
        werden sollen.
  \item Im "`no-His"'-Pfad (kein Histidin im Cluster) werden die
        Kan\"ale NH und HA \emph{vor} der Reorg-Schleife auf null
        gesetzt -- nicht erst durch null-$\Delta r$, sondern bereits
        durch den \"ubergeordneten Aufruf in \texttt{runner.py}.
        Das spart Geometrie-Berechnungen und macht den Pfad
        deterministisch sauber.
\end{itemize}

\subsection{Reorg-Aggregation pro Mode}
\label{app:pseudo-reorg}

Pro Mode werden die Bindungs-Modulationen $\Delta r_X$ aus den
thermisch skalierten Eigenvektoren berechnet.

\begin{algorithm}[H]
\caption{\textsc{ComputeBondModulation} (FeFe-Beispiel)}
\begin{algorithmic}[1]
\Require Bindungs-Atome $a, b$ (mit Positionen + thermisch skalierten
         Auslenkungen $\mathbf{e}_a, \mathbf{e}_b$)
\Ensure Bindungs-Modulation $\Delta r$ (vorzeichenbehaftet)
\State $\mathbf{r}_a^\text{neu} \gets \mathbf{r}_a + \mathbf{e}_a$
\State $\mathbf{r}_b^\text{neu} \gets \mathbf{r}_b + \mathbf{e}_b$
\State $r_\text{neu} \gets \|\mathbf{r}_a^\text{neu} - \mathbf{r}_b^\text{neu}\|$
\State $r_\text{eq} \gets \|\mathbf{r}_a - \mathbf{r}_b\|$
\State $\Delta r \gets r_\text{neu} - r_\text{eq}$
\State \textbf{return} $\Delta r$
\end{algorithmic}
\end{algorithm}

F\"ur Multi-Bond-Kan\"ale (z.\,B.\ FeS mit vier Fe-S-Bindungen) wird
\textsc{ComputeBondModulation} pro Bindung aufgerufen, und die
Beitr\"age werden \emph{vorzeichenfrei} per RSS aggregiert
(\S\ref{app:notation}):
\begin{equation}
  \Delta r_X^\text{rss} = \sqrt{\sum_{(a,b) \in X} (\Delta r_{(a,b)})^2}
\end{equation}
Damit ist der Kanal-Beitrag $\lambda_X = \tfrac{1}{2}\mu\omega^2
(\Delta r_X^\text{rss})^2$ identisch mit der Summe der
Pair-Beitr\"age:
\begin{equation}
  \lambda_X = \sum_{(a,b) \in X} \tfrac{1}{2}\mu\omega^2 (\Delta r_{(a,b)})^2.
  \label{eq:reorg-additivity}
\end{equation}
Diese Marcus-Hush-Additivit\"at ist die zentrale Konsistenz-Bedingung
und Smoke-Test (\S\ref{sec:workflow-worked-example}).

\subsection{HA-Reaktionskoordinaten-Berechnung}
\label{app:pseudo-ha}

Der HA-Kanal misst die H-Akzeptor-Distanzmodulation entlang der
Reaktionskoordinate (statt der naiven H-Akz-Distanz).

\begin{algorithm}[H]
\caption{\textsc{ComputeHAReactionCoord}}
\begin{algorithmic}[1]
\Require N$_\text{Donor}$, H, Akzeptor $A$ (Positionen + Auslenkungen)
\Ensure Reaktions-Koordinaten-Modulation $\Delta r_\text{HA}$
\State $\hat{\mathbf{n}}_{N\to A} \gets
       \frac{\mathbf{r}_A - \mathbf{r}_N}
            {\|\mathbf{r}_A - \mathbf{r}_N\|}$
       \Comment{Donor-Akzeptor-Achse, Einheitsvektor}
\State $\Delta\mathbf{e} \gets \mathbf{e}_H - \mathbf{e}_N$
       \Comment{H-Bewegung relativ zum Donor}
\State $\Delta r_\text{HA} \gets \Delta\mathbf{e} \cdot
       \hat{\mathbf{n}}_{N\to A}$
       \Comment{Projektion auf Reaktionskoordinate}
\State \textbf{return} $\Delta r_\text{HA}$
\end{algorithmic}
\end{algorithm}

\paragraph{Warum dieser Ansatz?}
Die naive H-Akz-Distanz misst die \emph{Gesamtbewegung} des H-Atoms
relativ zum Akzeptor-Atom. Damit zeigen reine Skelett-Modes (in
denen der ganze Komplex sich kollektiv bewegt) ein gro{\ss}es
$\Delta r$, obwohl das Proton sich \emph{relativ} zum Donor gar nicht
bewegt -- sie sind keine echten PCET-relevanten Modes. Die
Reaktionskoordinaten-Konvention extrahiert nur die echte
Protonen-Bewegung.

\subsection{Modulations-Spektrum und kumulative Reorg-Kurve}
\label{app:pseudo-spektren}

Aus den pro-Mode-Werten werden frequenzaufgel\"oste Spektren gebaut.

\begin{algorithm}[H]
\caption{\textsc{ModulationsSpectrum}}
\begin{algorithmic}[1]
\Require Mode-Liste mit $(\omega_i, \Delta r_X(i))$, Verbreiterung $\sigma$,
         Frequenzraster $\boldsymbol{\omega}$
\Ensure Modulations-Spektrum $M_X(\omega)$
\State $M_X(\omega) \gets 0 \quad \forall \omega \in \boldsymbol{\omega}$
\ForAll{Modes $i$}
  \State $g_i(\omega) \gets \frac{1}{\sigma\sqrt{2\pi}}
         \exp\!\left(-\frac{(\omega - \omega_i)^2}{2\sigma^2}\right)$
  \State $M_X(\omega) \gets M_X(\omega) + |\Delta r_X(i)| \cdot g_i(\omega)$
\EndFor
\State \textbf{return} $M_X(\boldsymbol{\omega})$
\end{algorithmic}
\end{algorithm}

\begin{algorithm}[H]
\caption{\textsc{KumulativeReorg}}
\begin{algorithmic}[1]
\Require Mode-Liste mit $(\omega_i, \lambda_X(i))$, sortiert aufsteigend nach $\omega$
\Ensure Kumulative Kurve $\Lambda_X(\omega)$
\State $\Lambda \gets 0$
\ForAll{Modes $i$ (sortiert aufsteigend)}
  \State Schreibe $(\omega_i^-, \Lambda)$ und $(\omega_i^+, \Lambda + \lambda_X(i))$
         \Comment{Stufenfunktion: Sprung an $\omega_i$}
  \State $\Lambda \gets \Lambda + \lambda_X(i)$
\EndFor
\State Schreibe Endwert $(\omega_\text{max}, \Lambda)$
\State \textbf{return} $\Lambda_X(\omega)$
\end{algorithmic}
\end{algorithm}

Beide Spektren sind im Sheet \texttt{Modulations\_Spektren} bzw.\
\texttt{Lambda\_kumulativ} ausgegeben (\S\ref{sec:ausgabe}).

\subsection{HDBSCAN-Clusterung des Embeddings}
\label{app:pseudo-hdbscan}

Nach der UMAP-Projektion auf 2D werden die Punkte mit HDBSCAN
\cite{Campello2013} dichtebasiert geclustert.

\begin{algorithm}[H]
\caption{\textsc{HDBSCANCluster}}
\begin{algorithmic}[1]
\Require 2D-Embedding-Punkte $\{x_1, \ldots, x_N\}$,
         min\_cluster\_size, min\_samples
\Ensure Cluster-Labels $\ell_1, \ldots, \ell_N \in \{-1, 0, 1, \ldots\}$
       (-1 = Rauschen)
\State \Comment{1. Gewichteter Graph: jeder Punkt hat Distanz zu seinen K n\"achsten Nachbarn}
\State $\text{core\_dist}(x_i) \gets$ Distanz zum \texttt{min\_samples}-ten Nachbarn
\State $\text{mutual\_reachability}(x_i, x_j) \gets$
       $\max(\text{core\_dist}(x_i), \text{core\_dist}(x_j),
       \|x_i - x_j\|)$
\State \Comment{2. Minimum-Spanning-Tree der mutual\_reachability-Matrix}
\State $\text{MST} \gets \textsc{KruskalMST}(\text{mutual\_reachability})$
\State \Comment{3. Hierarchischer Cluster-Baum}
\State $\text{Tree} \gets \textsc{ConvertMSTtoHierarchy}(\text{MST})$
\State \Comment{4. Stabile Cluster ausw\"ahlen (Excess of Mass)}
\State $\text{cluster\_list} \gets \textsc{ExcessOfMassSelection}(\text{Tree},
       \text{min\_cluster\_size})$
\State \Comment{5. Punkte den Clustern zuweisen oder als Rauschen markieren}
\ForAll{Punkte $x_i$}
  \If{$x_i$ ist Mitglied eines stabilen Clusters $C_k$}
    \State $\ell_i \gets k$
  \Else
    \State $\ell_i \gets -1$  \Comment{Rauschen}
  \EndIf
\EndFor
\State \textbf{return} $\boldsymbol{\ell}$
\end{algorithmic}
\end{algorithm}

\paragraph{Warum HDBSCAN statt DBSCAN?}
\begin{itemize}\setlength{\itemsep}{0pt}
  \item HDBSCAN braucht keine globale $\epsilon$-Schwelle (DBSCAN
        braucht eine -- die ist bei verschiedener Cluster-Dichte
        problematisch).
  \item HDBSCAN findet Cluster verschiedener Dichte automatisch (durch
        die hierarchische Konstruktion).
  \item HDBSCAN ist robust gegen Outlier (sie landen automatisch im
        Rauschen, ohne dass eine Schwelle angepasst werden muss).
\end{itemize}

In der Standard-Konfiguration verwendet das Tool
\texttt{min\_cluster\_size = 8} und \texttt{min\_samples = 5}; siehe
\S\ref{sec:embeddings} f\"ur die Begr\"undung.


% -------------------------------------------------------------------------
\section{Excel-Sheet-Vollreferenz}
\label{app:sheet-referenz}

Pro Lauf entstehen bis zu \emph{vier} Excel-Workbooks. Dieser Anhang
dokumentiert jede Spalte jedes Sheets, sodass Tabellenwerte ohne
R\"uckgriff auf den Code interpretierbar sind. Sheet-Namen, die in
\texttt{export.py} per \texttt{wb.create\_sheet(...)} erzeugt werden,
sind hier in der gleichen Reihenfolge gelistet.

\subsection{Hauptauswertungs-Workbook (\texttt{*\_analysis.xlsx})}
\label{app:wb-hauptauswertung}

Enth\"alt 16 Sheets mit allen modenbezogenen Auswertungen, Cluster-
Geometrie und SCSD-Symmetriezerlegung. Pflicht-Sheet f\"ur jede
Veroeffentlichung.

\subsubsection*{\texttt{Modenanalyse} -- Mode-by-mode Detail-Tabelle}

Eine Zeile pro Schwingungsmode (typisch 60--800 Modes je nach
\texttt{freq\_max}-Filter). 32 Spalten:

\begin{center}
\small
\begin{tabular}{p{3.6cm}p{1.4cm}p{8.4cm}}
\toprule
\textbf{Spalte} & \textbf{Einheit} & \textbf{Bedeutung} \\
\midrule
\texttt{Modus \#}        & --       & Mode-Index aus dem Logfile (1-basiert) \\
\texttt{Frequenz}        & cm$^{-1}$ & Frequenz $\omega_i$ \\
\texttt{u\_rms}          & $\angstrom$ & Thermische RMS-Amplitude
                                       $\sqrt{\hbar/(2\mu\omega)\coth(\hbar\omega/2k_BT)}$ \\
\texttt{Red.Masse}       & AMU      & Reduzierte Masse $\mu_i$ \\
\texttt{Frc.Const.}      & mdyn/$\angstrom$ & Kraftkonstante \\
\texttt{Sym.}            & --       & Mulliken-Symbol aus Gaussian (z.\,B.\ A, B$_g$) \\
\texttt{Typ}             & --       & OOP/INP-Klassifikation
                                       (bin\"ar, Schwelle 60\,\%) \\
\texttt{Typ (fein)}      & --       & 7-stufige OOP/INP-Klassifikation \\
\texttt{Praez.}          & --       & HP/Std-Eigenvektor-Quelle \\
\texttt{Lig OOP\%}       & \%       & OOP-Anteil der Mode auf
                                       Liganden-Sph\"are \\
\texttt{$\pm$sigma}      & \%       & Bootstrap-Streuung der Sch\"atzung \\
\texttt{Lig INP\%}       & \%       & INP-Anteil auf Liganden \\
\texttt{$\pm$sigma}      & \%       & Streuung \\
\texttt{Lig |d|}         & $\angstrom$ & Liganden-Auslenkung
                                          (Norm der Verschiebung) \\
\texttt{$\pm$sigma}      & $\angstrom$ & Streuung \\
\texttt{2nd OOP\%}       & \%       & wie \texttt{Lig OOP\%} f\"ur
                                       2.\ Liganden-Sph\"are
                                       (z.\,B.\ Backbone-N) \\
\texttt{$\pm$sigma}      & \%       & Streuung \\
\texttt{2nd INP\%}       & \%       & 2.\ Liganden-Sph\"are INP \\
\texttt{$\pm$sigma}      & \%       & Streuung \\
\texttt{2nd |d|}         & $\angstrom$ & 2.\ Liganden-Sph\"are
                                          Auslenkung \\
\bottomrule
\end{tabular}
\end{center}

Spalten 21--32 (nicht oben gelistet) enthalten Kern-Anteile
(\texttt{Kern OOP\%}, \texttt{Kern INP\%}, etc.) sowie OOP/INP-
Anteile der einzelnen Cluster-Atome (Fe$_1$, Fe$_2$, S$_1$, S$_2$).

\subsubsection*{\texttt{Gruppen\_OOP, \_INP, \_Winkel, \_Tors}}

Vier separate Sheets, jeweils eine Zeile pro Mode + Spalte pro
Liganden-Gruppe (Cys$_n$, His$_n$, etc.). Werte sind die OOP/INP-
Anteile bzw.\ Bind-Winkel pro Gruppe. Nutzbar f\"ur die
Liganden-spezifische Analyse einer einzelnen Mode.

\subsubsection*{\texttt{Fe\_S\_Cys} und \texttt{Fe\_N\_His}}

Pro Fe-Liganden-Bindung (z.\,B.\ Fe1-Cys72-S, Fe2-Cys83-S, etc.)
eine eigene Spalte. Eine Zeile pro Mode, Inhalt:
\begin{itemize}\setlength{\itemsep}{0pt}
  \item \texttt{stretch}: Anteil der Mode-Bewegung entlang der
        Bindungsachse.
  \item \texttt{bend\_inp}: In-plane-Bend-Anteil (in Cluster-Ebene).
  \item \texttt{bend\_oop}: OOP-Bend-Anteil (senkrecht zur Cluster-
        Ebene).
\end{itemize}

Diese Sheets sind die Datenquelle f\"ur die Mode-Loop-Aggregation
(\S\ref{sec:reorg}); sie zeigen, wie eine einzelne Mode auf jede
Fe-Liganden-Bindung wirkt.

\subsubsection*{\texttt{Kern\_Scores} -- Heuristische Cluster-Zerlegung}

Eine Zeile pro Mode, 10 Spalten mit den heuristischen Cluster-Score-
Funktionen aus \S\ref{sec:scsd}:

\begin{itemize}\setlength{\itemsep}{0pt}
  \item \texttt{Translation}, \texttt{Umbrella}, \texttt{Hinge-Fe},
        \texttt{Hinge-S}: Cluster-Bewegungs-Klassen.
  \item \texttt{OOP-Twist}, \texttt{Fe-Streckung},
        \texttt{S-Streckung}, \texttt{Breathing},
        \texttt{Rhombus-Scherung}, \texttt{Rotation-ip}:
        Sechs in-plane / OOP-Bewegungsmodi.
\end{itemize}

\textbf{Wichtiger Hinweis im Sheet-Header}: Die Kern-Scores sind
\emph{heuristische geometrische Projektionen}, summieren sich
i.\,A.\ nicht zu 1, und sind nicht-orthogonal. F\"ur eine
\textbf{rigorose Symmetriezerlegung} verweist das Sheet auf
\texttt{SCSD} (siehe unten).

\subsubsection*{\texttt{Abstaende\_Gleichgewicht}}

Tabelle der Gleichgewichts-Bindungsl\"angen aller relevanten
Cluster- und Liganden-Bindungen. Eine Zeile pro Bindung. Spalten:
Bindungs-Bezeichnung, Distanz in $\angstrom$, Atom-Indizes.

\subsubsection*{\texttt{Coordination} -- Liganden- und Gruppen-Atomzuordnung
                  \emph{(seit v1.0.3)}}
\label{app:wb-main-coordination}

Diagnose-Sheet, das die PDB-zu-Gaussian-Atomzuordnungen offenlegt, die
alle nachfolgenden Analysen verwenden. Drei Bl\"ocke, durch
Leerzeilen getrennt:
\begin{itemize}\setlength{\itemsep}{0pt}
  \item \textbf{Liganden}: eine Zeile pro Fe-koordinierendem Liganden
        (\texttt{LigandInfo} in \texttt{geometry.py}). Spalten:
        \texttt{label} (z.\,B.\ ``Cys 207''), \texttt{element}
        (S/N/O), \texttt{fe\_center} (1-basierter Gaussian-Center des
        gebundenen Fe), \texttt{lig\_center} (Gaussian-Center des
        Donor-Atoms), \texttt{bond\_len\_A},
        \texttt{his\_protonated} (Boolean), \texttt{h\_center}
        (Gaussian-Center des Imidazol-NH-H, falls vorhanden),
        \texttt{his\_hn\_center} (das Donor-N, das dieses H tr\"agt,
        NE2 oder ND1).
  \item \textbf{Gruppen (\texttt{group\_map})}: eine Zeile pro
        Atomgruppe, die in den \texttt{Groups\_OOP/INP/Winkel/Tors}-
        Bl\"attern verwendet wird. Spalten: \texttt{group}
        (Bezeichnung), \texttt{centers} (kommagetrennte Liste
        1-basierter Gaussian-Centers).
  \item \textbf{Cluster-Atome}: ein 4-Zeilen-Block mit Fe$_1$,
        Fe$_2$, S$_1$, S$_2$, ihren Gaussian-Centers und der
        Gleichgewichts-$(x, y, z)$ in $\angstrom$.
\end{itemize}

Dieses Sheet ist die erste Anlaufstelle, wenn eine Excel-Zeile
unerwartet leer aussieht: eine leere \texttt{centers}-Liste oder ein
fehlendes \texttt{h\_center} erkl\"art eine nachgelagerte
Null-Zeile in den \texttt{Groups\_*}-, \texttt{Fe\_*\_*}- oder
\texttt{His\_HN}-Bl\"attern (\S\ref{sec:migration-v104}).

\subsubsection*{\texttt{SCSD} -- Rigorose D$_{2h}$-Symmetriezerlegung}

Pro Mode 8 Spalten f\"ur die Anteile auf die D$_{2h}$-Irreps:
$A_g, B_{1g}, B_{2g}, B_{3g}, A_u, B_{1u}, B_{2u}, B_{3u}$. Werte
in \% des Cluster-Mode-Beitrags. Im Gegensatz zu \texttt{Kern\_Scores}
sind diese Werte \emph{orthogonal} und summieren sich exakt zu 100\,\%.

Methode: Kingsbury \& Senge (Chem.\ Sci.\ 15, 13638, 2024) mit
kanonischer D$_{2h}$-Referenz (Fe-Fe=2.73\,$\angstrom$,
Fe-S=2.20\,$\angstrom$, Achsen $x$=Fe-Fe, $y$=S-S, $z$=Normale).

\subsubsection*{\texttt{Reorganisationsenergie}}
\label{app:sheet-reorg}

Hauptsheet f\"ur die Reorg-Pipeline. Eine Zeile pro Mode, 21 Spalten:

\begin{center}
\small
\begin{tabular}{p{4.0cm}p{1.4cm}p{7.6cm}}
\toprule
\textbf{Spalte} & \textbf{Einheit} & \textbf{Bedeutung} \\
\midrule
\texttt{freq\_cm1}                  & cm$^{-1}$ & Mode-Frequenz \\
\texttt{dr\_FeFe\_rss\_a}           & $\angstrom$ & RSS-Aggregat
                                                    der Fe-Fe-Modulation \\
\texttt{dr\_FeFe\_sum\_signed\_a}   & $\angstrom$ & Signed-Sum-Aggregat
                                                    (kann negativ sein) \\
\texttt{lambda\_FeFe\_pair\_cm1}    & cm$^{-1}$ & Reorg pro Bindungspaar
                                                  (Sub-Aufschl\"usselung) \\
\texttt{lambda\_FeFe\_mode\_cm1}    & cm$^{-1}$ & Reorg pro Mode
                                                  (Aggregat \"uber alle
                                                  Pairs des Kanals) \\
\midrule
\multicolumn{3}{l}{\textit{Wiederholt f\"ur Kan\"ale: FeN, FeS, NH, HA
       (jeweils 4 Spalten)}} \\
\bottomrule
\end{tabular}
\end{center}

\textbf{Konvention}: \texttt{rss} (Root-Sum-Squared) ist immer
positiv und macht die Marcus-Hush-Additivit\"at exakt
(\S\ref{app:pseudo-reorg}). \texttt{sum\_signed} ist die vorzeichenfreie
Summe \"uber alle Bindungs-Modulationen des Kanals; sie zeigt, ob
die Bindungs-Modulationen koh\"arent (Sum $\approx$ rss) oder
inkoh\"arent (Sum $\ll$ rss) sind.

\textbf{Mode-vs-Pair-Konvention}: \texttt{lambda\_X\_pair} = einzelne
Bindungspaar-Beitr\"age summiert; \texttt{lambda\_X\_mode} = mit
$\mu_i\omega_i^2$-Gewichtung pro Mode aggregiert. Beide sind
erhalten f\"ur unterschiedliche Auswertungs-Zwecke.

\subsubsection*{\texttt{Reorg\_Total} -- System-Aggregate pro Kanal}

Vier Zeilen (eine pro Kanal: FeFe, FeN, FeS, NH, HA -- 5 Zeilen
total). Spalten:

\begin{itemize}\setlength{\itemsep}{0pt}
  \item \texttt{Kanal}: Bezeichnung des Bindungs-Kanals.
  \item \texttt{Lambda\_total\_pair\_cm1}: $\Lambda_X^\text{total}$
        aus Pair-Aggregation.
  \item \texttt{Lambda\_total\_mode\_cm1}: $\Lambda_X^\text{total}$
        aus Mode-Aggregation.
  \item \texttt{n\_modes\_contributing}: Anzahl Modes mit
        nicht-trivialem Beitrag ($\lambda_X(i) > 10^{-6}$\,cm$^{-1}$).
\end{itemize}

Dies sind die zentralen \emph{publikationsfertigen Werte} der
Reorg-Pipeline.

\subsubsection*{\texttt{Reorg\_pro\_Bindung} -- Aufschl\"usselung
                  pro einzelner Bindung}

Eine Zeile pro Einzel-Bindung (z.\,B.\ Fe1-Cys72-S, Fe1-Fe2,
Fe1-S$_b$1, etc.). Sechs Spalten:

\begin{itemize}\setlength{\itemsep}{0pt}
  \item \texttt{Bindung}: Bindungs-Bezeichnung.
  \item \texttt{Parent\_Kanal}: zu welchem Kanal (FeFe, FeS, FeN, ...)
        die Bindung geh\"ort.
  \item \texttt{Akzeptor\_Gewicht}: f\"ur HA-Kanal das Gauss-Gewicht
        des H-Akzeptor-Pairs (Default 1.0 f\"ur Hauptakzeptor).
  \item \texttt{Lambda\_total\_pair\_cm1}, \texttt{Lambda\_total\_mode\_cm1}:
        wie in \texttt{Reorg\_Total}, aber pro Bindung.
  \item \texttt{n\_modes\_contributing}: Modes mit Beitrag $> 10^{-6}$\,cm$^{-1}$.
\end{itemize}

Erlaubt die Frage zu beantworten: \emph{Welche einzelne Fe-S-Bindung
tr\"agt am meisten zu $\Lambda_\text{FeS}$ bei?}

\subsubsection*{\texttt{Modulations\_Spektren}}

Frequenzaufgel\"oste Modulations-Spektren $M_X(\omega)$, plus
Co-Spektren f\"ur PCET-Analyse. Eine Zeile pro Frequenz-Rasterpunkt
(Default: 0.5\,cm$^{-1}$ Schritt). 8 Spalten:

\begin{itemize}\setlength{\itemsep}{0pt}
  \item \texttt{freq\_cm1}: Frequenz-Rasterpunkt.
  \item \texttt{M\_FeFe, M\_FeN, M\_FeS, M\_NH, M\_HA}: Modulations-
        Spektren $M_X(\omega) = \sum_i |\Delta r_X(i)| \cdot
        \mathcal{N}(\omega - \omega_i, \sigma)$ mit Default
        $\sigma=5$\,cm$^{-1}$.
  \item \texttt{C\_ET\_FeS}: Co-Spektrum
        $\sqrt{M_\text{FeFe}(\omega) \cdot M_\text{FeS}(\omega)}$
        -- markiert Frequenzregionen, in denen FeFe und FeS
        gleichzeitig moduliert werden. Hoch-Indikator f\"ur
        ET-relevante Cluster-Atmungs-Modes.
  \item \texttt{C\_PCET}: Co-Spektrum
        $\sqrt{M_\text{NH}(\omega) \cdot M_\text{HA}(\omega)}$
        -- markiert PCET-relevante Modes (gleichzeitige
        N-H- und H-Akzeptor-Modulation).
\end{itemize}

\subsubsection*{\texttt{Lambda\_kumulativ}}

Kumulative Reorg-Kurven $\Lambda_X(\omega) = \sum_{\omega_i \leq \omega}
\lambda_X(i)$. Eine Zeile pro Frequenz-Rasterpunkt. 6 Spalten:
\texttt{freq\_cm1}, \texttt{Lambda\_FeFe}, \texttt{Lambda\_FeN},
\texttt{Lambda\_FeS}, \texttt{Lambda\_NH}, \texttt{Lambda\_HA}.

Jede Spalte ist eine \textbf{monotone Stufenfunktion}: Spr\"unge
markieren Modes mit $\lambda_X$-Beitrag. Direkt mit NRVS-Spektrum
\"uberlagerbar f\"ur Bandenidentifikation.

\subsubsection*{\texttt{B\_Faktoren}}

Pro Atom: thermische B-Faktoren $B_a = \frac{8\pi^2}{3} \langle u_a^2
\rangle$ aus der QHO-Mittelung \"uber alle Modes. Spalten: Atom-Index,
Element, B-Faktor in $\angstrom^2$.

\subsubsection*{\texttt{Info}}

Pro-Lauf-Metadaten: Logfile-Pfad, Anzahl Modes, Frequenz-Range,
Cluster-Geometrie, Liganden-Liste, Konfigurations-Parameter.
Volltext-Beschreibung des Laufs in einem Sheet.

\subsection{Embeddings-Workbook
            (\texttt{*\_analysis\_Embeddings.xlsx})}
\label{app:wb-embeddings}

Enth\"alt die UMAP-Projektionen und HDBSCAN-Clusterungen. Das
Workbook beherbergt \emph{drei} parallele Einbettungen derselben
Mode-Menge, jede in ihrem eigenen Paar von
\texttt{*\_clusters}- / \texttt{*\_profile}-Sheets, plus die
pro-Mode-Roh-C$_\alpha$-Amplituden-Matrix als Eingabe f\"ur das
Ca-UMAP.

\subsubsection*{\texttt{Projektionen}}

Globale UMAP. Eine Zeile pro Mode, Spalten: \texttt{Modus \#},
\texttt{Frequenz}, \texttt{UMAP\_x}, \texttt{UMAP\_y}. Berechnet auf
der vollen Feature-Matrix aus
\texttt{embedding.build\_feature\_matrix} (Marcus-Hush-Reorg-Features
+ Bindungs-Modulations-Features + Kernel-Scores). Direkt plotbar
als 2D-Embedding.

\subsubsection*{\texttt{Ca\_Amplituden}}

Pro Mode und pro C$_\alpha$-Residuum: die thermisch skalierte
$\lVert\evec\rVert$-Amplitude des R\"uckgrat-$\alpha$-Kohlenstoffs
des Residuums. Verwendet als Feature-Matrix f\"ur das Ca-UMAP
(siehe unten). Format: Zeilen = Modes, Spalten = Residuennummern.

\subsubsection*{\texttt{Cluster} / \texttt{Cluster\_Profil}}
\label{app:wb-embed-cluster}

Globale UMAP-Cluster-Zuordnungen. \texttt{Cluster} hat eine Zeile pro
Mode mit ihrer HDBSCAN-Cluster-ID (\texttt{cluster\_label}: $-1$ =
Rauschen; $0, 1, \ldots$ = Cluster-Index), Distanz zum
Cluster-Zentrum, und allen Feature-Werten.
\texttt{Cluster\_Profil} hat eine Zeile pro HDBSCAN-Cluster mit dem
mittleren Feature-Vektor, Fisher-F-Statistik pro Feature, und einer
geordneten Liste repr\"asentativer Modes (n\"achster zum Zentrum).

\subsubsection*{\texttt{SSE\_UMAP\_clusters} / \texttt{SSE\_UMAP\_profile}
                  \emph{(seit v1.0.3)}}
\label{app:wb-embed-ssumap}

Sekund\"arstruktur-UMAP, beschr\"ankt auf Modes, die mindestens einen
nicht-null Eintrag in der pro-SSE-Element-Feature-Zeile haben
(\texttt{sse}-Dict nach \texttt{analyze\_all\_sse} nicht leer).
N\"utzlich zur Identifikation kollektiver
Protein-R\"uckgrat-Modes, die mehrere Sekund\"arstrukturen
gleichzeitig anregen. Das \texttt{*\_clusters}-Sheet listet Modes mit
ihren SSE-UMAP-Koordinaten und HDBSCAN-Cluster-ID; das
\texttt{*\_profile}-Sheet gibt das Z-Score-Profil der
diskriminierenden Top-30-SSE-Elemente pro Cluster aus.

\subsubsection*{\texttt{Ca\_UMAP\_clusters} / \texttt{Ca\_UMAP\_profile}
                  \emph{(seit v1.0.3)}}
\label{app:wb-embed-caumap}

C$_\alpha$-Amplituden-UMAP. Jede Mode wird als Punkt in einem
$N_{C_\alpha}$-dimensionalen Raum dargestellt, eine Dimension pro
R\"uckgrat-$\alpha$-Kohlenstoff (die Spalten von
\texttt{Ca\_Amplituden}). Features werden vor UMAP pro Residuum
z-skaliert, sodass kein einzelnes hochamplitudisches Residuum die
Embedding-Distanz dominiert.

Konzeptionell beantworten die drei UMAPs komplement\"are Fragen:
\begin{itemize}\setlength{\itemsep}{0pt}
  \item \emph{Globales UMAP} (\texttt{Projektionen}): Welche Modes
        koppeln an dieselben Reorg-Kan\"ale (Fe-X, NH, HA)?
        N\"utzlich zur Identifikation ET-aktiver bzw.\
        PCET-aktiver Modes.
  \item \emph{SSE-UMAP} (\texttt{SSE\_UMAP\_*}): Welche Modes haben
        denselben Sekund\"arstruktur-Fu\ss abdruck? N\"utzlich zur
        Identifikation dom\"anenkollektiver Modes.
  \item \emph{Ca-UMAP} (\texttt{Ca\_UMAP\_*}): Welche Modes haben
        \"ahnliche r\"aumliche Verteilungen der R\"uckgrat-Bewegung?
        N\"utzlich zur Identifikation langreichweitiger
        allosterischer Pfade und auf bestimmte Protein-Regionen
        lokalisierter Modes.
\end{itemize}

Begleitende PNGs (nur wenn \texttt{cfg.export\_embedding\_plots = true}):
\texttt{\_embedding\_UMAP.png} (2 Panele),
\texttt{\_embedding\_SSE\_UMAP.png} (3 Panele: Frequenz / Mode-Typ /
Cluster), \texttt{\_embedding\_Ca\_UMAP.png} (3 Panele), und
\texttt{\_ca\_amplitudes\_heatmap.png} (log-skalierte
Residuum $\times$ Frequenz Heatmap der C$_\alpha$-Amplituden).

\subsection{Interpoliertes-pDOS-Workbook
            (\texttt{*\_analysis\_interp0.05.xlsx})}
\label{app:wb-interp}

Ein einziges Sheet \texttt{Interpoliert} mit der pDOS-Auswertung
auf einem feinen Frequenzraster (Default 0.05\,cm$^{-1}$). Spalten:
\texttt{freq\_cm1}, \texttt{pDOS\_Fe1}, \texttt{pDOS\_Fe2},
\texttt{pDOS\_S1}, \texttt{pDOS\_S2}, \texttt{pDOS\_total}.
Verwendet f\"ur den direkten NRVS-Spektrums-Vergleich.

\subsection{Sekund\"arstruktur-Workbook
            (\texttt{*\_analysis\_SSE.xlsx}, optional)}
\label{app:wb-sse}

Wird nur erzeugt, wenn ein PDB-File mit (oder DSSP-extrahierten)
Sekund\"arstruktur-Records mitgeliefert wird. Enth\"alt typischerweise
9--11 Sheets, je nach SSE-Komplexit\"at:

\subsubsection*{Pro-SSE-Element-Sheets (\texttt{SSE\_*})}

Pro SSE-Metrik ein eigenes Sheet:
\texttt{SSE\_amplitude\_mean},
\texttt{SSE\_amplitude\_max},
\texttt{SSE\_com\_amplitude} (Schwerpunkts-Amplitude),
\texttt{SSE\_lateral\_std} (laterale Bewegungs-Streuung),
\texttt{SSE\_lateral\_amplitude},
\texttt{SSE\_stretching} (axiale L\"angen-\"Anderung),
\texttt{SSE\_axial\_amplitude},
\texttt{SSE\_tilting\_angle} (Helix-Tilt),
\texttt{SSE\_internal\_amplitude} (interne Bewegung).

Jedes Sheet hat eine Zeile pro Mode + Spalte pro SSE-Element.

\subsubsection*{\texttt{SSE\_UMAP\_Cluster} und \texttt{SSE\_UMAP\_Profil}}

Wie \texttt{UMAP\_Cluster} und \texttt{UMAP\_Cluster\_Profil}, aber
auf der SSE-Feature-Matrix (statt der Cluster-Liganden-Matrix). Findet
strukturelle Cluster der Modes basierend auf ihrer Wirkung auf
Sekund\"arstruktur-Elemente.

\begin{infobox}
\textbf{Welche Sheets sind Pflicht f\"ur eine Veroeffentlichung?}
F\"ur eine Reorg-Manuskript-Tabelle reichen typisch \texttt{Reorg\_Total}
und \texttt{Reorg\_pro\_Bindung} (Hauptauswertung). F\"ur die
Bandenzuordnung kommt \texttt{Lambda\_kumulativ} dazu (NRVS-
\"Uberlagerung). F\"ur die Mode-Charakterisierung das Sheet
\texttt{Modenanalyse} und (bei symmetrie-relevanten Fragen)
\texttt{SCSD}. Der Rest dient internen Diagnose-Zwecken.
\end{infobox}


% =========================================================================
% Literatur
% =========================================================================

\begin{thebibliography}{99}

\bibitem{Kingsbury2024}
C.~J. Kingsbury, M.~O. Senge,
\emph{Quantifying near-symmetric molecular distortion using
symmetry-coordinate structural decomposition},
\textbf{Chem.\ Sci.} \textbf{15}, 13638--13649 (2024).
DOI: 10.1039/D4SC01670J. Web-Implementation:
\url{https://www.kingsbury.id.au/scsd}.

\bibitem{KingsburyPreprint2024}
C.~J. Kingsbury, M.~O. Senge,
\emph{Quantifying near-symmetric molecular distortion using
symmetry-coordinate structural decomposition},
ChemRxiv preprint, doi:10.26434/chemrxiv-2024-6b25s (2024).

\bibitem{Iwata1996}
S.~Iwata, M.~Saynovits, T.~A. Link, H.~Michel,
\emph{Structure of a water soluble fragment of the {`Rieske'} iron-sulfur
protein of the bovine heart mitochondrial cytochrome bc$_1$ complex
determined by {MAD} phasing at 1.5 \AA{} resolution},
\textbf{Structure} \textbf{4}, 567--579 (1996).

\bibitem{SauerHeyden2023}
P.~Sauer, M.~Heyden,
\emph{Frequency-selective anharmonic mode analysis of thermally excited
vibrations in proteins},
\textbf{J.\ Chem.\ Phys.} (2023).

\bibitem{Achterhold2002}
K.~Achterhold, F.~G. Parak,
\emph{Protein dynamics: determination of anisotropic vibrations at the
heme iron of myoglobin},
\textbf{J.\ Phys.: Condens.\ Matter} \textbf{14}, S1399--S1421 (2002).

\bibitem{Paulsen1999}
H.~Paulsen, H.~Winkler, A.~X. Trautwein \emph{et al.},
\emph{Measurement and simulation of nuclear inelastic-scattering spectra
of molecular crystals},
\textbf{Phys.\ Rev.\ B} \textbf{59}, 975 (1999).

\bibitem{Gee2021}
L.~B. Gee, V.~Pelmenschikov, C.~Mons, N.~Mishra, H.~Wang, Y.~Yoda,
K.~Tamasaku, M.-P. Golinelli-Cohen, S.~P. Cramer,
\emph{NRVS and DFT of mitoNEET: Understanding the special vibrational
structure of a [2Fe--2S] cluster with (Cys)$_3$(His)$_1$ ligation},
\textbf{Biochemistry} \textbf{60}, 2419--2424 (2021).
DOI: 10.1021/acs.biochem.1c00252.

\bibitem{Conlan2011}
A.~R. Conlan, M.~L. Paddock, C.~Homer \emph{et al.},
\emph{Mutation of the His ligand in mitoNEET stabilizes the 2Fe--2S
cluster despite conformational heterogeneity in the ligand environment},
\textbf{Acta Cryst.\ F} \textbf{67}, 516--523 (2011).

\bibitem{GoNoguti1983}
N.~Go, T.~Noguti, T.~Nishikawa,
\emph{Dynamics of a small globular protein in terms of low-frequency
vibrational modes},
\textbf{Proc.\ Natl.\ Acad.\ Sci.\ USA} \textbf{80}, 3696 (1983).

\bibitem{ParakKnapp1984}
F.~Parak, E.~W. Knapp,
\emph{A consistent picture of protein dynamics},
\textbf{Proc.\ Natl.\ Acad.\ Sci.\ USA} \textbf{81}, 7088 (1984).

\bibitem{WilsonDeciusCross1955}
E.~B. Wilson, J.~C. Decius, P.~C. Cross,
\emph{Molecular Vibrations: The Theory of Infrared and Raman
Vibrational Spectra},
McGraw-Hill, New York (1955); Dover Reprint (1980).

\bibitem{Sturhahn2004}
W.~Sturhahn,
\emph{Nuclear resonant spectroscopy},
\textbf{J.\ Phys.: Condens.\ Matter} \textbf{16}, S497--S530 (2004).

\bibitem{Sturhahn2000}
W.~Sturhahn, A.~Chumakov,
\emph{Lamb-M\"ossbauer factor and second-order Doppler shift from
inelastic nuclear resonant absorption},
\textbf{Hyperfine Interact.} \textbf{123/124}, 809--824 (2000).

\bibitem{Lipkin1995}
H.~J. Lipkin,
\emph{M\"ossbauer sum rules for use with synchrotron sources},
\textbf{Phys.\ Rev.\ B} \textbf{52}, 10073--10079 (1995).

\bibitem{SingwiSjolander1960}
K.~S. Singwi, A.~Sj\"olander,
\emph{Resonance absorption of nuclear gamma rays and the dynamics of
atomic motions},
\textbf{Phys.\ Rev.} \textbf{120}, 1093--1102 (1960).

\bibitem{KohnRuocco1998}
V.~G. Kohn, A.~I. Chumakov, R.~R\"uffer,
\emph{Nuclear resonant inelastic absorption of synchrotron radiation
in an anisotropic single crystal},
\textbf{Phys.\ Rev.\ B} \textbf{58}, 8437--8444 (1998).

\bibitem{Kabsch1976}
W.~Kabsch,
\emph{A solution for the best rotation to relate two sets of vectors},
\textbf{Acta Cryst.\ A} \textbf{32}, 922--923 (1976).

\bibitem{HammesSchifferStuchebrukhov2010}
S.~Hammes-Schiffer, A.~A. Stuchebrukhov,
\emph{Theory of coupled electron and proton transfer reactions},
\textbf{Chem.\ Rev.} \textbf{110}, 6939--6960 (2010).

\bibitem{Marcus1956}
R.~A. Marcus,
\emph{On the Theory of Oxidation-Reduction Reactions Involving
Electron Transfer. I},
\textbf{J.\ Chem.\ Phys.} \textbf{24}, 966--978 (1956).
DOI: 10.1063/1.1742723.

\bibitem{Marcus1993}
R.~A. Marcus,
\emph{Electron transfer reactions in chemistry. Theory and experiment}
(Nobel Lecture),
\textbf{Angew.\ Chem.\ Int.\ Ed.\ Engl.} \textbf{32}, 1111--1121 (1993).
DOI: 10.1002/anie.199311113.

\bibitem{Marcus1985}
R.~A. Marcus, N.~Sutin,
\emph{Electron transfers in chemistry and biology},
\textbf{Biochim.\ Biophys.\ Acta} \textbf{811}, 265--322 (1985).

\bibitem{Sturhahn1999}
W.~Sturhahn, V.~G.~Kohn,
\emph{Theoretical aspects of incoherent nuclear resonant scattering},
\textbf{Hyperfine Interact.} \textbf{123/124}, 367--377 (1999).

\bibitem{Chumakov1999}
A.~I.~Chumakov, W.~Sturhahn,
\emph{Experimental aspects of inelastic nuclear resonant scattering},
\textbf{Hyperfine Interact.} \textbf{123/124}, 781--808 (1999).

\bibitem{MaradudinFlinn1963}
A.~A.~Maradudin, P.~A.~Flinn,
\emph{Anharmonic contributions to the Debye-Waller factor},
\textbf{Phys.\ Rev.} \textbf{129}, 2529 (1963).

\bibitem{Lipkin1960}
H.~J.~Lipkin,
\emph{Some simple features of the M\"o\ss bauer effect},
\textbf{Ann.\ Phys.} (N.Y.) \textbf{9}, 332 (1960).

\bibitem{Hu2003}
M.~Y.~Hu, W.~Sturhahn, T.~S.~Toellner \emph{et al.},
\emph{Measuring velocity of sound with nuclear resonant inelastic
X-ray scattering},
\textbf{Phys.\ Rev.\ B} \textbf{67}, 094304 (2003).

\bibitem{Leu2008}
B.~M.~Leu, M.~Z.~Hu, J.~Zhao, S.~P.~Cramer,
\emph{Quantitative vibrational dynamics of iron in carbonyl compounds},
\textbf{J.\ Am.\ Soc.\ Mass Spectrom.} \textbf{17}, 235 (2008).

\bibitem{ThompsonCoxHastings1987}
P.~Thompson, D.~E.~Cox, J.~B.~Hastings,
\emph{Rietveld refinement of Debye-Scherrer synchrotron X-ray data
from $\mathrm{Al}_2\mathrm{O}_3$},
\textbf{J.\ Appl.\ Crystallogr.} \textbf{20}, 79 (1987).

\bibitem{ina32}
K.~Achterhold, F.~G.~Parak,
\emph{ina32.f -- Fortran source code for analysis of nuclear inelastic
scattering spectra}; interne Publikation, TU M\"unchen, Abteilung
Biophysik (2005--).

\bibitem{Codata2018}
E.~Tiesinga, P.~J.~Mohr, D.~B.~Newell, B.~N.~Taylor,
\emph{CODATA recommended values of the fundamental physical constants:
2018},
\textbf{Rev.\ Mod.\ Phys.} \textbf{93}, 025010 (2021).

% --- Embedding-Theorie -------------------------------------------------

\bibitem{Pearson1901}
K.~Pearson,
\emph{On lines and planes of closest fit to systems of points in space},
\textbf{Philos.\ Mag.} \textbf{2}, 559--572 (1901).

\bibitem{Hotelling1933}
H.~Hotelling,
\emph{Analysis of a complex of statistical variables into principal
components},
\textbf{J.\ Educ.\ Psychol.} \textbf{24}, 417--441 (1933).

\bibitem{Jolliffe2002}
I.~T. Jolliffe,
\emph{Principal Component Analysis},
2nd ed., Springer (2002).

\bibitem{vanDerMaaten2008}
L.~van der Maaten, G.~Hinton,
\emph{Visualizing data using t-SNE},
\textbf{J.\ Mach.\ Learn.\ Res.} \textbf{9}, 2579--2605 (2008).

\bibitem{vanDerMaaten2014}
L.~van der Maaten,
\emph{Accelerating t-SNE using tree-based algorithms},
\textbf{J.\ Mach.\ Learn.\ Res.} \textbf{15}, 3221--3245 (2014).

\bibitem{McInnes2018}
L.~McInnes, J.~Healy, J.~Melville,
\emph{UMAP: Uniform Manifold Approximation and Projection for
Dimension Reduction},
\textbf{arXiv} 1802.03426 (2018).

\bibitem{Becht2019}
E.~Becht, L.~McInnes, J.~Healy \emph{et al.},
\emph{Dimensionality reduction for visualizing single-cell data using
UMAP},
\textbf{Nat.\ Biotechnol.} \textbf{37}, 38--44 (2019).

\bibitem{Campello2013}
R.~J.~G.~B. Campello, D.~Moulavi, J.~Sander,
\emph{Density-based clustering based on hierarchical density estimates},
in \textbf{Adv.\ Knowl.\ Discov.\ Data Min.} (PAKDD 2013), pp.~160--172.

\bibitem{McInnes2017}
L.~McInnes, J.~Healy, S.~Astels,
\emph{hdbscan: Hierarchical density based clustering},
\textbf{J.\ Open Source Softw.} \textbf{2}, 205 (2017).

% --- Erweiterte PCET/ET-Theorie ----------------------------------------

\bibitem{HammesSchiffer2010}
S.~Hammes-Schiffer,
\emph{Theory of proton-coupled electron transfer in energy conversion
processes},
\textbf{Acc.\ Chem.\ Res.} \textbf{42}, 1881--1889 (2009).
DOI: 10.1021/ar9001284.

\bibitem{EdwardsSoudackov2010}
S.~J. Edwards, A.~V. Soudackov, S.~Hammes-Schiffer,
\emph{Driving force dependence of rates for nonadiabatic proton and
proton-coupled electron transfer: Conditions for inverted region behavior},
\textbf{J.\ Phys.\ Chem.\ B} \textbf{113}, 14545--14548 (2009).

\bibitem{LayfieldHammesSchiffer2014}
J.~P. Layfield, S.~Hammes-Schiffer,
\emph{Hydrogen tunneling in enzymes and biomimetic models},
\textbf{Chem.\ Rev.} \textbf{114}, 3466--3494 (2014).

\bibitem{CuiKarplus2008}
Q.~Cui, M.~Karplus,
\emph{Allostery and cooperativity revisited},
\textbf{Protein Sci.} \textbf{17}, 1295--1307 (2008).

\bibitem{Bergner2017}
M.~Bergner, S.~Dechert, S.~Demeshko, C.~Kupper, J.~M.~Mayer, F.~Meyer,
\emph{Model of the {MitoNEET} {[}2Fe-2S{]} cluster shows proton coupled
electron transfer},
\textbf{J.\ Am.\ Chem.\ Soc.} \textbf{139}, 701--707 (2017).
DOI: 10.1021/jacs.6b09180.

\bibitem{Kuila1992}
D.~Kuila, J.~R. Schoonover, R.~B. Dyer, C.~J. Batie, D.~P. Ballou,
J.~A. Fee, W.~H. Woodruff,
\emph{Resonance {R}aman studies of {R}ieske-type proteins},
\textbf{Biochim.\ Biophys.\ Acta} \textbf{1140}, 175--183 (1992).

\bibitem{Wang2025}
H.~Wang, V.~Pelmenschikov, Y.~Yoda, S.~P. Cramer,
\emph{NRVS of {Fe-S} cluster proteins and models -- A bestiary of
nifty normal modes},
\textbf{J.\ Inorg.\ Biochem.} \textbf{270}, 112935 (2025).
DOI: 10.1016/j.jinorgbio.2025.112935.

\bibitem{Zuris2011}
J.~A. Zuris, D.~A. Halim, A.~R. Conlan \emph{et al.},
\emph{Engineering the redox potential over a wide range within a
new class of FeS proteins},
\textbf{J.\ Am.\ Chem.\ Soc.} \textbf{132}, 13120--13122 (2010).

\bibitem{KurodaHammesSchiffer2020}
T.~Kuroda, S.~Hammes-Schiffer,
\emph{Vibronic coupling for proton-coupled electron transfer reactions:
Recent insights from theory},
\textbf{Curr.\ Opin.\ Electrochem.} \textbf{29}, 100789 (2021).

% --- PCET/ET-Literatur, Eisen-Schwefel-Spektroskopie ------------------

\bibitem{HammesSchifferSoudackov2008}
S.~Hammes-Schiffer, A.~V.~Soudackov,
\emph{Proton-coupled electron transfer in solution, proteins, and
electrochemistry},
\textbf{J.\ Phys.\ Chem.\ B} \textbf{112}, 14108--14123 (2008).

\bibitem{HuangRhys1950}
K.~Huang, A.~Rhys,
\emph{Theory of light absorption and non-radiative transitions in
$F$-centres},
\textbf{Proc.\ R.\ Soc.\ Lond.\ A} \textbf{204}, 406--423 (1950).

\bibitem{Beinert1997}
H.~Beinert, R.~H.~Holm, E.~M\"unck,
\emph{Iron-sulfur clusters: nature's modular, multipurpose structures},
\textbf{Science} \textbf{277}, 653--659 (1997).

\bibitem{Han1989}
S.~Han, R.~S.~Czernuszewicz, T.~G.~Spiro,
\emph{Vibrational spectra and normal mode analysis for {[}2Fe-2S{]}
protein analogs using sulfur-34, iron-54 and deuterium substitution:
coupling of iron-sulfur stretching and sulfur-carbon-carbon bending modes},
\textbf{J.\ Am.\ Chem.\ Soc.} \textbf{111}, 3496--3504 (1989).

\bibitem{Czernuszewicz1986}
R.~S.~Czernuszewicz, J.~LeGall, I.~Moura, T.~G.~Spiro,
\emph{Resonance Raman spectra of rubredoxin: new assignments and
vibrational coupling mechanism from iron-54/iron-56 isotope shifts
and variable-wavelength excitation},
\textbf{Inorg.\ Chem.} \textbf{25}, 696--700 (1986).

\bibitem{Bergmann2003}
U.~Bergmann, W.~Sturhahn, D.~E.~Linn, F.~E.~Jenney, M.~W.~W.~Adams,
K.~Rupnik, B.~J.~Hales, E.~E.~Alp, A.~Mayse, S.~P.~Cramer,
\emph{Observation of {Fe-H/D} modes by nuclear resonant vibrational
spectroscopy},
\textbf{J.\ Am.\ Chem.\ Soc.} \textbf{125}, 4016--4017 (2003).

\bibitem{Saouma2014}
C.~T.~Saouma, M.~M.~Pinney, J.~M.~Mayer,
\emph{Electron transfer and proton-coupled electron transfer reactivity
and self-exchange of synthetic {[}2Fe-2S{]} complexes:
Models for Rieske and mitoNEET clusters},
\textbf{Inorg.\ Chem.} \textbf{53}, 3153--3161 (2014).

\bibitem{Albers2014Rieske}
A.~Albers, S.~Demeshko, S.~Dechert, C.~T.~Saouma, J.~M.~Mayer, F.~Meyer,
\emph{Fast proton-coupled electron transfer observed for a high-fidelity
structural and functional {[}2Fe-2S{]} {Rieske} model},
\textbf{J.\ Am.\ Chem.\ Soc.} \textbf{136}, 3946--3954 (2014).

% --- Methodische Quellen fuer Vibrationsanalyse -----------------------

\bibitem{Cyvin1968}
S.~J.~Cyvin,
\emph{Molecular Vibrations and Mean Square Amplitudes},
Universitetsforlaget, Oslo / Elsevier, Amsterdam (1968).

\bibitem{Ochterski1999}
J.~W.~Ochterski,
\emph{Vibrational Analysis in Gaussian},
Gaussian, Inc., White Paper (1999).
\url{https://gaussian.com/vib/}

\bibitem{Frisch2016Gaussian16}
M.~J.~Frisch \emph{et al.},
\emph{Gaussian 16, Revision C.01},
Gaussian, Inc., Wallingford CT (2016).

\bibitem{NeeseORCAmanual2023}
F.~Neese,
\emph{The ORCA program system: user's manual},
Version 6.0, Max-Planck-Institut f\"ur Kohlenforschung,
M\"ulheim a.\,d.\,Ruhr (2023).
\url{https://www.faccts.de/docs/orca/6.0/manual/}

\bibitem{Neese2022WIRES}
F.~Neese,
\emph{Software update: The ORCA program system---Version 5.0},
\textbf{WIREs Comput.\ Mol.\ Sci.} \textbf{12}, e1606 (2022).

\end{thebibliography}

\end{document}
